% concatenated from main.tex
%%
%% This is LaTeX2e input.
%%

%% The following tells LaTeX that we are using the 
%% style file amsart.cls (That is the AMS article style
%%
\documentclass{amsart}
\usepackage{amssymb}
\usepackage{scalerel}
\usepackage{comment}
%% This has a default type size 10pt.  Other options are 11pt and 12pt
%% This are set by replacing the command above by
%% \documentclass[11pt]{amsart}
%%
%% or
%%
%% \documentclass[12pt]{amsart}
%%

%%
%% Some mathematical symbols are not included in the basic LaTeX
%% package.  Uncommenting the following makes more commands
%% available. 
%%

%\usepackage{amssymb}

%%
%% The following is commands are used for importing various types of
%% grapics.
%% 

%\usepackage{epsfig}  		% For postscript
%\usepackage{epic,eepic}       % For epic and eepic output from xfig

%%
%% The following is very useful in keeping track of labels while
%% writing.  The variant   \usepackage[notcite]{showkeys}
%% does not show the labels on the \cite commands.
%% 

%\usepackageshowkeys}


%%%%
%%%% The next few commands set up the theorem type environments.
%%%% Here they are set up to be numbered section.number, but this can
%%%% be changed.
%%%%

\newtheorem{thm}{Theorem}[section]
\newtheorem{prop}[thm]{Proposition}
\newtheorem{lem}[thm]{Lemma}
\newtheorem{cor}[thm]{Corollary}


%%
%% If some other type is need, say conjectures, then it is constructed
%% by editing and uncommenting the following.
%%

%\newtheorem{conj}[thm]{Conjecture} 


%%% 
%%% The following gives definition type environments (which only differ
%%% from theorem type invironmants in the choices of fonts).  The
%%% numbering is still tied to the theorem counter.
%%% 

\theoremstyle{definition}
\newtheorem{definition}[thm]{Definition}
\newtheorem{example}[thm]{Example}

%%
%% Again more of these can be added by uncommenting and editing the
%% following. 
%%

%\newtheorem{note}[thm]{Note}


%%% 
%%% The following gives remark type environments (which only differ
%%% from theorem type invironmants in the choices of fonts).  The
%%% numbering is still tied to the theorem counter.
%%% 


\theoremstyle{remark}

\newtheorem{remark}[thm]{Remark}


%%%
%%% The following, if uncommented, numbers equations within sections.
%%% 



\numberwithin{equation}{section}


%%%
%%% The following show how to make definition (also called macros or
%%% abbreviations).  For example to use get a bold face R for use to
%%% name the real numbers the command is \mathbf{R}.  To save typing we
%%% can abbreviate as

\newcommand{\ip}[2]{\left\langle {#1} , {#2} \right\rangle} %inner product
\newcommand{\norm}[1]{\left\lVert #1 \right\rVert} %standard norm
\newcommand{\Lone}[4]{\int_{#1}^{#2} {#3} d{#4}} %Integral
\newcommand{\seq}[1]{\{{#1}_{n}\}_{n=1}^{\infty}} %sequence notation
\newcommand{\s}[3]{\sum_{n=#1}^{#2}#3} %sum notation
\newcommand{\R}{\mathbb{R}} %reals
\newcommand{\Z}{\mathbb{Z}} %integers
%\newcommand{\N}{\mathbb{N}} %naturals
\newcommand{\C}{\mathbb{C}} %complexs
\newcommand{\D}{\mathbb{D}} %disc
\newcommand{\T}{\mathbb{T}} %torus
\newcommand{\W}[2]{W_{1}\left(#1 , #2\right)}
\newcommand{\Wp}[3]{W_{#1}\left(#2, #3\right)}
\newcommand{\Hi}{\mathcal{H}} %hilbert space
\newcommand{\tr}{\text{tr}} %trace \newcommand{\calF}{\mathcal{F}}
\newcommand{\Falpha}{\mathcal{F}^{2}_{\alpha}}
\newcommand{\calF}{\mathcal{F}}
\newcommand{\Weyl}{W}
\newcommand{\sep}{\mathrm{sep}}


%%
%% The comment after the defintion is not required, but if you are
%% working with someone they will likely thank you for explaining your
%% definition.  
%%
%% Now add you own definitions:
%%

%%%
%%% Mathematical operators (things like sin and cos which are used as
%%% functions and have slightly different spacing when typeset than
%%% variables are defined as follows:
%%%

\DeclareMathOperator{\dist}{dist} % The distance.



%%
%% This is the end of the preamble.
%% 

\begin{document}

%%
%% The title of the paper goes here.  Edit to your title.
%%

\title{A Semi-Classical Szeg\H{o}-type Limit Theorem for Toeplitz Operators}

%%
%% Now edit the following to give your name and address:
%% 

\author{Trevor Camper}
\address{Department of Mathematics, Dartmouth College,
Hanover, NH 03755}
\email{trevor.camper@dartmouth.edu}

%%
%% If there is another author uncomment and edit the following.
%%

\author{Mishko Mitkovski}\thanks{Supported in part by NSF Grant DMS-2453810.}
\address{School of Mathematical and Statistical Sciences, Clemson University,
Clemson, SC 29630}
\email{mmitkov@clemson.edu}

%%
%% If there are three of more authors they are added in the obvious
%% way. 
%%

%%%
%%% The following is for the abstract.  The abstract is optional and
%%% if not used just delete, or comment out, the following.
%%%


\begin{abstract}
We establish an abstract Szeg\H{o}-type limit theorem for Toeplitz operators associated with normalized continuous Parseval frames. Our assumptions consist essentially of a natural normalization of the frame measures and an approximate-identity condition for the normalized frame vectors. Under these hypotheses, the Szeg\H{o} asymptotics follow from a general argument based on Berezin--Lieb inequalities and
convergence of Berezin transforms, independently of the particular geometry of the underlying space.

As a principal application, we obtain a Szeg\H{o}-type limit theorem for Toeplitz operators on weighted Bergman spaces of the unit ball. This setting lies outside the standard Widom-type localization approach because the underlying hyperbolic measure has exponential volume growth. We also obtain Szeg\H{o} limit theorems for metrizable compact and locally compact Abelian groups, extending the classical compact-group results to more general symbols and recovering recent results for multidimensional tori as special cases. Further
applications include Gabor localization operators, Paley--Wiener spaces, eigenvalue distribution, and sampling density results.

\end{abstract}

%%
%%  LaTeX will not make the title for the paper unless told to do so.
%%  This is done by uncommenting the following.
%%

\maketitle

%%
%% LaTeX can automatically make a table of contents.  This is done by
%% uncommenting the following:
%%

%\tableofcontents

%%
%%  To enter text is easy.  Just type it.  A blank line starts a new
%%  paragraph. 
%%



%%
%%  To put mathematics in a line it is put between dollor signs.  That
%%  is $(x+y)^2=x^2+2xy+y^2$
%%

%%
%%% Displayed mathematics is put between double dollar signs.  
%%



%%
%% Its proof is set off by
%% 
%%
%% A new section is started as follows:
%%
%%%%%%%%%%%%%%%%%%%%%%%%%%%%%%%%%%%%%%%%%%%%%%%%%%%%%%%%%%%%%%%%%%%%%%
\section{Introduction}
%%%%%%%%%%%%%%%%%%%%%%%%%%%%%%%%%%%%%%%%%%%%%%%%%%%%%%%%%%%%%%%%%%%%%%
%\cite{SZ15,Wid06,laptev1996szego}

The classical Szeg\H o limit theorem \cite{SZ15}, sometimes also called the first Szeg\H o theorem, describes the asymptotic behavior of the spectra of Toeplitz matrices generated by a fixed symbol as the matrix size tends to infinity. More precisely, let $\sigma:\mathbb{T} \to \mathbb{R}$ be a symbol, and consider the compressions $P_n M_\sigma P_n : L^2(\mathbb{T}) \to L^2(\mathbb{T})$ of the multiplication operator $M_\sigma : L^2(\mathbb{T}) \to L^2(\mathbb{T})$, $M_\sigma f = \sigma f$, onto the span of the complex exponentials $\operatorname{span}\{ e^{ikx} : k = 0, \pm 1, \pm 2, \dots, \pm n \} \subseteq L^2(\mathbb{T})$. The matrices of these compressions, represented in this exponential basis, are the Toeplitz matrices $[\hat{\sigma}(i-j)]_{i,j=-n,\dots,n}$.

 One of the most classical forms of the Szeg\H o limit theorem states that for a strictly positive symbol $\sigma\in L^{1}(\mathbb{T})$ we have 

\begin{equation}\label{Szegolog}
\lim_{n\to\infty}\frac{1}{2n+1}\text{tr}\left(\log\left(P_{n}M_{\sigma}P_{n}\right)\right)=\frac{1}{2\pi}\int_{0}^{2\pi}\log\sigma\left(e^{i\theta}\right)d\theta. \end{equation}

\noindent Another, equally classical, version of the Szeg\H{o} limit theorem requires the symbol to be continuous, but not necessarily positive. In this case, for any function $\psi$ continuous on a closed interval containing the range of $\sigma$, we have 
\begin{equation}\label{Szegocont}
    \lim_{n\to\infty}\frac{1}{2n+1}\text{tr}\left(\psi\left(P_{n}M_{\sigma}P_{n}\right)\right)=\frac{1}{2\pi}\int_{0}^{2\pi}\psi\left(\sigma\left(e^{i\theta}\right)\right)d\theta. 
\end{equation}
Here, as everywhere above, $P_n$ denotes the orthogonal projection onto 
$\operatorname{span} \{ e^{ikx} : k = 0, \pm 1, \pm 2, \dots, \pm n \}$.
 
 
Numerous authors have explored extensions and generalizations of the first Szeg\H{o} limit theorem in various contexts and settings. A few notable works in this area are~\cite{guille79,kryein1983some,laptev1996szego,widom1979eigenvalue,widom2006asymptotic}. For a more comprehensive overview of the literature and the historical development of the problem, see  ~\cite{bottcher2013analysis,guille79,simon2005orthogonal}. Despite its relatively advanced age, the subject still attracts interest among mathematicians~\cite{guo2024first,MO24,nikolski2020szeg}. 
 
%In particular, Feichtinger and Nowak \cite{FN01} obtained results of a similar form in the context of time-frequency analysis. Let $\sigma\in L^{1}(\mathbb{R}^{2n})\cap L^{\infty}(\R^{2n})$ with $\sigma$ non-negative. Consider the Gabor-Toeplitz operator $T_{\sigma}$ on $L^{2}(\mathbb{R}^{n})$ by 
%\begin{align*} T_{\sigma}f=\int_{\mathbb{R}^{2n}}\sigma(p,q)V_{\phi}f(p,q)\phi_{(p,q)}dpdq,
%\end{align*}
%where $\{\phi_{(p,q)}(x)\}_{(p,q)\in\mathbb{R}^{2n}}=\{e^{2\pi ipx}\phi(x-q)\}_{(p,q)\in\mathbb{R}^{2n}}$ is the associated Gabor system with window function $\phi$ and $V_{\phi}f(p,q)=\int f(x)e^{2\pi ipx}\phi(x-q)dx=\ip{f}{\phi_{(p,q)}}$ is the short-time Fourier transform \cite{grochenig2013foundations}. They consider the asymptotic behavior of these operators, instead under dilations of the symbol to obtain 
%\begin{align*}    \lim_{R\to\infty}\frac{1}{R^{2n}}\text{tr}\left(T_{\sigma_{R}}\psi\left(T_{\sigma_{R}}\right)\right)=\int_{\mathbb{R}^{2n}}\sigma(\eta)\psi\left(\sigma(\eta)\right)d\eta, \tag{1.2}
%\end{align*}
%where $\sigma_{R}(p,q)=\sigma(p/R,q/R)$.\\
\begin{comment} Our first goal in this paper is to extend the classical Szeg\H{o} limit theorem to a more general setting. 
We introduce an abstract framework and establish several abstract forms of the first Szeg\H{o} limit theorem, 
which recover many classical Szeg\H{o} theorems as special cases. 
In particular, we obtain Szeg\H{o} limit theorems for higher-dimensional (including infinite-dimensional) tori, 
recently proved in~\cite{guo2024first} and motivated by conjectures in~\cite{nikolski2020szeg}. 
Notably, our approach does not require a group or complex structure on the domain. 
As an application, we derive the following complement to the well-known result of B\'edos~\cite{bedos1996folner}, 
recently established in~\cite{guo2024first} for the infinite-dimensional torus:
\end{comment}

The purpose of the present paper is both conceptual and concrete. Conceptually, we identify a common mechanism behind a broad class of first Szeg\H{o} limit theorems. We formulate an abstract theory for Toeplitz operators associated with normalized continuous Parseval frames. The hypotheses consist essentially of two properties: a normalization condition for the underlying measures and an approximate-identity condition expressing localization of the normalized frame vectors. Once these properties are verified, the proof of the Szeg\H{o} asymptotic is independent of the particular geometry of the underlying space and follows from the same argument based on Berezin--Lieb inequalities and convergence of Berezin transforms. Thus our approach separates the geometric input, which is needed to verify localization of the kernels, from the operator-theoretic argument producing the asymptotic trace formula.

\subsection{The abstract theorem}

The common mechanism behind all of our results is a Szeg\H{o} theorem for
Toeplitz operators associated with normalized continuous Parseval frames.
Let $X$ be a locally compact Hausdorff space and, for each parameter
$\alpha$, let $\mathcal{H}_\alpha$ be a Hilbert space equipped with a normalized
continuous Parseval frame
\[
  \{k_x^\alpha:x\in X\}\subset \mathcal{H}_\alpha
\]
and frame Radon measure $\nu_\alpha$.  Thus $\|k_x^\alpha\|_\alpha=1$ and
\[
  f=\int_X \langle f,k_x^\alpha\rangle_\alpha k_x^\alpha\,d\nu_\alpha(x),
  \qquad f\in \mathcal{H}_\alpha,
\]
where the integral is understood weakly.  We make two assumptions.
First, the frame measures have a common scale:
\begin{equation}
  d\nu_\alpha(x)=c(\alpha)\,d\nu(x),
  \label{intro:measure-scaling}
\end{equation}
where $\nu$ is a Radon measure independent of $\alpha$.  Second, the normalized frame
vectors become asymptotically localized: for every $x\in X$ and every
neighborhood $U$ of $x$,
\begin{equation}
  \int_{X\setminus U}
  \bigl|\langle k_x^\alpha,k_y^\alpha\rangle_\alpha\bigr|^2
  \,d\nu_\alpha(y)\longrightarrow 0
  \qquad (\alpha\to\infty).
  \label{intro:localization}
\end{equation}

For a real-valued symbol $\sigma\in L^1(X,\nu)$, set
\[
  T_\sigma^\alpha f
  =\int_X \sigma(x)\langle f,k_x^\alpha\rangle_\alpha
  k_x^\alpha\,d\nu_\alpha(x).
\]
The associated Berezin transform is
\[
  \widetilde\sigma_\alpha(x)
  =\int_X \sigma(y)
  \bigl|\langle k_x^\alpha,k_y^\alpha\rangle_\alpha\bigr|^2
  \,d\nu_\alpha(y).
\]

\medskip
\noindent\textbf{Main abstract theorem.}
Assumptions \eqref{intro:measure-scaling} and
\eqref{intro:localization} imply that
$\widetilde\sigma_\alpha\to\sigma$ in the relevant $L^p$ spaces.  Combined
with the Berezin--Lieb inequalities, this gives the following two trace
formulas.

\smallskip
\noindent\textbf{\emph{Compact form.}}
Suppose that $X$ is compact, that each $\mathcal{H}_\alpha$ is finite-dimensional, and
that $c(\alpha)=\dim \mathcal{H}_\alpha$.  If either
\begin{enumerate}
  \item $\sigma$ is bounded and $\psi$ is continuous on
  $[\operatorname*{ess\,inf}\sigma,\operatorname*{ess\,sup}\sigma]$, or
  \item $\sigma\geq 0$, $\psi\in C([0,\infty))$, and
  $\lim_{t\to\infty}\psi(t)/t$ exists and is finite,
\end{enumerate}
then
\begin{equation}
  \frac{1}{c(\alpha)}\operatorname{Tr}\bigl(\psi(T_\sigma^\alpha)\bigr)
  \longrightarrow
  \int_X\psi(\sigma(x))\,d\nu(x).
  \label{intro:compact-szego}
\end{equation}

\smallskip
\noindent\textbf{\emph{Locally compact form.}}
Suppose that
$\sigma,\eta\in L^1(X,\nu)\cap L^2(X,\nu)$ are real-valued,
$\eta\geq0$, and $T_\sigma^\alpha$ and $T_\eta^\alpha$ commute for every
$\alpha$.  If either
\begin{enumerate}
  \item $\sigma$ is bounded and $\psi$ is continuous on
  $[\operatorname*{ess\,inf}\sigma,\operatorname*{ess\,sup}\sigma]$, or
  \item $\sigma\geq0$, $\psi\in C([0,\infty))$, and
  $\lim_{t\to\infty}\psi(t)/t$ exists and is finite,
\end{enumerate}
then
\begin{equation}
  \frac{1}{c(\alpha)}
  \operatorname{Tr}\bigl(T_\eta^\alpha\psi(T_\sigma^\alpha)\bigr)
  \longrightarrow
  \int_X\eta(x)\psi(\sigma(x))\,d\nu(x).
  \label{intro:local-szego}
\end{equation}

The point is that all geometric information is confined to verifying
\eqref{intro:localization}.  Once this is done, the operator-theoretic proof
of \eqref{intro:compact-szego} and \eqref{intro:local-szego} is identical
across the examples.  We now record representative consequences for compact
Abelian groups, weighted Bergman spaces, and locally compact Abelian groups.

Concretely, the framework yields new Szeg\H{o} limit theorems as well as unified proofs and extensions of several known ones. In particular, we obtain a Szeg\H{o}-type limit theorem of this form for Toeplitz operators on weighted Bergman spaces of the unit ball. We also extend the compact Abelian-group setting of B\'edos~\cite{bedos1996folner} to the class of symbols considered below, recover the recent trace theorem of Guo, Li, and Zhou~\cite{guo2024first} as a special case, and obtain a corresponding theorem for locally compact Abelian groups. Further applications recover first-order asymptotics for Gabor localization and Paley--Wiener operators and yield eigenvalue-distribution consequences related to classical results of Landau--Widom and Lindholm.

We now state three representative consequences of the abstract theory.





\subsection{Weighted Bergman spaces}
 Our main application is a Szeg\H{o} limit theorem for Toeplitz operators on weighted Bergman spaces and other locally compact settings. 
The weighted Bergman spaces, denoted $A_\alpha^2(\mathbb{B}^{n})$ with $\alpha > -1$, consist of holomorphic functions on the unit ball $\mathbb{B}^{n}$ of $\mathbb{C}^n$ that are square-integrable with respect to the measure $d\mu_\alpha(z)=\frac{\Gamma(n+\alpha+1)}{\pi^{n}\Gamma(\alpha+1)}(1-|z|^2)^\alpha \, dV(z)$, 
where $dV(z)$ is the volume measure on $\mathbb{C}^n$. 
It is well known (see, e.g., \cite{Z90}) that $A_\alpha^2(\mathbb{B}^{n})$ is a reproducing kernel Hilbert space (RKHS), 
and that Toeplitz operators on these spaces take the form:

\begin{align*}
    T_{\sigma}^{\alpha}f=\frac{\Gamma(n+\alpha+1)}{\pi^{n}\Gamma(\alpha+1)}\int_{\mathbb{B}^{n}}\sigma(z)\ip{f}{k_{z}^{\alpha}}k_{z}^{\alpha}(1-|z|^{2})^{-n-1}dV(z), 
\end{align*}
where $k_{z}^{\alpha}(w)$ are the kernel sections. We prove the following analog of \eqref{Szegocont}. 

% \begin{thm}\label{Sze1} Let $\sigma:\D\to \R$ be a strictly positive symbol such that $\sigma\in L^{p}(\mathbb{D},(1-|z|^{2})^{-2}dA(z))$, for some $p\in [1,\infty)$. Then, we have 
%     \begin{align*}
%         \frac{1}{\alpha+1}\tr\left(\log\left(T_{\sigma}^{\alpha}\right)\right)\to\int_{\mathbb{D}}\log\left(\sigma(z)\right)\frac{1}{(1-|z|^{2})^{2}}dA(z).
%     \end{align*}
% \end{thm}

\begin{thm}\label{Sze2} Let $\sigma:\mathbb{B}^{n}\to \R$ be such that either
\begin{itemize}
    \item[(a)] $\sigma$ is bounded, $\sigma\in L^{1}(\mathbb{B}^{n},(1-|z|^{2})^{-n-1}dV(z))$, and $\psi:\R \to \R$ is continuous on $[\mathrm{ess\,inf}\,\sigma,\mathrm{ess\,sup}\,\sigma]$, or 
    \item[(b)] $\sigma\geq0$ and
\[
\sigma\in
L^1\!\left(B_n,(1-|z|^2)^{-n-1}dV(z)\right)
\cap
L^2\!\left(B_n,(1-|z|^2)^{-n-1}dV(z)\right),
\]
and $\psi\in C([0,\infty))$ satisfies
\[
\lim_{x\to\infty}\frac{\psi(x)}{x}\in\mathbb{R}.
\]
\end{itemize}
Then, 
    \begin{align*}
        \frac{\pi^{n}\Gamma(\alpha+1)}{\Gamma(n+\alpha+1)}\tr\left(T_{\sigma}^{\alpha}\psi\left(T_{\sigma}^{\alpha}\right)\right)\to\int_{\mathbb{B}^{n}}\sigma(z)\psi\left(\sigma(z)\right)(1-|z|^{2})^{-n-1}dV(z),
    \end{align*}
    as $\alpha\to\infty$.
\end{thm}

\noindent Our analog of \eqref{Szegocont} has slightly different form due to the non-compactness of the underlying domain $\mathbb{B}^{n}$. 

A corresponding \(n=1\) result was established in the first arXiv
version of the present work, posted in November 2024. Subsequently,
Ram{\'i}rez Monta{\~n}o~\cite{RamirezMontano} proved a related
Szeg\H{o} limit theorem on \(\mathbb B^n\) for smooth compactly supported
symbols and, more generally, for symbols supported on isotropic or
co-isotropic submanifolds. In the absolutely continuous case, his result
overlaps with the present theorem only for smooth compactly supported
symbols. By contrast, Theorem~\ref{Sze2} allows bounded \(L^1\) symbols, and
non-negative \(L^1\cap L^2\) symbols, without smoothness or compact-support
assumptions. The results are therefore complementary: the present theorem
extends the classical-symbol result in regularity, while
\cite{RamirezMontano} also treats singular symbol measures.


\begin{comment}The analogous result for the Bargmann-Fock space also holds. 
It was essentially established by Feichtinger and Nowak in~\cite{FN01} within the broader context of time-frequency analysis. 
However, the method employed in~\cite{FN01}, which traces back to Widom~\cite{widom1979eigenvalue}, 
does not extend to the Bergman setting due to the exponential growth of the underlying hyperbolic measure. 
As a result, the subsequent work~\cite{DMN02}, which considers the Bergman setting, produces weaker results than the original paper of Nowak and Feichtinger.
\end{comment}

Besides providing a new result in this setting, the Bergman example illustrates the advantage of the abstract approach: the same proof applies despite the failure of the geometric assumptions underlying the standard Widom-type method.


\subsection{Compact Abelian groups}

\begin{thm}[Szeg\H{o} Limit on Compact Abelian Groups]
\label{thm:compact_ab_group}
Let $G$ be a compact metrizable Abelian group with normalized Haar measure $\mu$ and let $\{\Gamma_N\}$ be a sequence of finite subsets of $\widehat G$ satisfying the F\o lner condition
\[
\frac{|(\gamma\Gamma_N)\triangle\Gamma_N|}{|\Gamma_N|}
\longrightarrow0
\qquad \gamma\in\widehat G.
\]
Let $\mathcal{K}_{N}=\text{span}\{\Gamma_{N}\}$ and $\sigma \in L^1(G)$ be a real-valued, non-negative function. Let $T_\sigma^N : \mathcal{K}_N \to \mathcal{K}_N$ denote the corresponding Toeplitz operator
\[
T_\sigma^N f = P_N (\sigma f), \quad f \in \mathcal{K}_N,\]
where $P_N : L^2(G) \to \mathcal{K}_N$ is the orthogonal projection. Then, for any continuous function $\psi : [0,\infty) \to \mathbb{R}$ such that $\psi(x)/x$ has finite limit as $x \to \infty$, we have
\[
\lim_{N\to\infty} \frac{1}{|\Gamma_N|} \mathrm{Tr} \, \psi(T_\sigma^N) = \int_G \psi(\sigma(x)) \, d\mu(x).
\]
\end{thm}

\begin{comment}

This result should be compared with the theorem of B\'edos~\cite{bedos1996folner} and the recent work of Guo, Li, and Zhou ~\cite{guo2024first}. In particular, the trace theorem of ~\cite{guo2024first} for tori is obtained as a special case. Our emphasis is different: the compact group plays no essential role in the proof once the normalized kernels are known to form an approximate identity. The same abstract argument therefore applies to arbitrary compact Abelian groups and, as we show below, to settings having neither a group structure nor a complex structure.

\end{comment}

An important feature of our approach is that it also yields the
determinant formulation of the first Szeg\H{o} limit theorem.  Although
the logarithm is not covered directly by our continuous functional
calculus results because of its singularity at the origin, the
determinant asymptotic follows from the same abstract framework by a
regularization argument.  More precisely, for a non-negative
integrable symbol $\sigma$ in the compact setting, we obtain
\[ \lim_{N\to\infty} \frac{1}{|\Gamma_N|} \log\det T_\sigma^N = \int_G \log\sigma(x)\,d\mu(x), \]
where the right-hand side, and hence the limit, is allowed to take the
value $-\infty$.  Thus no assumption that the symbol be bounded away
from zero is required.  In particular, our abstract framework yields
both the trace and determinant formulations of the first Szeg\H{o}
limit theorem.



%As an additional benefit of our abstract Szeg\H{o} limit theorems, we obtain abstract Weyl laws of the form:
% \indent The main appeal of studying these limits is that the asymptotic behavior of the eigen-counting functions can be obtained as a biproduct. These are known as the Weyl laws and have been well studied for the Laplace-Beltrami Operator and other pseudo-differential operators \cite{BZ71,Geis14,EM07, guille79}. These results have also been considered in the context of Toeplitz operators. For example, Berezin \cite{BZ71} obtained asymptotics on the semiclassical eigen-counting function $N_{h}(\lambda)$ as $h\to 0$ in terms of the distribution of the symbol of a class of Toeplitz operators on the weighted Fock space. Using $(1.2)$, Feichtinger and Nowak \cite{FN01} obtained 
% \begin{align*}\label{mylabel}
%     \lim_{R\to\infty}\frac{1}{R^{2n}}N_{T_{b_{R}}}(\lambda)=\left|\{b(\eta)>\lambda\}\right| \tag{1.5}
% \end{align*}
% as $R\to\infty$, where $|\cdot|$ denotes the $n$-dimensional Lebesgue measure and $N_{R}(\lambda)$ is the eigen-counting function of $T_{\sigma_{R}}$. Using this Weyl law, they obtain estimates for the plunge regions of Gabor-Toeplitz and Calder\'on-Toeplitz operators (see \cite{DMN02,DMFN02,FN01}). Similarly, we will obtain Weyl laws of the following form: 
%\begin{align*}  \lim_{\alpha\to\infty}\frac{1}{c(\alpha)}N_{\alpha}(\lambda)=\nu\left(\{\sigma>\lambda\}\right).\end{align*}


% similarities between the operators in $(1.3)$ and $(1.4)$ are apparent: essentially there is a collection of vectors indexed over a parameter, with an associated indexed measure. Given the results obtained in the setting of Gabor-Toeplitz operators, it would be of interest to know if similar Szeg\H o limits can be obtained in the setting of weighted Bergman spaces. The authors currently do not know of any such results obtained in the setting of the weighted Bergman spaces. This question has been considered in a variety of settings, for example see \cite{GVas04,EU20} when considering the Bargman-Fock space. In particular, the case when the symbol is integrable in some sense was addressed in \cite{FN01,EU20} in a variety of spaces, and the problem has been considered for multiplication operators in \cite{janssen1983szegHo, laptev1996szego} for example.\\
%The usual strategy for obtaining such limits was developed by Widom \cite{widom1979eigenvalue}, where one is concerned with projections of multiplication operators onto subspaces of $L^{2}$ generated by eigen-functions of the Laplace-Beltrami operator. This strategy requires a growth estimate on the Hilbert-Schmidt norm of the Hankel operator $P_{n}A(I-P_{N})$ or ultimately the compactness/boundedness of the commutator $[-\Delta,A]$, and has been employed by several authors to obtain similar limits (see \cite{janssen1983szegHo,laptev1996szego} for example). In fact, this approach was slightly modified and used by Feichtinger and Nowak \cite{FN01} as well by considering Gabor-Toeplitz operators as projections of the multiplication operator. \\
%In this paper, we will take a different and more simplistic approach to obtain these limits than the strategy established by Widom \cite{widom1979eigenvalue}. This approach avoids the need to compute the Hilbert-Schmidt norm of the associated Hankel operator by instead checking a simple approximate identity condition on the kernel of the framed Hilbert space on which our Toeplitz operator acts. The approximate identity condition, denoted Assumption (II) below, will then imply suitable convergence of the Berezin transform \cite{BZ72} and through a squeezing argument imply the limits which we discuss more in Section 2. Additionally, this argument allows us to consider the problem in a much more general setting and will demonstrate that both the classical Szeg\H{o} Limit Theorem and the theorems of Feichtinger and Nowak \cite{FN01} are actually part of the same theory. \\

\subsection{Locally compact Abelian groups}

The same framework also produces a non-compact analogue of the compact-group theorem.

\begin{thm}[Szeg\H{o} theorem for LCA groups]\label{thm:szego_lca}
Let $G$ be a locally compact Abelian group with Haar measure $m$
and dual group $\widehat{G}$ with Haar measure $\widehat{m}$.
Let $\{F_N\}_{N\ge1}\subset \widehat{G}$ be a F\o lner sequence and define
\[
\mathcal{H}_N
=
\{f\in L^2(G): \operatorname{supp}(\widehat{f})\subset F_N\}.
\]
Let $c(N)=\widehat{m}(F_N)$.

For $\sigma\in L^1(G,m)$ real-valued,
let $T_\sigma^N$ be the Toeplitz operator
\[
T_\sigma^N f = P_N(\sigma f).
\]

If either

\begin{itemize}
\item[(a)] $\sigma$ is bounded, $\sigma\in L^{1}(G)$, and
$\psi$ is continuous on
$[\mathrm{ess\,inf}\,\sigma,\mathrm{ess\,sup}\,\sigma]$, or
\item[(b)] $\sigma\ge0$, $\sigma\in L^{2}(G)$, and
$\psi\in C([0,\infty))$ such that the limit
$\lim_{x\to\infty}\psi(x)/x$ exists,
\end{itemize}

then
\[
\lim_{N\to\infty}
\frac{1}{c(N)}
\operatorname{tr}\!\big(
T_\sigma^N \psi(T_\sigma^N)
\big)
=
\int_G \sigma(x)\psi(\sigma(x))\, dm(x).
\]
\end{thm}

This may be regarded as a locally compact counterpart of the compact Abelian-group Szeg\H{o} theorem. More importantly for our purposes, both results arise from exactly the same abstract argument.

\begin{comment}
\subsection{The abstract mechanism} We now briefly describe the abstract mechanism behind our results. Let \[ \{(H_\alpha,\langle\cdot,\cdot\rangle_\alpha):\alpha>0\} \] be a family of Hilbert spaces, let $X$ be a locally compact Hausdorff space, and suppose that, for each $\alpha$, we are given a normalized continuous Parseval frame \[ \{k_x^\alpha:x\in X\}\subset H_\alpha \] with associated Radon measure $\nu_\alpha$. Thus \[ f=\int_X \langle f,k_x^\alpha\rangle_\alpha k_x^\alpha \,d\nu_\alpha(x), \qquad f\in H_\alpha, \] where the integral is understood in the weak sense. Given a real-valued symbol $\sigma$ on $X$, we define the associated \emph{Toeplitz operator} by \[ T_\sigma^\alpha f = \int_X \sigma(x)\langle f,k_x^\alpha\rangle_\alpha k_x^\alpha \,d\nu_\alpha(x), \] on its natural domain. This definition includes the usual Toeplitz operators arising in the examples considered below, including compressions of multiplication operators, Bergman Toeplitz operators, and time--frequency localization operators. The Berezin transform of $T_\sigma^\alpha$ is \[ \widetilde{\sigma}_\alpha(x) = \langle T_\sigma^\alpha k_x^\alpha,k_x^\alpha\rangle_\alpha = \int_X \sigma(y) \left|\langle k_x^\alpha,k_y^\alpha\rangle_\alpha\right|^2 \,d\nu_\alpha(y). \] Our abstract Szeg\H{o} theorem is based on two assumptions. First, the frame measures have the form \[ d\nu_\alpha(x)=c(\alpha)\,d\nu(x), \] where $\nu$ is independent of $\alpha$. Second, the normalized frame vectors become asymptotically localized: for every $x\in X$ and every neighborhood $U$ of $x$, \[ \int_{U^c} \left|\langle k_x^\alpha,k_y^\alpha\rangle_\alpha\right|^2 \,d\nu_\alpha(y) \longrightarrow 0 \qquad (\alpha\to\infty). \] The second condition says that the kernels \[ \left|\langle k_x^\alpha,k_y^\alpha\rangle_\alpha\right|^2 \] form an approximate identity. Its principal consequence is convergence of the Berezin transforms, \[ \widetilde{\sigma}_\alpha\longrightarrow\sigma, \] in $L^p$, $1\leq p<\infty$, for the appropriate class of symbols. Combining this convergence with Berezin--Lieb inequalities yields the Szeg\H{o} trace asymptotics.

The important point is that all geometric information is isolated in the verification of the approximate-identity condition. Once this condition has been established, the remainder of the proof is independent of the particular realization of the Toeplitz operators. The same argument therefore applies to compact Abelian groups, weighted Bergman spaces, locally compact Abelian groups, Gabor localization operators, and Paley--Wiener spaces. In this sense, the abstract theorem provides a ``proof by verification'' principle: to establish a first-order Szeg\H{o} theorem in a new continuous-frame setting, it suffices to identify the natural scaling $c(\alpha)$, realize the operators as Toeplitz operators associated with the frame, and verify the localization property of the normalized frame vectors.

\end{comment}

\subsection{Relation to previous work and scope}

The results above are first-order trace formulae.  Their purpose is to isolate a localization principle that produces the leading term in several different settings; they are not intended to subsume the finer, model-dependent asymptotics available in all of those settings.

For compact abelian groups, B\'edos~\cite{bedos1996folner} proved a Szeg\H{o}-type trace formula for bounded real-valued symbols along F{\o}lner sequences. For \(G=\mathbb T^d\), \(1\leq d\leq\infty\), Guo, Li, and Zhou~\cite{guo2024first} extended the trace formula to nonnegative \(L^1\)-symbols and continuous test functions of at most linear growth (motivated in part by questions of Nikolski and Pushnitski~\cite{nikolski2020szeg}); they also obtained determinant formulae and studied certain non-F{\o}lner
truncations. Our Theorem~\ref{thm:compact_ab_group} and Corollary~\ref{DetGroup} recover the trace and determinant part of their results in the torus case and extend the underlying compact group to an arbitrary compact
metrizable abelian group.  It does not address their 
non-F{\o}lner results, so those results are complementary to ours.

The comparison with the weighted Bergman theorem is of a different nature.
Ramirez Monta{\~n}o~\cite{RamirezMontano} treats positive smooth,
compactly supported symbols associated with isotropic and co-isotropic
submanifolds; in particular, his setting includes geometrically defined
singular measure symbols.  By contrast, Theorem~\ref{Sze2} concerns absolutely
continuous symbols and permits the rough symbol classes stated there, with
no smoothness or compact-support assumption.  Thus, in the common case of a
positive smooth compactly supported density, the two results give compatible
first-order formulae after the normalizations are matched.  Neither result
contains the other: the former allows singular geometric symbols, whereas
the latter allows rough diffuse symbols.

In the Gabor setting, the theorem of Feichtinger and
Nowak~\cite{FN01} assumes \(b\in L^1\) with \(0\leq b\leq1\).
Since this implies \(b\in L^1\cap L^2\), Corollary~\ref{FeiNow} recovers that
first-order asymptotic and permits the broader nonnegative symbol class
covered by our local theorem, together with the corresponding
linear-growth class of test functions.  Likewise, the Paley--Wiener
specialization recovers the leading Landau--Widom counting law
\cite{landau1980eigenvalue}.  It does not recover the logarithmic second-order term,
which depends on boundary information beyond the localization criterion used
here.


\begin{comment}
\subsection{Comparison with previous work}
The abstract point of view clarifies the relation between the present results and several earlier theorems. B\'edos~\cite{bedos1996folner} developed Szeg\H{o} and eigenvalue-distribution results using F\o lner approximations. Our compact Abelian-group theorem recovers this phenomenon in the continuous-frame language and extends it to unbounded symbol classes appearing above, while the locally compact theorem shows that the same mechanism persists beyond the compact setting.
Guo, Li, and Zhou~\cite{guo2024first} recently proved the first Szeg\H{o} limit theorem on multidimensional and infinitedimensional tori, motivated in part by questions of Nikolski and Pushnitski~\cite{nikolski2020szeg}. Their work also contains determinant-type results and an analysis of non-F\o lner approximations. Our framework addresses a different aspect of the problem: their trace asymptotic becomes a special case of a theorem in which neither the torus nor its Fourier structure is essential. In particular, the same argument applies to arbitrary compact Abelian groups and to non-group settings such as weighted Bergman spaces.

In time-frequency analysis, Feichtinger and Nowak~\cite{FN01}, following ideas of Widom, proved a Szeg\H{o}-type theorem for Gabor--Toeplitz localization operators. Their method exploits localization properties of the Euclidean phase space and does not extend directly to the Bergman geometry considered here. Our theorem recovers their first-order asymptotic while at the same time applying to the Bergman setting, where the invariant measure has exponential volume growth. Similarly, the Paley--Wiener application recovers the leading term in the classical Landau--Widom asymptotics; the second-order boundary term in the latter theorem is, of course, beyond the scope of our first-order abstract method.
Finally, the eigenvalue-distribution consequences of the abstract theorem yield further classical applications, including a Wiener-type limit on locally compact Abelian groups and the necessary density condition for sampling in the Bargmann--Fock space.
Thus, the principal contribution of the paper is not merely another instance of a Szeg\H{o} limit theorem. Rather, the results show that a large family of apparently different first-order Szeg\H{o} asymptotics is governed by the same simple localization mechanism. The weighted Bergman theorem provides a particularly useful illustration: the abstract method continues to apply precisely in a setting where the standard geometric localization approach encounters a genuine obstruction.

\end{comment}

\begin{comment}
We now introduce our abstract setting. Let $\alpha > 0$ be an indexing parameter and let 
$\{ (\mathcal{H}_\alpha, \langle \cdot, \cdot \rangle_\alpha) : \alpha > 0 \}$ be a family of Hilbert spaces indexed by $\alpha$. 
Let $X$ be a locally compact topological space, and let $\{ \nu_\alpha \}_{\alpha > 0}$ be a collection of Radon measures on $X$. 
We always assume that the measures $\nu_\alpha$ satisfy the criterion:
\begin{itemize}
    \item \textbf{Assumption (I).} $d\nu_\alpha(x) = c(\alpha) \, d\nu(x)$, where $\nu$ is a Radon measure on $X$ and $c(\alpha) > 0$ depends only on $\alpha$.
\end{itemize}

For each fixed $\alpha$, we assume the existence of a family of vectors $\{ k_x^\alpha \}_{x \in X} \subset \mathcal{H}_\alpha$ such that $\| k_x^\alpha \| = 1$ for all $x \in X$, and
\begin{align*}
    f = \int_X \langle f, k_x^\alpha \rangle_\alpha \, k_x^\alpha \, d\nu_\alpha(x), \quad \text{for all } f \in \mathcal{H}_\alpha.
\end{align*}
Here, the integral is understood in the weak sense; that is, for any $f, g \in \mathcal{H}_\alpha$, we have
\begin{align*}
    \langle f, g \rangle_\alpha = \int_X \langle f, k_x^\alpha \rangle_\alpha \, \langle k_x^\alpha, g \rangle_\alpha \, d\nu_\alpha(x).
\end{align*}
A family of vectors $\{ k_x^\alpha \}_{x \in X}$ that satisfy the above conditions is called a \emph{normalized continuous Parseval frame} for $\mathcal{H}_\alpha$.
% \begin{itemize}
%     \item[(I)] $d\nu_{\alpha}(x)=c(\alpha)d\nu(x)$, where $\nu$ is a Radon measure on $X$ and $c(\alpha)>0$ depends only on $\alpha$.
%     \item[(II)] For any $x\in X$ and open neighborhood $N_{x}$ of $x$,
%     \begin{align*}
%         \int_{N_{x}^{c}}\left|\ip{k_{x}^{\alpha}}{k_{y}^{\alpha}}_{\alpha}\right|^{2}d\nu_{\alpha}(y)\to 0,
%     \end{align*}
%     as $\alpha\to\infty$. 
% \end{itemize}
%  The first condition is a separability condition which gives us the associated scaling constant for our limit. The second condition is essentially a statement saying that $c(\alpha)\left|\ip{k_{x}^{\alpha}}{k_{y}^{\alpha}}\right|^{2}$ is an approximate identity. 
We now define the Toeplitz operators we will consider going forward.
\begin{definition}
    Let $\sigma\in L^{p}(X,\nu_{\alpha})$ for some $1\leq p\leq\infty$, with $\sigma:X\to\R$. If $p<\infty$, we further assume that for every $x\in X$ \[ k_{x}^{\alpha}\in \{f\in\mathcal{H}_{\alpha}:\int_{X}|\sigma(y)|\left|\ip{f}{k_{y}^{\alpha}}\right|^{2}d\nu_{\alpha}(y)<\infty\}=:\mathcal{D}(T^{\alpha}_{\sigma}).\] Then the \emph{Toeplitz operator} $T_{\sigma}^{\alpha}:\mathcal{D}(T_{\sigma}^{\alpha})\to\mathcal{H}_{\alpha}$ is defined by
\begin{align*}
    T_{\sigma}^{\alpha}f=\int_{X}\sigma(x)\ip{f}{k_{x}^{\alpha}}_{\alpha}k_{x}^{\alpha}d\nu_{\alpha}(x),
\end{align*}
where this integral is understood in the same weak sense as above. 
\end{definition}
\indent Clearly, each Toeplitz operator $T_{\sigma}^{\alpha}$ is a densely defined symmetric operator, which is bounded if (but not necessarily only if) $\sigma$ is bounded. Moreover, if $\sigma\in L^1(X,\nu_\alpha)$, then $T_{\sigma}^{\alpha}:\mathcal{H}_\alpha\to\mathcal{H}_{\alpha}$ is of trace class. An important transformation associated to a Toeplitz operator is its $\alpha$-Berezin transform~\cite{BZ72}.
\begin{definition}
    For $\sigma$ satisfying the above assumptions, the \emph{Berezin transform} of $T_{\sigma}^{\alpha}$ is the function $\Tilde{\sigma}^{\alpha}:X\to\mathbb{R}$ defined by 
\begin{align*}
    \Tilde{\sigma}^{\alpha}(x)=\ip{T_{\sigma}^{\alpha}k_{x}^{\alpha}}{k_{x}^{\alpha}}_{\alpha}=\int_{X}\sigma(y)\left|\ip{k_{x}^{\alpha}}{k_{y}^{\alpha}}_{\alpha}\right|^{2}d\nu_{\alpha}(y).
\end{align*}
\end{definition}
\end{comment}

\begin{comment}
We now introduce the class of Toeplitz operators in our abstract setting. We restrict attention to Toeplitz operators with real-valued integrable symbols.

\begin{definition}
Let $\sigma \in L^1(X,\nu_\alpha)$ be real-valued. 
The \emph{Toeplitz operator} with symbol $\sigma$ is the operator 
$T_\sigma^\alpha : \mathcal H_\alpha \to \mathcal H_\alpha$ defined weakly by
\begin{align*}
    \langle T_\sigma^\alpha f , g \rangle_\alpha
    =
    \int_X \sigma(x)\,
    \langle f, k_x^\alpha \rangle_\alpha
    \overline{\langle g, k_x^\alpha \rangle_\alpha}
    \, d\nu_\alpha(x),
    \qquad f,g \in \mathcal H_\alpha .
\end{align*}
Equivalently,
\begin{align*}
    T_\sigma^\alpha f
    =
    \int_X \sigma(x)\,
    \langle f, k_x^\alpha \rangle_\alpha
    k_x^\alpha \, d\nu_\alpha(x),
\end{align*}
where the integral is understood in the weak sense.
\end{definition}

Since $\sigma \in L^1(X,\nu_\alpha)$, the operator $T_\sigma^\alpha$ is well defined on all of 
$\mathcal H_\alpha$ and is trace class. 
Moreover, because $\sigma$ is real-valued, $T_\sigma^\alpha$ is symmetric. 
Its trace is given by
\begin{align*}
    \operatorname{Tr}(T_\sigma^\alpha)
    =
    \int_X \sigma(x)\, d\nu_\alpha(x).
\end{align*}

An important associated object is the $\alpha$-Berezin transform.

\begin{definition}
The \emph{Berezin transform} of $T_\sigma^\alpha$ is the function 
$\widetilde{\sigma}^\alpha : X \to \mathbb R$ defined by
\begin{align*}
    \widetilde{\sigma}^\alpha(x)
    =
    \langle T_\sigma^\alpha k_x^\alpha , k_x^\alpha \rangle_\alpha
    =
    \int_X \sigma(y)\,
    \left| \langle k_x^\alpha , k_y^\alpha \rangle_\alpha \right|^2
    \, d\nu_\alpha(y).
\end{align*}
\end{definition}
    
\end{comment}
\begin{comment}
The main goal of our paper is to prove Szeg\H o-type limit results in the above general abstract setting, making distinctions when the underlying phase space $X$ is locally compact with potentially infinite measure $\nu$, and when $X$ is compact with finite measure $\nu$. Our main result in the locally compact setting will be of the following form: for $\sigma:X\to\mathbb{R}$ integrable, $\eta:X\to\mathbb{R}$ non-negative and integrable, with $T_{\eta}^{\alpha}$ and $T_{\sigma}^{\alpha}$ commuting, and $\psi: [0,\infty)\to\mathbb{R}$ continuous with $\psi(x)/x\to \beta$ as $x\to\infty$ for some $\beta\in \mathbb{R}$, we have 
\begin{align*}
    \lim_{\alpha\to\infty}\frac{1}{c(\alpha)}\text{tr}\left(T_{\eta}^{\alpha}\psi\left(T_{\sigma}^{\alpha}\right)\right)=\int_{X}\eta(x)\psi\left(\sigma(x)\right)d\nu(x).
\end{align*}
In the case when $\sigma$ is bounded, $\psi$ can be taken to be any function continuous on the range of $\sigma$.

In the compact setting, the form is more classical: for $\sigma:X\to \R$  non-negative and integrable, and $\psi: [0,\infty)\to\mathbb{R}$ continuous with $\psi(x)/x\to \beta$ as $x\to\infty$ for some $\beta\in \mathbb{R}$, we have 
\begin{align*}
    \lim_{\alpha\to\infty}\frac{1}{c(\alpha)}\text{tr}\left(\psi(T_{\sigma}^{\alpha})\right)=\int_{X}\psi(\sigma(x))d\nu(x),
\end{align*}
 In the case when $\sigma$ is bounded, $\psi$ can be taken to be any function continuous on the range of $\sigma$.
\end{comment}

%The flow of the paper is as follows. In Section 2, we will present the motivation for our proofs and a variant of the Berezin-Lieb inequality which will be used throughout. In Section 3, we will prove key technical lemmas concerning the convergence of the $\alpha$-Berezin transform, a condition which is required in most of our abstract results. In Section 4, we will present and prove our abstract Szeg\H{o} limit theorems in the compact phase space setting, followed by a presentation of examples in this setting. Finally, in Section 5 we will prove variants of the Szeg\H{o} limit theorem in the locally compact setting, and derive several examples.  
% In Section 2, we will prove a general Szeg\H{o} limit theorem for Toeplitz operators, for which the convergence of the Berezin transform will depend. In Section 3, we prove the convergence of the associated Berezin transform for integrable symbols and derive various corollaries to the result of Section 2. In Section 4, we will obtain a Weyl law for integrable symbols, along with estimates of the associated plunge regions. Finally, Section 5 will present examples of both semiclassical function spaces where our assumptions hold, while also demonstrating that the classical Szeg\H{o} Limit Theorem is a corollary. We now make the following observation.
\begin{comment}
    

\section{Extended Berezin-Lieb Inequality}
 In the setting described above, the well-known Berezin-Lieb inequality \cite{BZ72} takes the following form:
\begin{align*}
    \int_{X}\psi\left(\Tilde{\sigma}^{\alpha}(x)\right)d\nu_{\alpha}(x)\leq\text{tr}\left(\psi\left(T_{\sigma}^{\alpha}\right)\right)\leq\int_{X}\psi\left(\sigma(x)\right)d\nu_{\alpha}(x),
\end{align*}
where $\psi$ is convex. Upon using Assumption (I), this inequality becomes
\begin{align*}
    \int_{X}\psi\left(\Tilde{\sigma}^{\alpha}(x)\right)d\nu(x)\leq\frac{1}{c(\alpha)}\text{tr}\left(\psi\left(T_{\sigma}^{\alpha}\right)\right)\leq\int_{X}\psi\left(\sigma(x)\right)d\nu(x).
\end{align*}
In particular, the upper bound is independent of $\alpha$. Hence, if one were to have $\Tilde{\sigma}^{\alpha}$ converge to $\sigma$ in some suitable manner, a squeezing argument along with an appropriate measure-theoretic convergence theorem could be applied to obtain an associated Szeg\H o limit. 
In fact, this argument can be extended further to include real-valued symbols using the following generalized Berezin-Lieb inequality. 
\begin{prop}
    Suppose $\sigma,\eta\in L^{\infty}(X,\nu)$ are real-valued, $\eta\geq 0$. Assume that the Toeplitz operators $T_{\eta}$ and $T_{\sigma}$ are simultaneously diagonalizable by the same orthonormal basis. If $\psi:\mathbb{R}\to\mathbb{R}$ is convex, then 
    \begin{align*}
        \int_{X}\psi\left(\Tilde{\sigma}(x)\right)\eta(x)d\nu(x)\leq\text{tr}\left(T_{\eta}\psi\left(T_{\sigma}\right)\right)\leq&\int_{X}\psi\left(\sigma(x)\right)\Tilde{\eta}(x)d\nu(x).
    \end{align*}
\end{prop}
\begin{proof}
     First, we prove the lower bound. Since $\sigma$ and $\eta$ are real-valued, $T_{\sigma}$ and $T_{\eta}$ are symmetric. By assumption, there is a common eigen-basis for $T_{\sigma}$ and $T_{\eta}$. Denote these eigenvectors as $\{\varphi_{n}\}$. By Parseval's identity and the definition of the Berezin transform, we have
    
    \begin{align*}
        \Tilde{\sigma}(x)=\ip{T_{\sigma}k_{x}}{k_{x}}=\sum_{n=1}^{\infty}\lambda_{n}(\sigma)\left|\ip{\varphi_{n}}{k_{x}}\right|^{2},
    \end{align*}
    where $\lambda_{n}(\sigma)$ denotes the eigenvalue of $T_{\sigma}$ with a corresponding eigenvector $\phi_n$. Now, note that by Jensen's inequality and an interchange of sums, we obtain
    \begin{align*}
        \int_{X}\psi\left(\Tilde{\sigma}(x)\right)\eta(x)d\nu(x)=&\int_{X}\psi\left(\sum_{n=1}^{\infty}\lambda_{n}(\sigma)\left|\ip{\varphi_{n}}{k_{x}}\right|^{2}\right)\eta(x)d\nu(x)\\
        \leq&\sum_{n=1}^{\infty}\psi(\lambda_{n}(\sigma))\int_{X}\eta(x)\left|\ip{\varphi_{n}}{k_{x}}\right|^{2}d\nu(x)\\
        =&\sum_{n=1}^{\infty}\psi\left(\lambda_{n}(\sigma)\right)\ip{T_{\eta}\varphi_{n}}{\varphi_{n}}\\
        =&\text{tr}\left(\psi\left(T_{\sigma}\right)T_{\eta}\right).
    \end{align*}
 For the lower bound, we have by similar means
    \begin{align*}
        \int_{X}\psi\left(\sigma(x)\right)\Tilde{\eta}(x)d\nu(x)=&\int_{X}\psi\left(\sigma(x)\right)\sum_{n=1}^{\infty}\lambda_{n}(\eta)\left|\ip{\varphi_{n}}{k_{x}}\right|^{2}d\nu(x)\\
        \geq&\sum_{n=1}^{\infty}\lambda_{n}(\eta)\psi\left(\ip{T_{\sigma}\varphi_{n}^{\alpha}}{\varphi_{n}^{\alpha}}\right)\\
        =&\sum_{n=1}^{\infty}\ip{T_{\eta}\varphi_{n}}{\varphi_{n}}\psi\left(\lambda_{n}(\sigma)\right),
    \end{align*}
    which completes the proof.  
\end{proof}
Note that in the above lemma also holds (but with a reverse of the bounds) if one assumes $\psi$ is concave.
\end{comment}
%%%%%%%%%%%%%%%%%%%%%%%%%%%%%%%%%%%%%%%%%%%

\section{Berezin--Lieb type inequalities}

We record two Berezin--Lieb type inequalities that will be used later. Let $\mathfrak{S}_{1}(\mathcal{H}_{\alpha})$ denote the trace-class operators.  

\subsection*{The classical inequality}

\begin{prop}\label{BL1}
Let $\sigma\in L^1(X,\nu_\alpha)$ be real-valued,
and let $\psi:\mathbb{R}\to\mathbb{R}$ be convex. Assume that
\[
\psi(T_\sigma^\alpha)\in\mathfrak{S}_1(\mathcal H_\alpha)
\qquad\text{and}\qquad
\psi(\sigma)\in L^1(X,\nu_\alpha).
\]
Then
\begin{equation}\label{BL-classical}
\int_X
\psi\bigl(\widetilde{\sigma}^{\,\alpha}(x)\bigr)
\,d\nu_\alpha(x)
\leq
\operatorname{Tr}\bigl(\psi(T_\sigma^\alpha)\bigr)
\leq
\int_X\psi\bigl(\sigma(x)\bigr)\,d\nu_\alpha(x).
\end{equation}
\end{prop}


\begin{proof} Since $\sigma\in L^1(X,\nu_\alpha)$, the operator
$T_\sigma^\alpha$ is trace class and therefore compact. Since
$\sigma$ is real-valued, it is self-adjoint. Hence there exists an
orthonormal eigenbasis $(\phi_j)_j$ of $\mathcal H_\alpha$, with
corresponding real eigenvalues $(\lambda_j)_j$.

Since the frame is normalized and Parseval,
\[
\sum_j
\bigl|\langle \phi_j,k_x^\alpha\rangle_\alpha\bigr|^2=1
\]
for every $x\in X$. Moreover,
\[
\widetilde{\sigma}^{\,\alpha}(x)
=
\left\langle T_\sigma^\alpha k_x^\alpha,k_x^\alpha
\right\rangle_\alpha
=
\sum_j \lambda_j
\bigl|\langle \phi_j,k_x^\alpha\rangle_\alpha\bigr|^2.
\]
Jensen's inequality therefore gives
\[
\psi\bigl(\widetilde{\sigma}^{\,\alpha}(x)\bigr)
\leq
\sum_j \psi(\lambda_j)
\bigl|\langle \phi_j,k_x^\alpha\rangle_\alpha\bigr|^2.
\]
Integrating and using the Parseval frame identity,
\[
\int_X
\psi\bigl(\widetilde{\sigma}^{\,\alpha}(x)\bigr)
\,d\nu_\alpha(x)
\leq
\sum_j\psi(\lambda_j)
=
\operatorname{Tr}\bigl(\psi(T_\sigma^\alpha)\bigr).
\]

For the reverse inequality, we have
\[
\lambda_j
=
\left\langle T_\sigma^\alpha\phi_j,\phi_j\right\rangle_\alpha
=
\int_X \sigma(x)
\bigl|\langle \phi_j,k_x^\alpha\rangle_\alpha\bigr|^2
\,d\nu_\alpha(x).
\]
Applying Jensen's inequality once more,
\[
\psi(\lambda_j)
\leq
\int_X \psi(\sigma(x))
\bigl|\langle \phi_j,k_x^\alpha\rangle_\alpha\bigr|^2
\,d\nu_\alpha(x).
\]
Summing over $j$ and using Parseval again yields
\[
\operatorname{Tr}\bigl(\psi(T_\sigma^\alpha)\bigr)
\leq
\int_X \psi(\sigma(x))\,d\nu_\alpha(x).
\]
This proves the result.
\end{proof}


Under Assumption~(I), dividing \eqref{BL-classical} by $c(\alpha)$ gives

\begin{equation}\label{renorm1}
\int_X
\psi\bigl(\widetilde{\sigma}^{\,\alpha}(x)\bigr)\,d\nu(x)
\leq
\frac{1}{c(\alpha)}
\operatorname{Tr}\bigl(\psi(T_\sigma^\alpha)\bigr)
\leq
\int_X\psi\bigl(\sigma(x)\bigr)\,d\nu(x).
\end{equation}

\subsection*{A weighted inequality}

\begin{prop}\label{BLweighted}
Let $\sigma,\eta\in L^1(X,\nu_\alpha)$, with
$\sigma$ real-valued and $\eta\geq 0$. Assume that
$T_\sigma^\alpha$ and $T_\eta^\alpha$ commute. Let
$\psi:\mathbb{R}\to\mathbb{R}$ be convex. Then

\begin{equation}\label{BL-weighted}
\int_X
\eta(x)\,
\psi\bigl(\widetilde{\sigma}^{\,\alpha}(x)\bigr)
\,d\nu_\alpha(x)
\leq
\operatorname{Tr}\!\left(
T_\eta^\alpha\psi(T_\sigma^\alpha)
\right)
\leq
\int_X
\widetilde{\eta}^{\,\alpha}(x)\,
\psi\bigl(\sigma(x)\bigr)
\,d\nu_\alpha(x),
\end{equation}
whenever the right-hand side is finite.
\end{prop}

\begin{proof}
Since $\sigma,\eta\in L^1(X,\nu_\alpha)$, the operators
$T_\sigma^\alpha$ and $T_\eta^\alpha$ are trace class. Moreover,
$T_\sigma^\alpha$ is self-adjoint and $T_\eta^\alpha$ is positive.
Since they commute, they admit a common orthonormal eigenbasis
$(\phi_j)_j$, with
\[
T_\sigma^\alpha\phi_j=\lambda_j\phi_j,
\qquad
T_\eta^\alpha\phi_j=\mu_j\phi_j,
\qquad
\mu_j\geq 0.
\]
Thus
\[
\operatorname{Tr}\!\left(
T_\eta^\alpha\psi(T_\sigma^\alpha)
\right)
=
\sum_j\mu_j\psi(\lambda_j).
\]

For each $x\in X$, Parseval's identity gives
\[
\widetilde{\sigma}^{\,\alpha}(x)
=
\sum_j\lambda_j
\bigl|\langle\phi_j,k_x^\alpha\rangle_\alpha\bigr|^2,
\]
and
\[
\sum_j
\bigl|\langle\phi_j,k_x^\alpha\rangle_\alpha\bigr|^2
=1.
\]
Therefore, Jensen's inequality implies
\[
\psi\bigl(\widetilde{\sigma}^{\,\alpha}(x)\bigr)
\leq
\sum_j\psi(\lambda_j)
\bigl|\langle\phi_j,k_x^\alpha\rangle_\alpha\bigr|^2.
\]
Multiplying by $\eta(x)$ and integrating gives
\[
\begin{aligned}
\int_X\eta(x)
\psi\bigl(\widetilde{\sigma}^{\,\alpha}(x)\bigr)
\,d\nu_\alpha(x)
&\leq
\sum_j\psi(\lambda_j)
\int_X\eta(x)
\bigl|\langle\phi_j,k_x^\alpha\rangle_\alpha\bigr|^2
\,d\nu_\alpha(x)\\
&=
\sum_j\mu_j\psi(\lambda_j).
\end{aligned}
\]

For the reverse inequality, we have
\[
\lambda_j
=
\int_X\sigma(x)
\bigl|\langle\phi_j,k_x^\alpha\rangle_\alpha\bigr|^2
\,d\nu_\alpha(x).
\]
Another application of Jensen's inequality yields
\[
\psi(\lambda_j)
\leq
\int_X\psi(\sigma(x))
\bigl|\langle\phi_j,k_x^\alpha\rangle_\alpha\bigr|^2
\,d\nu_\alpha(x).
\]
Multiplying by $\mu_j$ and summing over $j$, we obtain
\[
\begin{aligned}
\sum_j\mu_j\psi(\lambda_j)
&\leq
\int_X\psi(\sigma(x))
\sum_j\mu_j
\bigl|\langle\phi_j,k_x^\alpha\rangle_\alpha\bigr|^2
\,d\nu_\alpha(x)\\
&=
\int_X
\widetilde{\eta}^{\,\alpha}(x)
\psi(\sigma(x))\,d\nu_\alpha(x).
\end{aligned}
\]
This proves the result.
\end{proof}

Dividing by $c(\alpha)$ gives
\begin{equation}\label{renorm2}
\int_X
\eta(x)\,
\psi\bigl(\widetilde{\sigma}^{\,\alpha}(x)\bigr)
\,d\nu(x)
\leq
\frac{1}{c(\alpha)}
\operatorname{Tr}\!\left(
T_\eta^\alpha\psi(T_\sigma^\alpha)
\right)
\leq
\int_X
\widetilde{\eta}^{\,\alpha}(x)\,
\psi\bigl(\sigma(x)\bigr)
\,d\nu(x).
\end{equation}

\medskip

If $\psi$ is concave, the inequalities in
\eqref{BL-classical} and \eqref{BL-weighted} reverse.

\medskip

\noindent
Both inequalities are particularly useful after normalization.
If $\widetilde{\sigma}^\alpha$ converges to $\sigma$
in a suitable sense and the family
$\{\widetilde{\eta}^\alpha\}$ admits an analogous limit,
then \eqref{renorm1} and \eqref{renorm2}
yield Szeg\H{o}-type trace asymptotics by a standard squeezing argument
combined with an appropriate convergence theorem.
%%%%%%%%%%%%%%%%%%%%%%%%%%%%%%%%%%%%%%%%%%%%%%%%%%%%%%
\begin{comment}
    

\section{Convergence of the Berezin Transform}
\indent The purpose of this section will be to prove that under a simple approximate identity condition, the convergence in measure of the Berezin transform needed for the Szeg\H{o} limit theorems of Sections 4 and 5 are guaranteed. This condition, which we refer to as Assumption (II), is the following: 
\begin{itemize}
    \item \textbf{Assumption (II)} for every $x\in X$, and every neighborhood $N_{x}$ of $x$, 
\begin{align*}
    \int_{N_{x}^{c}}\left|\ip{k_{x}^{\alpha}}{k_{y}^{\alpha}}\right|^{2}d\nu_{\alpha}(x)\to 0
\end{align*}
as $\alpha\to\infty$. 
\end{itemize}

Before considering this, we note the following elementary boundedness property of the Berezin transform. 
\begin{lem}
    The $\alpha$-Berezin transform $\sim^{\alpha}:L^{p}(X,\nu)\to L^{p}(X,\nu)$ for $1\leq p\leq\infty$ is bounded uniformly in $\alpha$.
\end{lem}
\begin{proof}
    The fact that the map is linear trivially follows from the definition of the Berezin transform. Thus, we proceed to the boundedness. Let $\sigma\in L^{\infty}(X,\nu)$. Then, 
    \begin{align*}
        \norm{\Tilde{\sigma}^{\alpha}}_{\infty}=\sup_{x\in X}\left|\int_{X}\sigma(y)\left|\ip{k_{x}^{\alpha}}{k_{y}^{\alpha}}_{\alpha}\right|^{2}d\nu_{\alpha}(y)\right|\leq&\sup_{x\in X}\int_{X}\left|\sigma(y)\right|\left|\ip{k_{x}^{\alpha}}{k_{y}^{\alpha}}_{\alpha}\right|^{2}d\nu_{\alpha}(y)\\
        \leq&\norm{\sigma}_{\infty}\sup_{x\in X}\int_{X}\left|\ip{k_{x}^{\alpha}}{k_{y}^{\alpha}}_{\alpha}\right|^{2}d\nu_{\alpha}(y),
    \end{align*}
    from which we obtain $\norm{\sim^{\alpha}}_{L^{\infty}\to L^{\infty}}\leq 1$, since the above integral is $1$ uniformly in $(x,\alpha)$. Now, let $\sigma\in L^{1}(X,\nu)$. Then, by Fubini's Theorem and Assumption (I)
    \begin{align*}
        \norm{\Tilde{\sigma}^{\alpha}}_{1}=\int_{X}\left|\tilde{\sigma}^{\alpha}(x)\right|d\nu(x)\leq&\int_{X}\int_{X}\left|\sigma(y)\right|\left|\ip{k_{x}^{\alpha}}{k_{y}^{\alpha}}_{\alpha}\right|^{2}d\nu_{\alpha}(y)d\nu(x)\\
        =&\int_{X}\int_{X}\left|\sigma(y)\right|\left|\ip{k_{x}^{\alpha}}{k_{y}^{\alpha}}_{\alpha}\right|^{2}d\nu_{\alpha}(x)d\nu(y)\\
        =&\int_{X}\left|\sigma(y)\right|d\nu(y),
    \end{align*}
    from which we obtain $\norm{\sim^{\alpha}}_{L^{1}\to L^{1}}\leq 1$. Thus, the result follows by an application of the Riesz-Thorin Interpolation Theorem \cite{G10}. 
\end{proof}
We will hold off on applying Lemma $3.1$ until Proposition $3.3$. We will first show that the Berezin transform converges pointwise to the symbol for the set of continuous, compactly supported functions on $X$, $C_{c}(X)$. 
%In many of our example spaces, this class has been shown to converge; see \cite{Z90,Z12} for more details.
The assumption that our underlying measure is Radon ensures that these test functions are dense in our $L^{p}$ spaces. Then, applying a density argument will give $L^{p}$ convergence for $L^{p}$ symbols, which in turn will imply convergence in measure. Below, we use a slightly non-standard approach to $L^p$ convergence which at first sight might seem more convoluted than necessary. The reason is that we don't want to assume any group homogeneity on our space $X$. If $X$ is a homogeneous space, then the $L^p$ convergence can be proved in a standard, more direct, way. 

        \begin{lem}
            If $\sigma\in C_{b}(X)$, then $\Tilde{\sigma}^{\alpha}\to\sigma$ pointwise.
        \end{lem}
        \begin{proof}
            Let $x\in X$. By continuity, for $\varepsilon>0$, there is a neighborhood $N_{x}$ such that 
                \begin{align*}
                    \left|\sigma(y)-\sigma(x)\right|<\varepsilon/2,
                \end{align*}
                for each $y\in N_{x}$. By Assumption (II), we have that for $\alpha>0$ sufficiently large that 
                \begin{align*}
                    \int_{N_{x}^{c}}\left|\ip{k_{x}^{\alpha}}{k_{y}^{\alpha}}_{\alpha}\right|^{2}d\nu_{\alpha}(y)<\varepsilon/(2\norm{\sigma}_{\infty}).
                \end{align*}
                Now, note that 
                \begin{align*}
                    \left|\Tilde{\sigma}^{\alpha}(x)-\sigma(x)\right|=&\left|\int_{X}\left(\sigma(y)-\sigma(x)\right)\left|\ip{k_{x}^{\alpha}}{k_{y}^{\alpha}}_{\alpha}\right|^{2}d\nu_{\alpha}(y)\right|\\
                    \leq&\int_{X}\left|\sigma(y)-\sigma(x)\right|\left|\ip{k_{x}^{\alpha}}{k_{y}^{\alpha}}_{\alpha}\right|^{2}d\nu_{\alpha}(y)\\
                    =&\int_{N_{x}}\left|\sigma(y)-\sigma(x)\right|\left|\ip{k_{x}^{\alpha}}{k_{y}^{\alpha}}_{\alpha}\right|^{2}d\nu_{\alpha}(y)\\
                    &\quad+\int_{N_{x}^{c}}\left|\sigma(y)-\sigma(x)\right|\left|\ip{k_{x}^{\alpha}}{k_{y}^{\alpha}}_{\alpha}\right|^{2}d\nu_{\alpha}(y).
                \end{align*}
                Using the estimate on $\sigma$ on $N_{x}$ gives
                \begin{align*}
                     \left|\Tilde{\sigma}^{\alpha}(x)-\sigma(x)\right|\leq&(\varepsilon/2)\int_{N_{x}}\left|\ip{k_{x}^{\alpha}}{k_{y}^{\alpha}}_{\alpha}\right|^{2}d\nu_{\alpha}(y)\\
                    &\quad+\int_{N_{x}^{c}}\left|\sigma(y)-\sigma(x)\right|\left|\ip{k_{x}^{\alpha}}{k_{y}^{\alpha}}_{\alpha}\right|^{2}d\nu_{\alpha}(y)
                \end{align*}
                Finally, using the the fact that $\sigma$ is bounded and the second estimate establishes the result for sufficiently large $\alpha$. 
        \end{proof}
        Using Lemma $3.2$, we arrive at the following proposition:
        \begin{prop}
            For $\sigma\in L^{p}(X,\nu)$, $1\leq p<\infty$, we have $\Tilde{\sigma}^{\alpha}\to\sigma$ in $L^{p}$
        \end{prop}
        \begin{proof}
            First, recall that $C_{c}(X)\subset C_{b}(X)\subset L^{p}(X,\nu)$ is dense for $1\leq p<\infty$. We first prove this result on this dense subspace. By Lemma 3.2, we have that $\Tilde{\sigma}^{\alpha}\to\sigma$ pointwise. Hence, $\left|\Tilde{\sigma}^{\alpha}(x)-\sigma(x)\right|^{p}\to 0$ as $\alpha\to\infty$. Furthermore, note that by triangle inequality, Jensen's inequality, and others, we have 
            \begin{align*}
                \left|\Tilde{\sigma}^{\alpha}(x)-\sigma(x)\right|^{p}\leq&2^{p-1}\left(\left|\Tilde{\sigma}^{\alpha}(x)\right|^{p}+\left|\sigma(x)\right|^{p}\right)\\
                =&2^{p-1}\left(\left|\int_{X}\sigma(y)\left|\ip{k_{x}^{\alpha}}{k_{y}^{\alpha}}_{\alpha}\right|^{2}d\nu_{\alpha}(y)\right|^{p}+\left|\sigma(x)\right|^{p}\right)\\
                \leq&2^{p-1}\left(\int_{X}\left|\sigma(y)\right|^{p}\left|\ip{k_{x}^{\alpha}}{k_{y}^{\alpha}}_{\alpha}\right|^{2}d\nu_{\alpha}(y)+\left|\sigma(x)\right|^{p}\right)\\
                =&2^{p-1}\left(\widetilde{\left|\sigma\right|^{p}}^{\alpha}(x)+\left|\sigma(x)\right|^{p}\right).
            \end{align*}
            Note that $\left|\sigma\right|^{p}\in C_{c}(X)$. Thus, letting 
            \begin{align*}
                g_{\alpha}(x)=2^{p-1}\left(\widetilde{\left|\sigma\right|^{p}}^{\alpha}(x)+\left|\sigma(x)\right|^{p}\right).
            \end{align*}
            we have a family of dominating functions such that $g_{\alpha}\to g:=2^{p}\left(\left|\sigma\right|^{p}\right)$ pointwise as $\alpha\to\infty$. Furthermore, note that for any $\alpha>0$, we have 
            \begin{align*}
                \int_{X}\widetilde{\left|\sigma\right|^{p}}^{\alpha}(x)d\nu(x)=\int_{X}\int_{X}\left|\sigma(y)\right|^{p}\left|\ip{k_{x}^{\alpha}}{k_{y}^{\alpha}}_{\alpha}\right|^{2}d\nu_{\alpha}(y)d\nu(x)=\int_{X}\left|\sigma(y)\right|^{p}d\nu(y).
            \end{align*}
            Hence, 
            \begin{align*}
                \int_{X}g_{\alpha}(x)d\nu(x)=2^{p}\int_{X}\left|\sigma(x)\right|^{p}d\nu(x)=\int_{X}g(x)d\nu(x).
            \end{align*}
            Thus, by the Dominated Convergence Theorem, we have that 
            $\norm{\Tilde{\sigma}^{\alpha}-\sigma}_{p}\to0$ as $\alpha\to\infty$. Now, let $\varepsilon>0$ and $\sigma\in L^{p}(X,\nu)$. Then there is a $\psi\in C_{c}(X)$ such that $\norm{\psi-\sigma}_{p}<\varepsilon$. By Lemma 4.1, there is a constant $A_{p}$ depending only on $p$, such that $\norm{\Tilde{f}^{\alpha}}_{p}\leq A_{p}\norm{f}_{p}$ for $f\in L^{p}(X,\nu)$. Hence, letting $f=\psi-\sigma$, we obtain $\norm{\Tilde{\psi}^{\alpha}-\Tilde{\sigma}^{\alpha}}_{p}\leq A_{p}\norm{\psi-\sigma}_{p}<A_{p}\varepsilon$. Finally, by the first part of this proof, there is an $\alpha'>0$ such that for $\alpha>\alpha'$, we have $\norm{\Tilde{\psi}^{\alpha}-\psi}_{p}<\varepsilon$. Combining these three estimates, we obtain
            \begin{align*}
                \norm{\Tilde{\sigma}^{\alpha}-\sigma}_{p}\leq\norm{\Tilde{\psi}^{\alpha}-\Tilde{\sigma}^{\alpha}}_{p}+\norm{\Tilde{\psi}^{\alpha}-\psi}_{p}+\norm{\psi-\sigma}_{p}<(2+A_{p})\varepsilon,
            \end{align*}
            and the convergence is established. 
        \end{proof}

\end{comment}        
%%%%%%%%%%%%%%%%%%%%%%%%%%%%%%%%%%%%%%%%%%%
\section{Convergence of the Berezin transform}

The goal of this section is to show that an approximate identity
condition implies convergence in $L^p$ (hence in measure)
of the Berezin transform, which is the key input for the
Szeg\H{o} limit theorems of Sections~4 and~5.

\medskip

\noindent
\begin{itemize}
    \item \textbf{Assumption (II)} For every $x\in X$ and every neighborhood $N_x$ of $x$,
\[
    \int_{N_x^c}
    |\langle k_x^\alpha,k_y^\alpha\rangle|^2
    \, d\nu_\alpha(y)
    \;\longrightarrow\; 0
    \qquad (\alpha\to\infty).
\]
\end{itemize}





\begin{lem}
For $1 \le p \le \infty$, the Berezin transform
$\sigma \mapsto \widetilde{\sigma}^\alpha$
defines a linear contraction on $L^p(X,\nu)$,
uniformly in $\alpha$.
\end{lem}

\begin{proof}
The $L^\infty$ bound follows immediately from
\[
    |\widetilde{\sigma}^\alpha(x)|
    \le
    \|\sigma\|_\infty
    \int_X
    |\langle k_x^\alpha,k_y^\alpha\rangle|^2
    d\nu_\alpha(y)
    =
    \|\sigma\|_\infty.
\]

For $p=1$, Fubini's theorem and Assumption (I) give
\[
    \|\widetilde{\sigma}^\alpha\|_1
    \le
    \int_X |\sigma(y)|\, d\nu(y)
    =
    \|\sigma\|_1.
\]

The general case follows by interpolation.
\end{proof}



\begin{lem}
If $\sigma \in C_c(X)$, then
$\widetilde{\sigma}^\alpha(x) \to \sigma(x)$
for every $x\in X$.
\end{lem}

\begin{proof}
Fix $x\in X$ and $\varepsilon>0$.
By continuity, there exists a neighborhood $N_x$
such that $|\sigma(y)-\sigma(x)|<\varepsilon$
for $y\in N_x$.

Write
\[
    \widetilde{\sigma}^\alpha(x)-\sigma(x)
    =
    \int_X
    (\sigma(y)-\sigma(x))
    |\langle k_x^\alpha,k_y^\alpha\rangle|^2
    d\nu_\alpha(y).
\]
Splitting the integral over $N_x$ and $N_x^c$,
the contribution on $N_x$ is bounded by $\varepsilon$,
while the contribution on $N_x^c$
tends to $0$ by Assumption~(II)
since $\sigma$ is bounded.
\end{proof}


\begin{prop}\label{Lp}
Let $1 \le p < \infty$ and $\sigma \in L^p(X,\nu)$.
Then
\[
    \widetilde{\sigma}^\alpha \to \sigma
    \quad \text{in } L^p(X,\nu).
\]
In particular, $\widetilde{\sigma}^\alpha \to \sigma$
in measure.
\end{prop}

\begin{proof}
First assume $\sigma \in C_c(X)$.
By the previous lemma,
$\widetilde{\sigma}^\alpha(x) \to \sigma(x)$ pointwise.
Moreover, Jensen’s inequality gives
\[
    |\widetilde{\sigma}^\alpha(x)|^p
    \le
    \widetilde{|\sigma|^p}^{\,\alpha}(x).
\]
Hence
\[
    |\widetilde{\sigma}^\alpha(x)-\sigma(x)|^p
    \le
    2^{p-1}
    \big(
        \widetilde{|\sigma|^p}^{\,\alpha}(x)
        + |\sigma(x)|^p
    \big).
\]
The right-hand side has constant integral,
since
\[
    \int_X
    \widetilde{|\sigma|^p}^{\,\alpha}(x)
    d\nu(x)
    =
    \int_X |\sigma(x)|^p d\nu(x).
\]
Since $|\sigma|^p\in C_c(X)$, Lemma~3.2 also gives
\[
\widetilde{|\sigma|^p}^{\,\alpha}(x)\longrightarrow |\sigma(x)|^p
\qquad\text{for every }x\in X.
\]
Thus the generalized dominated convergence theorem applies to
\[
f_\alpha:=\bigl|\widetilde{\sigma}^{\alpha}-\sigma\bigr|^p
\quad\text{and}\quad
g_\alpha:=2^{p-1}\left(
\widetilde{|\sigma|^p}^{\,\alpha}+|\sigma|^p
\right).
\]
Indeed, $0\leq f_\alpha\leq g_\alpha$, $f_\alpha\to0$ pointwise,
$g_\alpha\to2^p|\sigma|^p$ pointwise, and the preceding identity gives
\[
\int_X g_\alpha\,d\nu
=
\int_X 2^p|\sigma|^p\,d\nu.
\]
Consequently,
\[
\bigl\|\widetilde{\sigma}^{\alpha}-\sigma\bigr\|_p\longrightarrow0.
\]

For general $\sigma \in L^p$,
approximate by $\psi \in C_c(X)$ and use
the uniform $L^p$ boundedness of the Berezin transform.
\end{proof}

%%%%%%%%%%%%%%%%%%%%%%%%%%%%%%%%%%%%%%%%%%%%%%%%
\begin{comment}        
The above results hold true for Toeplitz operators whose symbols correspond to integrable functions. However, these can be viewed as Toeplitz operators whose symbol is an absolutely continuous measure on the phase space. However, one can also define Toeplitz operators when the symbol is a measure. 
\begin{definition}
    Let $\mu:\mathcal{B}(X)\to\C$ be a complex-valued Borel measure. We call the operator weakly defined as 
    \begin{align*}
        T_{\mu}f=\int_{X}\ip{f}{k_{x}}k_{x}d\mu(x)
    \end{align*}
    is the Toeplitz operator with symbol $\mu$. 
\end{definition}
While the results to follow concerning Szeg\H{o} limits do not abstract well to this situation, the associated Berezin transforms are still well-defined and as we will see converge in the limit $\alpha\to\infty$. First, we make the following observations. \\
The convergence of the Berezin transform above essentially is due to the function $c(\alpha)|\ip{k_{x}^{\alpha}}{k_{y}^{\alpha}}_{\alpha}|^{2}$ acting as a $\nu$-approximate identity. However, the separability of the measure $\nu_{\alpha}$ no longer makes sense in this circumstance. Instead, we replace Assumptions (I) and (II) with the following assumption, which still capture this essence:
\begin{itemize}
    \item \textbf{Assumption (III)} For each $y\neq x$, $|\ip{k_{x}^{\alpha}}{k_{y}^{\alpha}}_{\alpha}|^{2}\to 0$ as $\alpha\to\infty$. 
\end{itemize}
From these assumptions, we obtain the following. 
\begin{prop}
    Let $\mu:\mathcal{B}(X)\to\C$ be a complex-valued Borel measure such that $|\mu|$ is a finite measure, and let $T_{\mu}:\mathcal{H}_{\alpha}\to\mathcal{H}_{\alpha}$ denote the corresponding Toeplitz operator. Then, 
    \begin{align*}
        \Tilde{\mu}^{\alpha}(x):=\ip{T_{\mu}k_{x}^{\alpha}}{k_{x}^{\alpha}}_{\alpha}\to\mu(\{x\}),
    \end{align*}
    as $\alpha\to\infty$, for any $x\in X$. 
\end{prop}
\begin{proof}
    Let $x\in X$. Define the measure $\hat{\mu}:=|\mu-\mu(\{x\})\delta_{x}|$. Then $\hat{\mu}$ is a finite measure. Proceeding along, let $U\subset X$ be an open neighborhood of $x$. Then, 
    \begin{align*}
        \left|\Tilde{\mu}^{\alpha}(x)-\mu(\{x\})\right|=&\left|\int_{X}\left|\ip{k_{x}^{\alpha}}{k_{y}^{\alpha}}_{\alpha}\right|^{2}d\mu(y)-\mu(\{x\})\int_{X}\left|\ip{k_{x}^{\alpha}}{k_{y}^{\alpha}}_{\alpha}\right|^{2}d\delta_{x}(y)\right|\\
        \leq&\int_{U}\left|\ip{k_{x}^{\alpha}}{k_{y}^{\alpha}}_{\alpha}\right|^{2}d\hat{\mu}(y)+\int_{U^{c}}\left|\ip{k_{x}^{\alpha}}{k_{y}^{\alpha}}_{\alpha}\right|^{2}d\hat{\mu}(y).
    \end{align*}
    Let $\varepsilon>0$. Now, as $\left|\ip{k_{x}^{\alpha}}{k_{y}^{\alpha}}_{\alpha}\right|^{2}\leq 1$, and $\hat{\mu}$ is a finite measure with $\hat{\mu}(\{x\})=0$, we may choose $U$ sufficiently small such that $\hat{\mu}(U)<\varepsilon/2$. Thus, 
    \begin{align*}
        \int_{U}\left|\ip{k_{x}^{\alpha}}{k_{y}^{\alpha}}_{\alpha}\right|^{2}d\hat{\mu}(y)\leq\hat{\mu}(U)<\varepsilon/2.
    \end{align*}
    Likewise, $\left|\ip{k_{x}^{\alpha}}{k_{y}^{\alpha}}_{\alpha}\right|^{2}\leq 1$ on $U^{c}$, and since $\hat{\mu}(\{x\})=0$ and Assumption (III), we have that $\left|\ip{k_{x}^{\alpha}}{k_{y}^{\alpha}}_{\alpha}\right|^{2}\to 0$ $\hat{\mu}$-a.e. as $\alpha\to\infty$. Thus, by the Bounded Convergence Theorem, there is an $\alpha$ sufficiently large such that 
    \begin{align*}
        \int_{U^{c}}\left|\ip{k_{x}^{\alpha}}{k_{y}^{\alpha}}_{\alpha}\right|^{2}d\hat{\mu}(y)<\varepsilon/2.
    \end{align*}
    Combining the results gives the desired limit. 
\end{proof}

\end{comment}
%%%%%%%%%%%%%%%%%%%%%%%%%%%%%%%%
%%%%%%%%%%%%%%%%%%%%%%%%%%%%%%%%%
The preceding results apply to Toeplitz operators with integrable symbols, which may be viewed as Toeplitz operators whose symbol is an absolutely continuous measure on phase space. More generally, one may define Toeplitz operators whose symbol is a (finite) complex Borel measure.

\begin{definition}
Let $\mu:\mathcal{B}(X)\to\C$ be a complex-valued Borel measure. The (weakly defined) Toeplitz operator with symbol $\mu$ is given by
\[
T_\mu f
= \int_X \langle f,k_x^\alpha\rangle_\alpha \, k_x^\alpha \, d\mu(x).
\]
\end{definition}

Although the Szeg\H{o}-type limit theorems do not extend in full generality to this setting, the associated Berezin transforms remain well defined. However, the separability of the measure $\nu_{\alpha}$ no longer makes sense in this circumstance. Instead, we replace Assumptions (I) and (II) with the following assumption, which still capture this essence:
\begin{itemize}
    \item \textbf{Assumption (III)} For each $y\neq x$, $|\ip{k_{x}^{\alpha}}{k_{y}^{\alpha}}_{\alpha}|^{2}\to 0$ as $\alpha\to\infty$. 
\end{itemize}
From this assumption, we obtain the following. 

\begin{prop}
Let $\mu$ be a complex Borel measure with $|\mu|(X)<\infty$, and let $T_\mu:\mathcal{H}_\alpha\to\mathcal{H}_\alpha$ denote the corresponding Toeplitz operators. Then, for every $x\in X$,
\[
\widetilde{\mu}^\alpha(x)
:= \langle T_\mu k_x^\alpha, k_x^\alpha\rangle_\alpha
\longrightarrow \mu(\{x\}),
\qquad \text{as } \alpha\to\infty.
\]
\end{prop}

\begin{proof}
% Fix $x\in X$. Since
% \[
% \widetilde{\mu}^\alpha(x)
% = \int_X |\langle k_x^\alpha,k_y^\alpha\rangle_\alpha|^2 \, d\mu(y),
% \]
% write $\mu=\mu(\{x\})\delta_x + (\mu-\mu(\{x\})\delta_x)$ and set
% \[
% \hat\mu := |\mu-\mu(\{x\})\delta_x|.
% \]
% Then $\hat\mu$ is a finite measure with $\hat\mu(\{x\})=0$, and
% \[
% \bigl|\widetilde{\mu}^\alpha(x)-\mu(\{x\})\bigr|
% \le
% \int_X |\langle k_x^\alpha,k_y^\alpha\rangle_\alpha|^2 \, d\hat\mu(y).
% \]

% Let $\varepsilon>0$. Choose an open neighbourhood $U$ of $x$ such that $\hat\mu(U)<\varepsilon/2$. Since
% $|\langle k_x^\alpha,k_y^\alpha\rangle_\alpha|^2 \le 1$, we obtain
% \[
% \int_U |\langle k_x^\alpha,k_y^\alpha\rangle_\alpha|^2 \, d\hat\mu
% \le \hat\mu(U)
% < \varepsilon/2.
% \]

% On $U^c$, Assumption (II) implies that
% $|\langle k_x^\alpha,k_y^\alpha\rangle_\alpha|^2 \to 0$
% for $y\neq x$; hence the convergence holds $\hat\mu$-a.e. By bounded convergence,
% \[
% \int_{U^c}
% |\langle k_x^\alpha,k_y^\alpha\rangle_\alpha|^2
% \, d\hat\mu
% \longrightarrow 0.
% \]
% For $\alpha$ sufficiently large this term is $<\varepsilon/2$, which yields the claim.
Let $x\in X$. Define the measure $\hat{\mu}:=|\mu-\mu(\{x\})\delta_{x}|$. Then $\hat{\mu}$ is a finite measure. Proceeding along, let $U\subset X$ be an open neighborhood of $x$. Then, 
    \begin{align*}
        \left|\Tilde{\mu}^{\alpha}(x)-\mu(\{x\})\right|=&\left|\int_{X}\left|\ip{k_{x}^{\alpha}}{k_{y}^{\alpha}}_{\alpha}\right|^{2}d\mu(y)-\mu(\{x\})\int_{X}\left|\ip{k_{x}^{\alpha}}{k_{y}^{\alpha}}_{\alpha}\right|^{2}d\delta_{x}(y)\right|\\
        \leq&\int_{U}\left|\ip{k_{x}^{\alpha}}{k_{y}^{\alpha}}_{\alpha}\right|^{2}d\hat{\mu}(y)+\int_{U^{c}}\left|\ip{k_{x}^{\alpha}}{k_{y}^{\alpha}}_{\alpha}\right|^{2}d\hat{\mu}(y).
    \end{align*}
    Let $\varepsilon>0$. Now, as $\left|\ip{k_{x}^{\alpha}}{k_{y}^{\alpha}}_{\alpha}\right|^{2}\leq 1$, and $\hat{\mu}$ is a finite measure with $\hat{\mu}(\{x\})=0$, we may choose $U$ sufficiently small such that $\hat{\mu}(U)<\varepsilon/2$. Thus, 
    \begin{align*}
        \int_{U}\left|\ip{k_{x}^{\alpha}}{k_{y}^{\alpha}}_{\alpha}\right|^{2}d\hat{\mu}(y)\leq\hat{\mu}(U)<\varepsilon/2.
    \end{align*}
    Likewise, $\left|\ip{k_{x}^{\alpha}}{k_{y}^{\alpha}}_{\alpha}\right|^{2}\leq 1$ on $U^{c}$, and since $\hat{\mu}(\{x\})=0$ and Assumption (III), we have that $\left|\ip{k_{x}^{\alpha}}{k_{y}^{\alpha}}_{\alpha}\right|^{2}\to 0$ $\hat{\mu}$-a.e. as $\alpha\to\infty$. Thus, by the Bounded Convergence Theorem, there is an $\alpha$ sufficiently large such that 
    \begin{align*}
        \int_{U^{c}}\left|\ip{k_{x}^{\alpha}}{k_{y}^{\alpha}}_{\alpha}\right|^{2}d\hat{\mu}(y)<\varepsilon/2.
    \end{align*}
    Combining the results gives the desired limit. 
\end{proof}

\begin{prop}\label{lkc}
Let $\mu$ be a finite complex Radon measure on $X$. Under
Assumption~(III), for every $x\in X$,
\[
\int_X
\langle k_x^\alpha,k_y^\alpha\rangle_\alpha\,d\mu(y)
\longrightarrow
\mu(\{x\})
\qquad
(\alpha\to\infty).
\]
\end{prop}

\begin{proof}
Set
\[
A_\alpha(x)
=
\int_X
\langle k_x^\alpha,k_y^\alpha\rangle_\alpha\,d\mu(y)
\]
and
\[
\widehat\mu
=
\left|\mu-\mu(\{x\})\delta_x\right|.
\]
Since $\langle k_x^\alpha,k_x^\alpha\rangle_\alpha=1$, we have
\[
A_\alpha(x)-\mu(\{x\})
=
\int_X
\langle k_x^\alpha,k_y^\alpha\rangle_\alpha
\,d\bigl(\mu-\mu(\{x\})\delta_x\bigr)(y).
\]

Let $\varepsilon>0$. Choose a neighborhood $U$ of $x$ such that
\[
\widehat\mu(U)<\frac{\varepsilon}{2}.
\]
Since the kernels are normalized,
\[
\left|\langle k_x^\alpha,k_y^\alpha\rangle_\alpha\right|
\leq 1.
\]
Therefore,
\[
\begin{aligned}
|A_\alpha(x)-\mu(\{x\})|
&\leq
\int_U
\left|\langle k_x^\alpha,k_y^\alpha\rangle_\alpha\right|
\,d\widehat\mu(y) \\
&\quad+
\int_{U^c}
\left|\langle k_x^\alpha,k_y^\alpha\rangle_\alpha\right|
\,d\widehat\mu(y).
\end{aligned}
\]
The first term is smaller than $\varepsilon/2$. For $y\neq x$,
Assumption~(III) implies
\[
\left|\langle k_x^\alpha,k_y^\alpha\rangle_\alpha\right|
\longrightarrow 0.
\]
Hence the second term tends to zero by dominated convergence. Thus
\[
\limsup_{\alpha\to\infty}
|A_\alpha(x)-\mu(\{x\})|
\leq \frac{\varepsilon}{2}.
\]
Since $\varepsilon$ is arbitrary, the result follows.
\end{proof}



% For the above proposition, the full strength of Assumptions (I) and (II) is not required. 

% The argument uses only the off--diagonal decay
% \[
% |\langle k_x^\alpha,k_y^\alpha\rangle_\alpha|^2 \longrightarrow 0
% \qquad (y\neq x),
% \]
% that is, the fact that the reproducing kernels concentrate at the base point as $\alpha\to\infty$. Thus the conclusion remains valid under the weaker assumption:

% \begin{itemize}
% \item[(III)] For each $y\neq x$, 
% \[
% |\langle k_x^\alpha,k_y^\alpha\rangle_\alpha|^2 \to 0
% \quad \text{as } \alpha\to\infty.
% \]
% \end{itemize}

% In the abstract framework considered here, Assumption (III) is already a consequence of (II). However, isolating it clarifies that only pointwise off--diagonal decay is needed for the atomic limit of the Berezin transform, whereas the stronger structural hypotheses (I)--(II) are essential for global Szeg\H{o}-type results.

%%%%%%%%%%%%%%%%%%%%%%%%%%%%%%%%%%%%%%%%%%%%%
\begin{comment}


\section{Szeg\H{o} Limit Theorems for Finite Measures}
        Using the reasoning in Section 2, we can obtain the following Szeg\H{o} limit theorems in the compact, finite measure setting. First, we require a more general result.
        % \begin{cor}
        %     Suppose $\sigma\in L^{p}(X,\nu)$ for some $1\leq p<\infty$ and $\sigma\geq 0$. Let $\psi:\R\to\R^{+}$ be uniformly continuous and convex. Then, 
        %     \begin{align*}
        %         \frac{\text{tr}\left(\psi \left(T_{\sigma}^{\alpha}\right)\right)}{c(\alpha)}\to\int_{X}\psi(\sigma(x))d\nu(x)
        %     \end{align*}
        %     as $\alpha\to\infty$.
        % \end{cor}
    %     \begin{remark}
    %         In Corollary $3.4$, the assumption that $T_{\sigma}^{\alpha}$ is trace class is not needed. In fact, the Berezin-Lieb inequality allows for the circumstance when $\tr(T_{\sigma}^{\alpha})=\infty$. Thus, in this situation, the above limit is taken to be infinity. However, under the additional assumption that $\sigma\in L^{\infty}(X,\nu)$, $T_{\sigma}^{\alpha}$ is trace class, and we obtain the following stronger result.  
    %     \end{remark}
    %     \begin{cor}
    %         Suppose $\sigma\in L^{p}(X,\nu)\cap L^{\infty}(X,\nu)$ for some $1\leq p<\infty$ and $\sigma\geq 0$. Let $\varphi:[0,\norm{\sigma}_{\infty}]\to\R$ be continuous with $\psi(0)=0$. Then, 
    %         \begin{align*}
    %             \frac{\text{tr}\left(\varphi \left(T_{\sigma}^{\alpha}\right)\right)}{c(\alpha)}\to\int_{X}\varphi(\sigma(x))d\nu(x)
    %         \end{align*}
    %         as $\alpha\to\infty$.
    %     \end{cor}
    %     \begin{proof}
    %         First, without loss of generality, we suppose that $\norm{\sigma}_{\infty}=1$. By our main theorem, we have that for each $n\in\N$ that the limit holds in the case when $\psi(x)=x^{n}$ as $\alpha\to\infty$, since $x\mapsto x^{n}$ is uniformly continuous and convex on $[0,1]$. The result then follows by taking linear combinations to form the polynomials, and a standard Stone-weierstrass argument. 
    %     \end{proof}
    % In fact, one can obtain a slightly stronger result in the bounded case of Corollary $3.6$. We require the following lemma first, which can be viewed as a generalization of the Berezin-Lieb inequality. 
% Finally, from this result, we obtain a corollary which is a generalization of \cite{FN01}. 
% \begin{cor}
%     Suppose $\psi:\mathbb{R}^{+}\to\mathbb{R}^{+}$ is continuous. Let $\sigma\in L^{1}(X,\nu)\cap L^{\infty}(X,\nu)$ be real-valued. Then, 
%     \begin{align*}
%         \lim_{\alpha\to\infty}\frac{1}{c(\alpha)}\text{tr}\left(T_{\sigma}^{\alpha}\psi\left(T_{\sigma}^{\alpha}\right)\right)=\int_{X}\sigma(x)\psi(\sigma(x))d\nu(x).
%     \end{align*}
% \end{cor}
% Now that we have obtained the Szeg\H{o} limits, we use these results to obtain corresponding Weyl laws.  
\begin{thm}
    Let $\psi:\mathbb{R}\to\mathbb{R}$ be uniformly continuous on the range of $\sigma$, convex, and non-negative. Suppose $\sigma:X\to\mathbb{R}$ is such that $T_{\sigma}^{\alpha}$ is diagonalizable and  $\Tilde{\sigma}^{\alpha}\to\sigma$ in $\nu$-measure as $\alpha\to\infty$. Then, 
    \begin{align*}
        \lim_{\alpha\to\infty}\frac{1}{c(\alpha)}\text{tr}\left(\psi\left(T_{\sigma}^{\alpha}\right)\right)=\int_{X}\psi\left(\sigma(x)\right)d\nu(x).
    \end{align*}
    %If $\psi$ is continuous, convex, and non-negative, the same result holds for $\sigma\in C_{c}(X)$.
\end{thm}
\begin{proof}
            %We will prove the result when $\Tilde{\sigma}^{\alpha}\to\sigma$ in $\nu$-measure. The case when $\sigma\in C_{c}(X)$ (i.e., when the Berezin transform converges pointwise) is similar. 
             Since $\Tilde{\sigma}_{\alpha}\to\sigma$ in $\nu$-measure, using that $\psi$ is uniformly continuous by continuity mapping \cite{Bart61}, we have $\psi\left(\Tilde{\sigma}_{\alpha}\right)\to\psi\left(\sigma \right)$ in $\nu$-measure. We now apply Proposition 2.1 for each $\alpha>0$, taking $\eta=1$:
                \begin{align*}
                    \int_{X}\psi\left(\Tilde{\sigma}_{\alpha}(x)\right)d\nu_{\alpha}(x)\leq \text{tr}\left(\psi\left(T_{\sigma}^{\alpha}\right)\right)\leq\int_{X}\psi\left(\sigma(x)\right)d\nu_{\alpha}(x).
                \end{align*}
                By Assumption (I), we have for any $\alpha>0$
                \begin{align*}
                    \int_{X}\psi\left(\Tilde{\sigma}_{\alpha}(x)\right)d\nu(x)\leq \frac{1}{c(\alpha)}\text{tr}\left(\psi\left(T_{\sigma}^{\alpha}\right)\right)\leq\int_{X}\psi\left(\sigma(x)\right)d\nu(x).
                \end{align*}
                Applying Fatou's lemma in the lower bound, we see that 
                \begin{align*}
                    \int_{X}\psi\left(\sigma(x)\right)d\nu(x)\leq\liminf_{\alpha\to\infty}\int_{X}\psi\left(\Tilde{\sigma}_{\alpha}(x)\right)d\nu(x)\leq\liminf_{\alpha\to\infty}\frac{1}{c(\alpha)}\text{tr}\left(\psi\left(T_{\sigma}^{\alpha}\right)\right).
                \end{align*}
                Now, using the upper inequality, we obtain
                \begin{align*}
                    \limsup_{\alpha\to\infty}\frac{1}{c(\alpha)}\text{tr}\left(\psi\left(T_{\sigma}^{\alpha}\right)\right)\leq\int_{X}\psi\left(\sigma(x)\right)d\nu(x).
                \end{align*}
                Combining these two inequalities, we obtain the desired limit. 
        \end{proof}
        %  We now obtain the following analog to Theorem 1.1. 
        % \begin{thm}\label{Sze_log}
        %     Suppose $\sigma\in L^{p}(X)$ for some $p\in [1,\infty)$ is strictly positive such that $\Tilde{\sigma}^{\alpha}\to\sigma$ in $\nu$-measure as $\alpha\to\infty$. Then, 
        %     \begin{align*}
        %         \lim_{\alpha\to\infty}\frac{1}{c(\alpha)}\text{tr}\left(\log\left(T_{\sigma}^{\alpha}\right)\right)=\int_{X}\log\left(\sigma(x)\right)d\nu(x).
        %     \end{align*}
        % \end{thm}
        % \begin{proof}
        %     The proof of this result is similar in nature to Theorem 3.1, but requires more care due to the logarithm not being uniformly continuous on $(0, \infty)$ and the concavity of the logarithm. For this, let $c>0$. 
        %     Note that $\log(x+c)$ is uniformly continuous and concave, and $\log(x+c)\leq x$ giving $\log(\sigma(x)+c)\leq\frac{1}{p}\sigma(x)^{p}\in L^{1}(X)$. Hence, using the Berezin-Lieb inequality for concave functions and the Fatou-Lebesgue Theorem, we have 
        %     \begin{align*}
        %         \limsup_{\alpha\to\infty}\frac{1}{c(\alpha)}\text{tr}\left(\log\left(T_{\sigma}^{\alpha}\right)\right)\leq&\limsup_{\alpha\to\infty}\frac{1}{c(\alpha)}\text{tr}\left(\log\left(T_{\sigma}^{\alpha}+cI\right)\right)\\
        %         \leq& \limsup_{\alpha\to\infty}\int_{X}\log(\Tilde{\sigma}^{\alpha}(x)+c)d\nu(x)\\
        %         \leq&\int_{X}\log(\sigma(x)+c)d\nu(x),
        %     \end{align*}
        %     using the convergence and uniform continuity of $\log(x+c)$. Since the left hand side is independent of $c$, taking $c\to 0$ and applying the Fatou-Lebesgue Theorem again gives the upper bound. For the lower bound, notice that by the Berezin-Lieb inequality for the concave function $\log(x)$ we have 
        %     \begin{align*}
        %         \int_{X}\log(\sigma(x))d\nu(x)\leq\liminf_{\alpha\to\infty}\frac{1}{c(\alpha)}\text{tr}\left(\log\left(T_{\sigma}^{\alpha}\right)\right).
        %     \end{align*}
        %     This proves the result. 
        % \end{proof}
       
We now specialize to the finite dimensional case. 

\begin{cor}\label{finite}
    Suppose $\text{dim}(\mathcal{H}_{\alpha})<\infty$, $\text{dim}(\mathcal{H}_{\alpha})/c(\alpha)$ is bounded in $\alpha$, and $\nu$ is a finite measure on $X$. Let $\sigma\in L^{\infty}(X,\nu)$ be real-valued such that $\Tilde{\sigma}^{\alpha}\to\sigma$ in $\nu$-measure as $\alpha\to\infty$, and let $\psi: \R \to\R$ be continuous on the range of $\sigma$. Then,
    \begin{align*}
        \lim_{\alpha\to\infty}\frac{1}{c(\alpha)}\text{tr}\left(\psi\left(T_{\sigma}^{\alpha}\right)\right)=\int_{X}\psi(\sigma(x))d\nu(x).
    \end{align*}
\end{cor}
Before proving this result, we state the following trace estimates which are standard and we include without proof. 
\begin{lem}
    Let $\mathcal{H}$ be a Hilbert space, let $A:\mathcal{H}\to\mathcal{H}$ be a self-adjoint operator, and let $B:\mathcal{H}\to\mathcal{H}$ be a positive operator. Let $\psi:\text{spec}(A)\to\R$ be continuous. Then
    \begin{itemize}
        \item[1.] If $\text{dim}(\mathcal{H})<\infty$,
        \begin{align*}
            \left|\text{tr}\left(\psi\left(A\right)\right)\right|\leq\text{dim}(\mathcal{H})\sup\{|\psi(x)|:x\in\text{spec}(A)\}.
        \end{align*}
        \item[2.] If $A$ is trace-class, then 
            \begin{align*}
                \left|\text{tr}\left(A\psi(A)\right)\right|\leq\norm{A}_{1}\sup\{|\psi(x)|:x\in\text{spec}(A)\},
            \end{align*}
            where $\norm{\cdot}_{1}$ is the trace-norm of $A$.
            \item[3.] If $B$ is trace-class and $A$ and $B$ are simultaneously diagonalizable by a common eigenbasis, then 
            \begin{align*}
                |\text{tr}(B\psi(A))|\leq \norm{B}_{1}\sup\{|\psi(x)|:x\in\text{spec}(A)\}.
            \end{align*}
    \end{itemize}
\end{lem}
\begin{proof}[Proof of Corollary~\ref{finite}]
    We will first prove this statement for the monomials $\psi(x)=x^{k}$. Since $\psi(x)=x^{2k}$ is uniformly continuous and convex on $[-\norm{\sigma}_\infty, \norm{\sigma}_\infty]$, the result in this case follows directly from Theorem 4.1. The case when $\psi(x)=x^{2k+1}$ requires more care. Let 
    \begin{align*}
        x^{+}=\begin{cases}
            x & x>0\\
            0 & x\leq 0.
        \end{cases} 
    \end{align*}
    and similarly define $x^{-}$. Note that $\psi^{+}(x)=\left(x^{2k+1}\right)^{+}$ is convex and continuous on $\R$. Hence, $\psi^{+}$ is convex and uniformly continuous on $[\inf\sigma, \sup\sigma]$ and by the Theorem 3.2 we have 
    \begin{align*}
        \frac{1}{c(\alpha)}\text{tr}_{\alpha}\left(\psi^{+}\left(T_{\sigma}\right)\right)\to\int_{X}\left(\sigma(x)^+\right)^{2k+1} d\nu(x).
    \end{align*}
    
    Likewise, the function $\psi^{-}=\left(x^{2k+1}\right)^{-}$ is convex and continuous. Thus, in the even case, we have
    \begin{align*}
        \frac{1}{c(\alpha)}\text{tr}_{\alpha}\left(\left(T_{\sigma}\right)^{2k}\right)\to\int_{X}\left(\sigma(x)\right)^{2k}d\nu(x),
    \end{align*}
    while in the odd case we obtain 
    \begin{align*}
        \frac{1}{c(\alpha)}\text{tr}_{\alpha}\left(\psi_{2k+1}^{\pm}\left(T_{\sigma}\right)\right)\to\int_{X}\psi_{2k+1}^{\pm}\left(\sigma(x)\right)d\nu(x).
    \end{align*}
    Thus, using that the trace is additive we see that the result holds for $\psi$ a polynomial. The result then follows from the first part of Lemma 4.3 and an application of the Stone-Weierstrass Theorem. 
\end{proof}

We now obtain the following analogy to Theorem 1.1, which is motivated by the corresponding form of the first Szeg\H o theorem appearing in~\cite{simon2005orthogonal}. 
\begin{cor}
    Suppose $c(\alpha)=\text{dim}(\mathcal{H}_{\alpha})<\infty$ for each $\alpha>0$ and $\nu$ is a probability measure. Let $\sigma\in L^{1}(X,\nu)$ be non-negative such that $\Tilde{\sigma}^{\alpha}\to\sigma$ in $\nu$-measure as $\alpha\to\infty$. Let $\psi:\R\to\R$ be continuous on the range of $\sigma$ and such that $\lim_{x\to\infty}\frac{\psi(x)}{x}$ exists. Then, 
    \begin{align*}
        \lim_{\alpha\to\infty}\frac{1}{c(\alpha)}\text{tr}\left(\psi\left(T_{\sigma}^{\alpha}\right)\right)=\int_{X}\psi(\sigma(x))d\nu(x).
    \end{align*}
\end{cor}
\begin{proof}
    We first consider a special class of $\psi$. For $\psi(x)=x$, the result clearly follows since $T_{\sigma}^{\alpha}$ is a Toeplitz operator. We now prove the case when $\psi(x)=\log(x+c)$ for $c>0$. First, note that for $x\geq 0$, $|\log(x+c)|\leq x+|\log(c)|$. Hence, if $\sigma\in L^{1}(X,\nu)$, then $\log(\sigma+c)\in L^{1}(X,\nu)$ since $\nu(X)=1$. Thus, applying Proposition 2.1 when $\eta=1$ in the concave case 
    \begin{align*}
        \int_{X}\log(\sigma(x)+c)d\nu(x)\leq\frac{1}{c(\alpha)}\tr\left(\log\left(T_{\sigma}^{\alpha}+c\right)\right)\leq\int_{X}\log(\Tilde{\sigma}^{\alpha}(x)+c)d\nu(x). 
    \end{align*}
    Now, note that the lower bound gives 
    \begin{align*}
        \int_{X}\log(\sigma(x)+c)d\nu(x)\leq\liminf_{\alpha\to\infty}\frac{1}{c(\alpha)}\tr\left(\log\left(T_{\sigma}^{\alpha}+c\right)\right).
    \end{align*}
    For the upper bound, using the logarithmic inequality along with the Dominated Convergence Theorem we see that 
    \begin{align*}
        \limsup_{\alpha\to\infty}\frac{1}{c(\alpha)}\tr\left(\log\left(T_{\sigma}^{\alpha}+c\right)\right)\leq&\limsup_{\alpha\to\infty}\int_{X}\log(\Tilde{\sigma}^{\alpha}(x)+c)d\nu(x)\\
        =&\int_{X}\log(\sigma(x)+c)d\nu(x),
    \end{align*}
    by the same reasoning as in the proof of Theorem 4.1. Thus, the limit is established in this case and we have 
    \begin{align*}
        \lim_{\alpha\to\infty}\frac{1}{c(\alpha)}\text{tr}\left(\log\left(T_{\sigma}+c\right)\right)=\int_{X}\log\left(\sigma(x)+c\right)d\nu(x).
    \end{align*}
    Now, let
    \begin{align*}
        L_{\alpha}(\psi)=\frac{1}{c(\alpha)}\text{tr}\left(\psi\left(T_{\sigma}\right)\right)
    \end{align*}
    and
    \begin{align*}
        L(\psi)=\int_{X}\psi(\sigma(x))d\nu(x).
    \end{align*}
    Clearly, both $L_{\alpha}$ and $L$ are linear in $\psi$. Also $|L_{\alpha}(\psi)|\leq\norm{\psi}_{\infty}$ and $|L(\psi)|\leq\norm{\psi}_{\infty}$. Let $\psi^{c}(x)=\log(x+c)$. By the above, $L_{\alpha}(\psi^{c})\to L(\psi^{c})$ for any $c>0$ and hence by Lemma 3.4 $L_{\alpha}(\psi)\to L(\psi)$ for any $\psi$ in the uniform closure of $\{\psi^{c}\}$. But the derivatives of $\psi^{c}$ with respect to $c$ at $c>0$ are uniform limits, so this closure must contain $\{1/(x+1)^{n}\}_{n=0}^{\infty}$. Thus, the limit is established for $\psi\in C_{0}([0,\infty))$. \\
    Now, suppose $\psi(x)/x\to 0$ as $x\to\infty$. Then for $\varepsilon>0$ we may write $\psi(x)=\psi_{1,\varepsilon}(x)+\psi_{2,\varepsilon}$ with $\psi_{1,\varepsilon}\in C_{0}([0,\infty))$ and $|\psi_{2,\varepsilon}(x)|\leq\varepsilon x$. Then, 
    \begin{align*}
        |L_{\alpha}(\psi_{2,\varepsilon})|\leq\varepsilon L_{\alpha}(x)=\varepsilon\int_{X}\sigma(x)d\nu(x). 
    \end{align*}
    Thus, 
    \begin{align*}
        \limsup_{\alpha\to\infty}\left|L_{\alpha}(\psi)-L(\psi)\right|\leq2\varepsilon \int_{X}\sigma(x)d\nu(x).
    \end{align*}
    But $\varepsilon>0$ is arbitrary. Taking $\varepsilon\to 0$ proves the result in this case. Finally, suppose $\psi(x)/x\to \beta<\infty$ as $x\to\infty$. Taking $\psi_{1}(x)=\psi(x)-\beta x$ gives the result in general by linearity, the case when $\psi(x)/x\to 0$ and when $\psi(x)=x$. 
\end{proof}


\end{comment}
%%%%%%%%%%%%%%%%%%%%%%%%%%%%%%%%%%%%%%%%%%%%%%%%%%%%%%

\section{Abstract Szeg\H{o} Limit Theorem for Compact Spaces}

Throughout this section, let $\sigma \in L^1(X,\nu)$ be real-valued, and $\nu$ is a finite measure. We also assume that all the Hilbert spaces $H_\alpha$ are finite-dimensional. Under these assumptions the associated Toeplitz operators $T_\sigma^\alpha$
are trace-class, self-adjoint and satisfy
\[
\widetilde{\sigma}^\alpha \longrightarrow \sigma
\quad\text{in $\nu$-measure as } \alpha\to\infty.
\]

We begin with the convex case, which forms the core of the argument.

\begin{thm}\label{thm:convex_szego}
\leavevmode
\begin{itemize}
\item[(i)] Suppose $\sigma \ge 0$ and $\psi \in C([0,\infty))$
is convex such that the limit 
\[
\lim_{x\to\infty} \frac{\psi(x)}{x}
\]
exists. Then
\[
\lim_{\alpha\to\infty}
\frac{1}{c(\alpha)}\tr\!\big(\psi(T_\sigma^\alpha)\big)
=
\int_X \psi(\sigma(x))\,d\nu(x).
\]

\item[(ii)] If $\sigma \in L^\infty(X,\nu)$ and
$\psi$ is convex and continuous on
$[\text{ess inf } \sigma, \text{ess sup }\sigma]$, then the same limit holds.
\end{itemize}
\end{thm}

\begin{proof}
We prove (i); part (ii) follows by a simpler argument.


For $\psi(x)=x$ the result holds since
\[
\frac{1}{c(\alpha)}\tr(T_\sigma^\alpha)
=
\int_X \sigma\, d\nu,
\]
which is independent of $\alpha$. 

Write
\[
\psi(x)=\beta x + \psi_0(x),
\qquad
\text{where } \frac{\psi_0(x)}{x}\to 0.
\]
By linearity it suffices to treat the case
$\psi(x)/x\to 0$.

For convex $\psi$, Proposition~\ref{BL1} gives
\[
\int_X \psi(\widetilde{\sigma}^\alpha)\,d\nu_\alpha
\le
\tr(\psi(T_\sigma^\alpha))
\le
\int_X \psi(\sigma)\,d\nu_\alpha.
\]
Dividing by $c(\alpha)$ and using Assumption~(I),
\[
\int_X \psi(\widetilde{\sigma}^\alpha)\,d\nu
\le
\frac{1}{c(\alpha)}\tr(\psi(T_\sigma^\alpha))
\le
\int_X \psi(\sigma)\,d\nu.
\]

Since $\psi(x)/x\to 0$, for every $\varepsilon>0$
there exists $C_\varepsilon$ such that
\[
|\psi(x)| \le \varepsilon x + C_\varepsilon.
\]
By Proposition~\ref{Lp}, $\widetilde{\sigma}^{\alpha}\to\sigma$ in
$L^1(X,\nu)$, and hence the family
$\{\widetilde{\sigma}^{\alpha}\}_{\alpha}$ is uniformly integrable.
Since $\nu(X)<\infty$, the preceding estimate implies that the family
$\{\psi(\widetilde{\sigma}^{\alpha})\}_{\alpha}$ is uniformly
integrable. Moreover,
\[
\psi(\widetilde{\sigma}^{\alpha})\longrightarrow\psi(\sigma)
\]
in measure. Therefore, Vitali's convergence theorem gives
\[
\int_X\psi(\widetilde{\sigma}^{\alpha})\,d\nu
\longrightarrow
\int_X\psi(\sigma)\,d\nu .
\]

\medskip
\noindent
\textbf{Proof of (ii).}
If $\sigma\in L^\infty$, its essential range
is compact. A convex continuous function on a compact
interval is uniformly continuous.
The same Berezin--Lieb argument as above,
now without growth considerations,
yields the limit directly.
\end{proof}

\subsection{From Convex Functions to General Continuous Symbols}

We extend the Szeg\H{o} theorem from convex functions to general continuous functions via a decomposition into differences of convex functions.

\begin{lem}\label{lem:diff_convex}
Let $\psi$ be real-valued. 
\begin{itemize}
    \item[(i)] If $\psi\in C[a,b]$, then for every $\varepsilon>0$ there exist convex functions $\psi_1,\psi_2$ on $[a,b]$ such that 
    \[
        \|\psi-(\psi_1-\psi_2)\|_{\infty} < \varepsilon.
    \]
    \item[(ii)] If $\psi\in C_0[0,\infty)$, then for every $\varepsilon>0$ there exist convex functions $\psi_1,\psi_2\in C_0[0,\infty)$ such that
    \[
        \|\psi-(\psi_1-\psi_2)\|_\infty < \varepsilon.
    \]
\end{itemize}
\end{lem}

\begin{proof}
(i) For a polynomial $p$ on $[a,b]$, let $M = \inf_{x\in[a,b]} p''(x)$. If $M\geq 0$ then $p$ is convex itself, so we can trivially decompose $p=p-0$. If $M<0$ then 
\[
    p(x) = (p(x)-\frac{M}{2}x^2) - (-\frac{M}{2})x^2,
\]
where both terms are convex. Any $\psi\in C[a,b]$ can be uniformly approximated by polynomials; applying this decomposition yields the claim.

(ii) Let
\[
\mathcal C
:=
\{\varphi\in C_0([0,\infty)):
\varphi\text{ is convex}\}
\]
and
\[
\mathcal A
:=
\{\varphi_1-\varphi_2:
\varphi_1,\varphi_2\in\mathcal C\}.
\]

We first show that every $\varphi\in\mathcal C$ is non-negative and
non-increasing. If $0\leq s<t<u$, convexity gives
\[
\frac{\varphi(t)-\varphi(s)}{t-s}
\leq
\frac{\varphi(u)-\varphi(t)}{u-t}.
\]
Letting $u\to\infty$ and using $\varphi(u)\to 0$, we obtain
\[
\frac{\varphi(t)-\varphi(s)}{t-s}\leq 0.
\]
Thus $\varphi$ is non-increasing. Since $\varphi(t)\to 0$ as
$t\to\infty$, it follows that $\varphi\geq 0$.

We next show that $\mathcal C$ is closed under multiplication. Let
$f,g\in\mathcal C$. Then $fg\in C_0([0,\infty))$. Suppose that
$x\leq y$ and let $z=\lambda x+(1-\lambda)y$, where
$0\leq\lambda\leq 1$. Set
\[
a=f(x),\quad b=f(y),\quad c=g(x),\quad d=g(y).
\]
Since $f$ and $g$ are non-negative and non-increasing,
$a\geq b\geq 0$ and $c\geq d\geq 0$. By convexity,
\[
f(z)\leq \lambda a+(1-\lambda)b,
\qquad
g(z)\leq \lambda c+(1-\lambda)d.
\]
Moreover,
\[
\begin{aligned}
&\lambda ac+(1-\lambda)bd
\\
&\quad-
\bigl(\lambda a+(1-\lambda)b\bigr)
\bigl(\lambda c+(1-\lambda)d\bigr)
\\
&=
\lambda(1-\lambda)(a-b)(c-d)\geq 0.
\end{aligned}
\]
Consequently,
\[
f(z)g(z)
\leq
\lambda f(x)g(x)+(1-\lambda)f(y)g(y),
\]
so $fg$ is convex. Therefore $\mathcal C$ is closed under
multiplication.

It follows that $\mathcal A$ is an algebra. Indeed, if
$f=f_1-f_2$ and $g=g_1-g_2$, with $f_j,g_j\in\mathcal C$, then
\[
fg
=
(f_1g_1+f_2g_2)-(f_1g_2+f_2g_1)\in\mathcal A.
\]

The function
\[
\rho(t)=\frac{1}{1+t}
\]
belongs to $\mathcal C\subset\mathcal A$, is strictly decreasing, and
is nonzero at every point. Hence $\mathcal A$ separates points and
vanishes nowhere.

By the Stone--Weierstrass theorem for $C_0([0,\infty))$, $\mathcal A$
is uniformly dense in $C_0([0,\infty))$. Therefore, for every
$\varepsilon>0$, there exist convex functions
$\psi_1,\psi_2\in C_0([0,\infty))$ such that
\[
\|\psi-(\psi_1-\psi_2)\|_\infty<\varepsilon.
\]
This proves (ii).
\end{proof}


\subsection{General Szeg\H{o} Theorem for compact $X$}

\begin{thm}\label{thm:general_szego}
Let $\sigma\in L^1(X,\nu)$ be real-valued. Suppose $c(\alpha)=\dim \mathcal{H}_\alpha$ for all $\alpha$.
\begin{itemize}
    \item[(i)] If $\sigma$ is bounded, let $\psi$ be continuous on $[\text{ess inf } \sigma, \text{ess sup }\sigma]$.
    \item[(ii)] If $\sigma$ is unbounded and non-negative, let $\psi\in C[0,\infty)$ such that the limit $\lim_{x\to\infty} \psi(x)/x$ exists.
\end{itemize}

Then
\[
    \lim_{\alpha\to\infty} \frac{1}{c(\alpha)} \tr\big(\psi(T_\sigma^\alpha)\big)
    = \int_X \psi(\sigma(x)) \, d\nu(x).
\]
\end{thm}

\begin{proof}
\textbf{Case (i):} By Lemma~\ref{lem:diff_convex} (i), for every $\varepsilon>0$ there exist convex functions $\psi_1,\psi_2$ on $[\text{ess inf } \sigma, \text{ess sup }\sigma]$ such that
\[
    \|\psi-(\psi_1-\psi_2)\|_\infty < \varepsilon.
\]
Then
\[
    \frac{1}{c(\alpha)} \tr(\psi(T_\sigma^\alpha)) 
    = \frac{1}{c(\alpha)} \tr(\psi_1(T_\sigma^\alpha)) 
      - \frac{1}{c(\alpha)} \tr(\psi_2(T_\sigma^\alpha)) 
      + R_\alpha,
\]
where $|R_\alpha| \le \varepsilon$. Applying the convex Szeg\H{o} theorem to $\psi_1,\psi_2$ and letting $\varepsilon\to 0$ gives the result.

\textbf{Case (ii):} 
Let
\[
\beta:=\lim_{x\to\infty}\frac{\psi(x)}{x},
\qquad
r(x):=\psi(x)-\beta x.
\]
Then
\[
\frac{r(x)}{x}\longrightarrow 0
\qquad (x\to\infty).
\]
It is therefore enough to prove the result for $r$.

Fix $\varepsilon>0$. Choose $R>0$ such that
\[
|r(t)|\leq \varepsilon t
\qquad\text{for all }t\geq R.
\]
Let $\chi_R\in C_c([0,\infty))$ satisfy
\[
0\leq\chi_R\leq 1,\qquad
\chi_R(t)=1\ \text{for }0\leq t\leq R,\qquad
\chi_R(t)=0\ \text{for }t\geq 2R.
\]
Set
\[
r_R:=\chi_R r.
\]
Then $r_R\in C_0([0,\infty))$. Moreover,
\[
|r(t)-r_R(t)|\leq \varepsilon t
\qquad\text{for all }t\geq 0.
\]

By Lemma~\ref{lem:diff_convex}, for every $\delta>0$ there exist convex functions
$q_1,q_2\in C_0([0,\infty))$ such that
\[
\|r_R-(q_1-q_2)\|_\infty<\delta.
\]
The convex Szegő theorem applied to $q_1$ and $q_2$ gives
\[
\lim_{\alpha\to\infty}
\frac{1}{c(\alpha)}
\operatorname{Tr}\bigl((q_1-q_2)(T_\sigma^\alpha)\bigr)
=
\int_X(q_1-q_2)(\sigma(x))\,d\nu(x).
\]
Since $c(\alpha)=\dim\mathcal H_\alpha$, the uniform approximation implies
\[
\left|
\frac{1}{c(\alpha)}
\operatorname{Tr}\bigl((r_R-q_1+q_2)(T_\sigma^\alpha)\bigr)
\right|
\leq \delta,
\]
and
\[
\left|
\int_X
\bigl(r_R-(q_1-q_2)\bigr)(\sigma(x))\,d\nu(x)
\right|
\leq \delta\,\nu(X).
\]
Letting $\delta\to0$, we obtain
\[
\lim_{\alpha\to\infty}
\frac{1}{c(\alpha)}
\operatorname{Tr}\bigl(r_R(T_\sigma^\alpha)\bigr)
=
\int_X r_R(\sigma(x))\,d\nu(x).
\]

Finally, since $T_\sigma^\alpha\geq0$,
\[
\left|
\frac{1}{c(\alpha)}
\operatorname{Tr}\bigl((r-r_R)(T_\sigma^\alpha)\bigr)
\right|
\leq
\varepsilon\,
\frac{1}{c(\alpha)}
\operatorname{Tr}(T_\sigma^\alpha)
=
\varepsilon\int_X\sigma\,d\nu,
\]
while
\[
\left|
\int_X(r-r_R)(\sigma(x))\,d\nu(x)
\right|
\leq
\varepsilon\int_X\sigma\,d\nu.
\]
Letting first $\alpha\to\infty$ and then $\varepsilon\to0$ proves
\[
\lim_{\alpha\to\infty}
\frac{1}{c(\alpha)}
\operatorname{Tr}\bigl(r(T_\sigma^\alpha)\bigr)
=
\int_X r(\sigma(x))\,d\nu(x).
\]
Adding the linear term $\beta x$ completes the proof.
\end{proof}

%%%%%%%%%%%%%%%%%%%%%%%%%%%%%%%%%%%%%%%%%%%%%%

\subsection{Determinant-type Szeg\H{o} theorem}

Although Theorem~\ref{thm:general_szego} is formulated for continuous functions, its
conclusion extends to the logarithmic function by a simple
regularization argument.  This yields the determinant formulation of
the first Szeg\H{o} theorem.

\begin{thm}\label{thm:determinant_szego}
Suppose that the assumptions of Theorem~\ref{thm:general_szego} hold and that
$c(\alpha)=\dim H_\alpha$.  Let $\sigma\in L^1(X,\nu)$ be non-negative. Then
\[
 \lim_{\alpha\to\infty}
 \frac{1}{c(\alpha)}
 \text{tr}\left( \log\left(T_\sigma^\alpha\right)\right)
 =
 \int_X \log\sigma(x)\,d\nu(x),
\]
where both sides are understood in the extended real sense. In
particular, the right-hand side belongs to $[-\infty,\infty)$.
\end{thm}

\begin{proof}
For $\varepsilon>0$, let
\[
 \psi_\varepsilon(x)=\log(x+\varepsilon).
\]
Since $\psi_\varepsilon$ is continuous on $[0,\infty)$ and
$\psi_\varepsilon(x)/x\to0$ as $x\to\infty$, Theorem~4.3 gives
\[
 \lim_{\alpha\to\infty}
 \frac{1}{c(\alpha)}
 \tr\log(T_\sigma^\alpha+\varepsilon I)
 =
 \int_X\log(\sigma(x)+\varepsilon)\,d\nu(x).
\]
Since
\[
 T_\sigma^\alpha\leq T_\sigma^\alpha+\varepsilon I,
\]
the monotonicity of the logarithm gives
\[
 \frac{1}{c(\alpha)}\text{tr}\left(\log\left( T_\sigma^\alpha\right)\right)
 \leq
 \frac{1}{c(\alpha)}
 \tr\log(T_\sigma^\alpha+\varepsilon I).
\]
Consequently,
\[
 \limsup_{\alpha\to\infty}
 \frac{1}{c(\alpha)}\text{tr}\left(\log\left( T_\sigma^\alpha\right)\right)
 \leq
 \int_X\log(\sigma+\varepsilon)\,d\nu.
\]
Letting $\varepsilon\downarrow0$ and an application of monotone convergence yields
\[
 \limsup_{\alpha\to\infty}
 \frac{1}{c(\alpha)}\text{tr}\left(\log\left( T_\sigma^\alpha\right)\right)
 \leq
 \int_X\log\sigma\,d\nu,
\]
where the integral is understood in the extended sense.

For the reverse inequality, fix $\varepsilon>0$ and define
\[
\phi_\varepsilon(t)
=
\begin{cases}
-\log(t+\varepsilon), & t\geq 0,\\[2mm]
-\log(\varepsilon)-\dfrac{t}{\varepsilon}, & t<0.
\end{cases}
\]
Then $\phi_\varepsilon$ is convex on $\mathbb{R}$ and agrees with
$-\log(t+\varepsilon)$ on $[0,\infty)$.

Applying the upper Berezin--Lieb inequality to
$\phi_\varepsilon$ gives
\[
\operatorname{Tr}\bigl(-\log(T_\sigma^\alpha+\varepsilon I)\bigr)
\leq
\int_X-\log(\sigma(x)+\varepsilon)\,d\nu_\alpha(x).
\]
Equivalently,
\[
\frac{1}{c(\alpha)}
\operatorname{Tr}\log(T_\sigma^\alpha+\varepsilon I)
\geq
\int_X\log(\sigma(x)+\varepsilon)\,d\nu(x).
\]
Letting $\varepsilon\downarrow0$ and using monotone convergence gives
\[
\frac{1}{c(\alpha)}
\operatorname{Tr}\log(T_\sigma^\alpha)
\geq
\int_X\log(\sigma(x))\,d\nu(x),
\]
where both sides are interpreted in the extended real sense.
\end{proof}



%%%%%%%%%%%%%%%%%%%%%%%%%%%%%%%%%%%%%%%%%%%%
\subsection{Eigenvalue counting consequences} We are in the setting of Theorem~\ref{thm:general_szego}. So $X$ is compact, $(X,\nu)$ is a finite measure space, $\mathcal{H}_\alpha$ are all finite-dimensional  with 
$\dim \mathcal{H}_\alpha=c_\alpha$ for all $\alpha$,
and $T_\sigma^\alpha$ are Toeplitz operators with real-valued $\sigma\in L^1(X,\nu)$ such that either $\sigma$ bounded or $\sigma$ non-negative, and eigenvalues
\(\lambda_1^\alpha\ge\cdots\ge\lambda_{d_\alpha}^\alpha\), where $d_{\alpha}$ counts multiplicity. 

\begin{cor}
For any $t\in\mathbb R$ such that
\(\nu(\sigma=t)=0\),
\[
\lim_{\alpha\to\infty} \frac{1}{\dim \mathcal{H}_\alpha} 
\#\{j:\lambda_j^\alpha>t\}
=
\nu(\{x\in X:\sigma(x)>t\}).
\]
\end{cor}

\begin{proof}
Define the empirical eigenvalue measures
\[
\mu_\alpha = \frac{1}{d_\alpha} \sum_{j=1}^{d_\alpha} \delta_{\lambda_j^\alpha}.
\]
The abstract Szegő theorem gives
\(\mu_\alpha \Rightarrow \sigma_* \nu\) weakly.
By the Portmanteau theorem, for any $t$ with $\nu(\sigma=t)=0$,
\[
\lim_{\alpha\to\infty} \mu_\alpha((t,\infty)) = (\sigma_* \nu)((t,\infty)) = \nu(\sigma>t),
\]
which is exactly the statement.
\end{proof}

%%%%%%%%%%%%%%%%%%%%%%%%%%%%%%%%%%%%%%%%%%%%

\begin{comment}
    

It should be noted that in the absence of Assumption (II) and the convergence of the $\alpha$-Berezin transform, the resulting Szeg\H{o} limit theorems still hold in the following sense.
       
\begin{prop}
    Suppose $\text{dim}(\mathcal{H}_{\alpha})<\infty$ for each $\alpha>0$, and $\text{dim}(\mathcal{H}_{\alpha})/c(\alpha)$ is a bounded sequence in $\alpha$. Additionally assume that $\nu$ is a finite measure. Let $\sigma\in L^{\infty}(X)$. Then there is a measure $\mu$ on $[\inf\sigma,\sup\sigma]$ and a subsequence $\{\alpha_{j}\}$ such that for any $\psi\in C([\inf\sigma,\sup\sigma])$
    \begin{align*}
        \lim_{j\to\infty}\frac{1}{c(\alpha_{j})}\text{tr}\left(\psi\left(T_{\sigma}^{\alpha_{j}}\right)\right)=\int_{[0,\sup\sigma]}\psi(x)d\mu(x).
    \end{align*}
\end{prop}
\begin{proof}
    The proof of this result follows along the exact same lines as the corresponding result in Section 5, with minor modifications. We delay the proof until Section 5. 
\end{proof}


%%%%%%%%%%%%%%%%%%%%%%%%%%%%%%%%%%%%%%%%%%%%%%%%%%%
\subsection{Compact Abelian Groups} The classical case above can easily be extended to general compact Abelian groups (in place of $\T$). This setting was first considered by B\'edos \cite{bedos1996folner} and served as the motivation for two problems posed by N. Nikolski and A. Pushnitski \cite{nikolski2020szeg} in the setting of infinite-dimensional tori. These problems were essentially answered in a recent paper by Guo, Li, and Zhou \cite{guo2024first}. Below, we provide an improvement on the results of B\'edos to non-negative integrable symbols, on compact Abelian groups. The locally compact setting will be considered in Section 5. 

Let $G$ be a compact Abelian group with a normalized Haar measure $\mu$. Let $\hat{G}$ be its dual group. Let $\Gamma_N\subseteq \hat{G}, N\in \mathbb{N}$ be a sequence of finite subsets such that $\Gamma_{N}=\Gamma_{N}^{-1}$ and $\{\Gamma_{N}\}$ satisfy the F\o lner condition, i.e. for any $\eta\in\hat{G}$
\begin{align*}
    \lim_{N\to\infty}\frac{|(\eta\cdot\Gamma_{N})\triangle\Gamma_{N}|}{|\Gamma_{N}|}=0.
\end{align*}
It is known that the F\o lner condition is equivalent to one where $\eta\in\hat{G}$ is replaced by any $K\subset\hat{G}$ compact. 
Let $\mathcal{K}_N$ be a subspace of $L^2(G)$ obtained as a linear span of the characters $\xi\in \Gamma_N$. These characters form an orthonormal basis for this finite-dimensional space. It's easy to see that $\mathcal{K}_N$ is a RKHS with a reproducing kernel $K^N:G\times G\to \C$ given by 
\[K^N(x,y)=\sum_{\xi\in\Gamma_N}\xi(x-y).\] Note that for every $x\in G$ we have \[\norm{K^N_x}^2=K^N(x,x)=|\Gamma_N|.\] We consider the topological space $G$ with Radon measures $\mu_N=|\Gamma_N|\mu$. The normalized reproducing kernels $\{k^N_x\}_{x\in G}$ form a continuous Parseval frame for $\mathcal{K}^N$ in our sense since \[\norm{f}^2=\int_G|f(x)|^2d\mu(x)=\int_G|\ip{f}{k^N_x}|^2d\mu_N(x).\] Assumption (I) is automatically satisfied by our choice of $\mu_N$. We have 

\[|\Gamma_N||\ip{k^N_x}{k^N_y}|^2=\frac{1}{|\Gamma_N|}\left|\sum_{\xi\in\Gamma_N}\xi(x-y)\right|^2.\] 
However, such a kernel is an approximate identity. In fact, this follows from Lemma 3.4 of \cite{coifman1973operators}. From this, we obtain Theorem 1.1.
\end{comment}
%%%%%%%%%%%%%%%%%%%%%%%%%%%%%%%%%%%%%%%%%%%%

\subsection{Compact Abelian Groups}

The classical Szeg\H{o} limit theorem on $\mathbb{T}$ extends naturally to general compact metrizable Abelian groups. This setting was first considered by B\'edos \cite{bedos1996folner} and motivated two problems posed by N. Nikolski and A. Pushnitski \cite{nikolski2020szeg} in the context of infinite-dimensional tori. These problems were essentially resolved in a recent work by Guo, Li, and Zhou \cite{guo2024first}. Here, we present an improvement of B\'edos' results to \emph{non-negative integrable symbols} on compact Abelian groups. The locally compact, non-compact case will be treated in Section~5.

Let $G$ be a compact Abelian group with normalized Haar measure $\mu$, and let $\hat{G}$ denote its dual group. Let $\{\Gamma_N\}_{N\in\mathbb{N}} \subset \hat{G}$ be a sequence of finite subsets satisfying the F\o lner condition:
\[
\forall \eta \in \hat{G}, \qquad \lim_{N\to\infty} \frac{|(\eta \cdot \Gamma_N) \triangle \Gamma_N|}{|\Gamma_N|} = 0.
\]
The metrizability assumption is exactly what permits the sequential formulation: $\widehat G$ is then countable and admits a F{\o}lner sequence. For nonmetrizable compact Abelian groups, the corresponding formulation uses F{\o}lner nets rather than sequences.
It is well known that this condition is equivalent to requiring the limit to hold for all compact subsets $K \subset \hat{G}$.

For each $N$, define the finite-dimensional subspace 
\[
\mathcal{K}_N := \mathrm{span}\{\xi : \xi \in \Gamma_N\} \subset L^2(G),
\] 
whose elements are linear combinations of characters in $\Gamma_N$. These characters form an orthonormal basis of $\mathcal{K}_N$, which is a reproducing kernel Hilbert space with kernel
\[
K^N(x,y) = \sum_{\xi \in \Gamma_N} \xi(x-y), \qquad x,y \in G.
\]
The norm of the kernel at any point satisfies
\[
\|K^N_x\|^2 = K^N(x,x) = |\Gamma_N|.
\]

Endow $G$ with the scaled Radon measures $\mu_N := |\Gamma_N| \, \mu$. Then the normalized reproducing kernels 
\[
k^N_x := \frac{K^N_x}{\|K^N_x\|}, \quad x \in G,
\]
form a continuous Parseval frame for $\mathcal{K}_N$, since for any $f \in \mathcal{K}_N$ we have
\[
\|f\|^2 = \int_G |f(x)|^2 \, d\mu(x) = \int_G |\langle f, k^N_x \rangle|^2 \, d\mu_N(x).
\]
In particular, Assumption (I) of our general Szeg\H{o} framework is automatically satisfied, with $c(N)=|\Gamma_N|=\dim \mathcal{K}_N$.

Moreover, the kernel satisfies
\[
|\Gamma_N| \, |\langle k^N_x, k^N_y \rangle|^2 = \frac{1}{|\Gamma_N|} \left| \sum_{\xi \in \Gamma_N} \xi(x-y) \right|^2,
\]
which is a standard approximate identity on $G$. This follows, for example, from Lemma 3.4 of \cite{coifman1973operators}, and verifies Assumption (II). Hence, the general Szeg\H{o} theorem applies in this setting. Besides Bedos's theorem we also obtain Theorem~\ref{thm:compact_ab_group}, which is a  generalization of~\cite{guo2024first} for all compact Abelian groups.

\begin{proof}[Proof of Theorem~\ref{thm:compact_ab_group}]
The space $\mathcal{K}_N$ and normalized kernels $k^N_x$ satisfy Assumptions (I) and (II) of our general Szeg\H{o} theorem. Moreover, the normalized kernel $|\langle k^N_x, k^N_y \rangle|^2$ acts as an approximate identity on $G$, ensuring that the Berezin transform of $T_\sigma^N$ converges in measure to $\sigma$. The result then follows directly by applying Theorem~\ref{thm:general_szego} to $\mathcal{K}_N$ and $\sigma$.
\end{proof}

As a further consequence of the abstract determinant theorem, we obtain
the determinant formulation of the first Szeg\H{o} theorem on compact
Abelian groups.  In contrast with formulations requiring the symbol to
be bounded away from zero, the following result allows arbitrary
non-negative integrable symbols.

\begin{cor}\label{DetGroup}
Let $G$ be a compact metrizable Abelian group with normalized Haar measure $\mu$,
and let $\{\Gamma_N\}$ be a F{\o}lner sequence as above.  Let
$\sigma\in L^1(G,\mu)$ be non-negative, and let
$T_\sigma^N:K_N\to K_N$ be the corresponding Toeplitz operator.  Then
\[
    \lim_{N\to\infty}
    \frac{1}{|\Gamma_N|}
    \text{tr}\left(\log\left( T_\sigma^N\right)\right)
    =
    \int_G \log\sigma(x)\,d\mu(x),
\]
where the equality is understood in the extended real sense.  In
particular, the right-hand side is allowed to equal $-\infty$.
\end{cor}

\begin{proof}
As shown above, the spaces $K_N$, together with their normalized
reproducing kernels, satisfy Assumptions~(I) and~(II), with
\[
    c(N)=|\Gamma_N|=\dim K_N.
\]
The result therefore follows directly from Theorem~\ref{thm:determinant_szego}.
\end{proof}

This extends the determinant formulation to non-negative integrable
symbols without requiring the symbol to be bounded away from zero.
Together with Theorem~\ref{thm:compact_ab_group}, it shows that both the trace and determinant
forms of the first Szeg\H{o} limit theorem follow from the same
continuous-frame argument.

%%%%%%%%%%%%%%%%%%%%%%%%%%%%%%%%%%%%%%%%%%%%%
%\subsection{Eigenvalue Distribution of Toeplitz Operators}
%As remarked in the introduction, using the Szeg\H o limits developed above one can derive certain Weyl laws for integrable Toeplitz operators. Below, we will derive Weyl laws for such symbols. From this, one can also derive estimates on the plunge region associated to an operator. This plunge region corresponds to the gap between eigenvalues corresponding to $1$ and those corresponding to $0$ for symbols bounded by $1$ \cite{DMFN02,DMN02}. The results to follow very much resemble the work of Feichtinger and Nowak \cite{FN01}, and in fact the proofs are essentially identical. We now define the eigen-counting function. 
%\begin{definition} Let \begin{align*} N_{\alpha}(\sigma,\lambda)=\#\{\lambda_{j}^{\alpha}(\sigma):\lambda_{j}^{\alpha}(\sigma)>\lambda\} \end{align*} denote the number of eigenvalues of $T_{\sigma}^{\alpha}$ greater than $\lambda$. Similarly, we denote the number bounded between $\gamma$ and $\beta$ as $N_{\alpha}(\sigma,\gamma,\beta)$. \end{definition}
% We then obtain the following Weyl law.
%\begin{thm} Let $\sigma\in L^{1}(X,\nu)\cap L^{\infty}(X,\nu)$ be non-negative. Suppose $0<\lambda<\norm{\sigma}_{\infty}$ with  $\nu\left(\{\sigma=\lambda\}\right)=0$. Then, \begin{align*}  \lim_{\alpha\to\infty}\frac{1}{c(\alpha)}N_{\alpha}(\sigma,\lambda)=\nu\left(\{\sigma(x)>\lambda\}\right). \end{align*}
%\end{thm}
%\begin{proof}

%Without loss of generality, we assume that $\norm{\sigma}_{\infty}=1$. Let $\chi=\chi_{\lambda}$ be the indicator function on $[\lambda,1)$ and let $\psi(x)=\chi/x$. Let $\varepsilon>0$. There are continuous functions $\psi_{\varepsilon,1},\psi_{\varepsilon,2}$ such that $\psi_{\varepsilon,1}(x)\leq\chi(x)\leq\psi_{\varepsilon,2}(x)\leq \frac{1}{\lambda}$ and $\psi_{\varepsilon,1}$ and $\psi_{\varepsilon,2}$ coincide with $\psi$ outside of the intervals $[\lambda-\varepsilon,\lambda+\varepsilon/2]$. By Theorem 3.7, we have 
    %\begin{align*}\frac{1}{c(\alpha)}\text{tr}\left(\left(\psi_{\varepsilon,2}-\psi_{\varepsilon,1}\right)\left(T_{\sigma}^{\alpha}\right)\right)\to\int_{X}\sigma(x)\left(\psi_{\varepsilon,2}-\psi_{\varepsilon,1}\right)\left(\sigma(x)\right)d\nu(x), \end{align*} and \begin{align*} \int_{X}\sigma(x)\left(\psi_{\varepsilon,2}-\psi_{\varepsilon,1}\right)\left(\sigma(x)\right)d\nu(x)\leq\nu\left(\{\lambda-\varepsilon<\sigma<\lambda+\varepsilon\}\right)\frac{1}{\lambda}. \end{align*}Thus, letting $\varepsilon\to 0$ we see that the limit in the statement of this theorem exists, and must equal the desired quantity. \end{proof}
%Using the same proof strategy will give the corresponding statements for plunge regions as well. Additionally, using the essentially the same proof strategy, we obtain the following corollary in the context of Corollary 3.7.\begin{cor} Let $\sigma\in L^{\infty}(X,\nu)$ with the assumptions of Corollary 4.4. Suppose $\inf\sigma<\gamma<\beta<\sup\sigma$ with either $\gamma>0$ or $\beta<0$. Additionally, assume that $\nu\left(\{\sigma=\gamma\}\right)=\nu\left(\{\sigma=\beta\}\right)=0$. Then,\begin{align*}  \lim_{\alpha\to\infty}\frac{1}{c(\alpha)}N_{\alpha}(\sigma,\gamma,\beta)=\nu\left(\gamma<\sigma(x)<\beta\right).\end{align*}\end{cor}
%We now discuss several applications of the above results. 
% \begin{thm}
%     Let $\sigma$ be non-negative, and $\sigma\in L^{p}(X,\nu)\cap L^{\infty}(X,\nu)$ for some $1\leq p<\infty$. Then, for $0<\lambda<\|\sigma\|_{\infty}$ such that $\nu\left(\{x:\sigma(x)=\lambda\}\right)=0$, we have 
%     \begin{align*}
%         N_{\alpha}^{>}(\lambda)\sim \nu_{\alpha}\left(\{\sigma>\lambda\}\right),
%     \end{align*}
%     as $\alpha\to\infty$. 
% \end{thm}
% \begin{proof}[Proof of Theorem $4.2$]
%     Without loss of generality, suppose $\norm{\sigma}_{\infty}=1$. Let $\psi(r)=\chi_{(\lambda,1]}$, let $\varepsilon>0$, and let $\psi_{-}$ and $\psi_{+}$ be two continuous functions such that 
%     \begin{align*}
%         0\leq\psi_{-}(x)\leq\psi(x)\leq\psi_{+}(x)\leq 1/\lambda,
%     \end{align*}
%     and $\psi_{-}$ and $\psi_{+}$ agree with $\psi$ outside of $[\lambda-\varepsilon,\lambda+\varepsilon]$. Clearly, we have that 
%     \begin{align*}
%         \text{tr}\left(\psi_{-}\left(T_{\sigma}^{\alpha}\right)\right)\leq N(\lambda)\leq&\text{tr}\left(\psi_{+}\left(T_{\sigma}^{\alpha}\right)\right).
%     \end{align*}
%     Thus, by Corollary $2.5$, we have that as $\alpha\to\infty$
%     \begin{align*}
%         \frac{\text{tr}\left(\left(\psi_{+}-\psi_{-}\right)\left(T_{\sigma}^{\alpha}\right)\right)}{c(\alpha)}\to&\int_{X}\left(\psi_{+}-\psi_{-}\right)(\sigma(x))d\nu(x)\\
%         \leq&\frac{1}{\lambda}\nu\left(\{\lambda-\varepsilon\leq \sigma(x)\leq\lambda+\varepsilon\}\right).
%     \end{align*}
%     Now, letting $\varepsilon\to0$, this clearly converges to zero. From this, it is clear to see the conclusion. 
% \end{proof}
% In a similar fashion to the proof above, one can also obtain the following asymptotic estimate on the plunge region.
% \begin{cor}
% Let $\lambda_{1},\lambda_{2}\in (0,1)$ and 
% \begin{align*}
%     M_{\alpha}(\lambda_{1},\lambda_{2})=\#\{\lambda_{j}(\alpha):\lambda_{1}<\lambda_{j}<\lambda_{2}\}.
% \end{align*}
% Then, if $\nu\left(\{\sigma=\lambda_{i}\}\right)=0$ for $i=1,2$,
% \begin{align*}
%     M_{\alpha}(\lambda_{1},\lambda_{2})\sim\nu_{\alpha}\left(\{\lambda_{1}<\sigma(x)<\lambda_{2}\}\right).
% \end{align*}
% \end{cor}

%%%%%%%%%%%%%%%%%%%%%%%%%%%%%%%%%%%%%%%%%%%%%%%%%%%%%
\begin{comment}
\section{Szeg\H{o} Limit Theorems for Infinite Measures}
 Applying the same lines of reasoning from the previous section directly with Proposition 2.1, we obtain the following result.
\begin{thm}
    Suppose $\sigma\in L^{1}(X,\nu)\cap L^{\infty}(X,\nu)$ with $\sigma$ real-valued and $\eta\in L^{1}(X,\nu)$ with $\eta\geq 0$. Suppose additionally that $T_{\sigma}^{\alpha}$ and $T_{\eta}^{\alpha}$ are simultaneously diagonalizable by a common orthonormal basis, and $\Tilde{\sigma}^{\alpha}\to\sigma$, $\Tilde{\eta}^{\alpha}\to\eta$ in $\nu$-measure as $\alpha\to\infty$. Let $\psi:\R\to\mathbb{R}$ be continuous on the range of $\sigma$. Then,
    \begin{align*}
        \lim_{\alpha\to\infty}\frac{1}{c(\alpha)}\text{tr}\left(T_{\eta}^{\alpha}\psi\left(T_{\sigma}^{\alpha}\right)\right)=\int_{X}\psi\left(\sigma(x)\right)\eta(x)d\nu(x).
    \end{align*}
\end{thm}
\begin{proof}
    We consider when $\psi$ is convex first. Using Proposition 2.1, we have 
    \begin{align*}
        \int_{X}\psi\left(\Tilde{\sigma}^{\alpha}(x)\right)\eta(x)d\nu(x)\leq\frac{1}{c(\alpha)}\text{tr}\left(T_{\eta}^{\alpha}\psi\left(T_{\sigma}^{\alpha}\right)\right)\leq\int_{X}\psi\left(\sigma(x)\right)\Tilde{\eta}^{\alpha}(x)d\nu(x).
    \end{align*}
    Note that $\psi$ is continuous on a compact interval and hence bounded. Additionally, the $\alpha$-Berezin transform is an $L^{\infty}(X,\nu)$-contraction (see Lemma 3.1). Thus,
    \begin{align*}
        \sup_{x\in[\inf\Tilde{\sigma}^{\alpha},\sup\Tilde{\sigma}^{\alpha}]}|\psi(x)|\leq\sup_{x\in[-\norm{\sigma}_{\infty},\norm{\sigma}_{\infty}]}|\psi(x)|.
    \end{align*}
    From this, we obtain
    \begin{align*}
        \left|\psi\left(\Tilde{\sigma}^{\alpha}(x)\right)\eta(x)\right|\leq\norm{\psi(\Tilde{\sigma}^{\alpha})}_{\infty}\eta(x)\leq \norm{\psi}_{L^{\infty}([-\norm{\sigma}_{\infty},\norm{\sigma}_{\infty}])}\eta(x).  
    \end{align*}
    By applying the Dominated Convergence Theorem the lower bound becomes
    \begin{align*}
        \int_{X}\psi\left(\sigma(x)\right)\eta(x)d\nu(x)\leq\liminf_{\alpha\to\infty}\frac{1}{c(\alpha)}\text{tr}\left(T_{\eta}^{\alpha}\psi\left(T_{\sigma}^{\alpha}\right)\right).
    \end{align*}
    The upper bound requires more work. Note that $\left|\psi(\sigma(x))\Tilde{\eta}^{\alpha}(x)\right|\leq\norm{\psi(\sigma)}_{\infty}\left|\Tilde{\eta}^{\alpha}(x)\right|$, with $\Tilde{\eta}^{\alpha}$ integrable, $\Tilde{\eta}^{\alpha}\to\eta$ in $\nu$-measure, and 
    \begin{align*}
        \int_{X}\Tilde{\eta}^{\alpha}(x)d\nu(x)=\int_{X}\eta(y)\int_{X}\left|\ip{k_{x}^{\alpha}}{k_{y}^{\alpha}}_{\alpha}\right|^{2}d\nu_{\alpha}(x)d\nu(y)=\int_{X}\eta(y)d\nu(y).
    \end{align*}
    Thus, by the Dominated Convergence Theorem, 
    \begin{align*}
        \limsup_{\alpha\to\infty}\frac{1}{c(\alpha)}\text{tr}\left(T_{\eta}^{\alpha}\psi\left(T_{\sigma}^{\alpha}\right)\right)\leq&\limsup_{\alpha\to\infty}\int_{X}\psi(\sigma(x))\Tilde{\eta}^{\alpha}(x)d\nu(x)\\
        =&\int_{X}\psi(\sigma(x))\eta(x)d\nu(x),
    \end{align*}
    and the result is established. The result is extended to $\psi$ continuous on the range of $\sigma$ by using the same strategy for the monomials $\psi(x)=x^k$ as in the proof of Corollary 4.2, extending to polynomials by linearity of the trace, appealing to Lemma 4.3 (3), and applying the Stone-Weierstrauss Theorem.
\end{proof}
Transitioning back to Theorems 1.2 and 1.3, we obtain the following result.
\begin{cor}\label{Sze_Berg}
    Suppose $\sigma\in L^{1}(X,\nu)\cap L^{\infty}(X,\nu)$ is non-negative such that $\Tilde{\sigma}^{\alpha}\to\sigma$ in $\nu$-measure as $\alpha\to\infty$. Then, for $\psi:\R\to\R$ continuous on the range of $\sigma$, 
    \begin{align*}
        \lim_{\alpha\to\infty}\frac{1}{c(\alpha)}\text{tr}\left(T_{\sigma}\psi\left(T_{\sigma}\right)\right)=\int_{X}\sigma(x)\psi(\sigma(x))d\nu(x).
    \end{align*}
\end{cor}
% \begin{proof}
%     For this proof, we follow the strategy set out in Corollary 4.2. However, things are simplified greatly due to the non-negativity of the symbol. In this case, the monomials $x^{n}$ are convex and uniformly continuous on the interval $[0,\norm{\sigma}_{\infty}]$ and by linearity of the trace we obtain 
%     \begin{align*}
%         \lim_{\alpha\to\infty}\frac{1}{c(\alpha)}\text{tr}\left(T_{\sigma}^{\alpha}p\left(T_{\sigma}^{\alpha}\right)\right)=\int_{X}\sigma(x)p(\sigma(x))d\nu(x)
%     \end{align*}
    % for any polynomial $p$. Thus, applying Lemma 4.3 and the Stone-Weierstrass Theorem gives the result. 
% \end{proof}
    Finally, we consider the case when the Berezin transform doesn't converge. 
 \begin{prop}
    Suppose $\sigma\in L^{1}(X,\nu)\cap L^{\infty}(X,\nu)$ is non-negative. Then there is a measure $\mu$ on $[0,\sup\sigma]$ and a subsequence $\{\alpha_{j}\}$ such that for any $\psi\in C([0,\sup\sigma])$
    \begin{align*}
        \lim_{j\to\infty}\frac{1}{c(\alpha_{j})}\text{tr}\left(T_{\sigma}^{\alpha_{j}}\psi\left(T_{\sigma}^{\alpha_{j}}\right)\right)=\int_{[0,\sup\sigma]}\psi(x)d\mu(x).
    \end{align*}
\end{prop}
\begin{proof}
    By Proposition 2.1, we have for any $\alpha>0$ and $k\in\mathbb{N}$
    \begin{align*}
        \frac{1}{c(\alpha)}\text{tr}\left(T_{\sigma}^{\alpha}\left(T_{\sigma}^{\alpha}\right)^{k}\right)\leq\int_{X}\Tilde{\sigma}^{\alpha}(x)\left(\sigma(x)\right)^{k}d\nu(x)\leq\norm{\sigma}_{\infty}^{k}\norm{\sigma}_{1}<\infty,
    \end{align*}
    since $x\mapsto x^{k}$ is convex. Thus, the sequence $\frac{1}{c(\alpha)}\text{tr}\left(T_{\sigma}^{\alpha}\right)$ is a bounded sequence. By the Bolzano-weierstrass Theorem, we can extract a convergent subsequence, say $\{\alpha_{0,j}\}$. Using this subsequence, we see that $\frac{1}{c(\alpha_{0,j})}\text{tr}\left(T_{\sigma}^{\alpha_{0,j}}T_{\sigma}^{\alpha_{0,j}}\right)$ is also a bounded sequence. Applying the Bolzano-weierstrass Theorem, we obtain a subsequence $\{\alpha_{1,j}\}$. Continuing this process, we obtain a family of subsequences $\{\alpha_{k,j}\}$. Letting $\{\alpha_{j}\}=\{\alpha_{j,j}\}$ be the diagonal sequence, we see that for each $k\in\mathbb{N}$, we have 
    \begin{align*}
        \lim_{j\to\infty}\frac{1}{c(\alpha_{j})}\text{tr}\left(T_{\sigma}^{\alpha_{j}}\left(T_{\sigma}^{\alpha_{j}}\right)^{k}\right)
    \end{align*}
    exists. The linearity of the trace then extends this result to polynomials, while the result is then extended to continuous functions by the same Stone-weierstrass argument as above. Then, applying the Riesz-Kakutani Representation theorem \cite{lax2014functional} gives the result. 
\end{proof}

\end{comment}
%%%%%%%%%%%%%%%%%%%%%%%%%%%%%%%%%%%%%%%%%%%%%%%

\section{Abstract Szeg\H{o} Limit Theorem for Locally Compact Spaces}

Throughout this section let $X$ be locally compact and let
$\sigma,\eta \in L^1(X,\nu)$ be real-valued.
As usual we assume that the normalized measures satisfy

\[
d\nu_\alpha = c(\alpha)\, d\nu,
\qquad\text{(Assumption (I))}.
\]

We also assume that the Berezin transforms satisfy

\[
\widetilde{\sigma}^{\alpha} \to \sigma,
\qquad
\widetilde{\eta}^{\alpha} \to \eta
\quad\text{in $\nu$-measure as } \alpha\to\infty,
\]
which we know is true under (Assumption (II)). 

Finally, we assume that the Toeplitz operators $T_\sigma^\alpha$ and $T_\eta^\alpha$ commute for all $\alpha$.

\subsection{Convex Szeg\H{o} Limit}

\begin{thm}\label{thm:convex_szego_lc} Let $\sigma,\eta\in L^1(X,\nu)\cap L^2(X,\nu)$ be real-valued with $\eta$ non-negative.
Assume that the Toeplitz operators $T_\sigma^\alpha$ and
$T_\eta^\alpha$ commute for all $\alpha$. 
\leavevmode
\begin{itemize}
\item[(i)] Suppose $\sigma \ge 0$ and $\psi \in C([0,\infty))$
is convex such that the limit
\[
\lim_{x\to\infty} \frac{\psi(x)}{x}
\]
exists. Then
\[
\lim_{\alpha\to\infty}
\frac{1}{c(\alpha)}
\tr\!\big(T_\eta^\alpha \psi(T_\sigma^\alpha)\big)
=
\int_X \eta(x)\psi(\sigma(x))\,d\nu(x).
\]

\item[(ii)] If $\sigma \in L^\infty(X,\nu)$ and
$\psi$ is convex and continuous on
$[\mathrm{ess\,inf}\,\sigma,\mathrm{ess\,sup}\,\sigma]$,
then the same limit holds.
\end{itemize}
\end{thm}

\begin{proof} We prove (i); part (ii) follows by a simpler argument.

\medskip
\noindent
\textbf{Step 1: The linear case.}

For $\psi(x)=x$ we have
\[
\frac{1}{c(\alpha)}
\tr\!\big(T_\eta^\alpha T_\sigma^\alpha\big)
=
\int_X \eta(x)\widetilde{\sigma}^\alpha(x)\, d\nu(x),
\]
using Assumption (I).

By Proposition~\ref{Lp},
\[
\widetilde{\sigma}^{\,\alpha}\longrightarrow\sigma
\qquad\text{in }L^2(X,\nu).
\]
Therefore, by the Cauchy--Schwarz inequality,
\[
\left|
\int_X\eta(x)
\bigl(\widetilde{\sigma}^{\,\alpha}(x)-\sigma(x)\bigr)
\,d\nu(x)
\right|
\leq
\|\eta\|_{L^2(X,\nu)}
\left\|
\widetilde{\sigma}^{\,\alpha}-\sigma
\right\|_{L^2(X,\nu)}
\longrightarrow 0.
\]
Consequently,
\[
\frac{1}{c(\alpha)}
\operatorname{Tr}\!\left(
T_\eta^\alpha T_\sigma^\alpha
\right)
\longrightarrow
\int_X\eta(x)\sigma(x)\,d\nu(x).
\]

\medskip
\noindent
\textbf{Step 2: Reduction to sublinear growth.}

Write
\[
\psi(x)=\beta x + \psi_0(x),
\qquad
\text{where } \frac{\psi_0(x)}{x}\to 0.
\]
By linearity it suffices to treat the case
$\psi(x)/x\to 0$.

\medskip
\noindent
\textbf{Step 3: Berezin--Lieb inequality.}

For convex $\psi$ the weighted Berezin--Lieb inequality yields
\[
\int_X \eta \psi(\widetilde{\sigma}^\alpha)\, d\nu_\alpha
\le
\tr\!\big(T_\eta^\alpha \psi(T_\sigma^\alpha)\big)
\le
\int_X \widetilde{\eta}^\alpha \psi(\sigma)\, d\nu_\alpha.
\]

Dividing by $c(\alpha)$ and using Assumption (I),
\[
\int_X \eta \psi(\widetilde{\sigma}^\alpha)\, d\nu
\le
\frac{1}{c(\alpha)}
\tr\!\big(T_\eta^\alpha \psi(T_\sigma^\alpha)\big)
\le
\int_X \widetilde{\eta}^\alpha \psi(\sigma)\, d\nu.
\]

\medskip

\noindent\textbf{Step 4: Convergence of both bounds.}

Since $\psi_0(x)/x\to0$, there exists a constant $C>0$ such that
\[
|\psi_0(t)|\leq t+C,
\qquad t\geq0.
\]

By Proposition~\ref{Lp},
\[
\widetilde{\sigma}^{\,\alpha}\to\sigma
\quad\text{and}\quad
\widetilde{\eta}^{\,\alpha}\to\eta
\qquad\text{in }L^2(X,\nu).
\]
Therefore,
\[
\left\|
\eta\bigl(\widetilde{\sigma}^{\,\alpha}-\sigma\bigr)
\right\|_{L^1}
\leq
\|\eta\|_{L^2}
\left\|
\widetilde{\sigma}^{\,\alpha}-\sigma
\right\|_{L^2}
\longrightarrow0.
\]
Thus
\[
\eta\,\widetilde{\sigma}^{\,\alpha}
\longrightarrow
\eta\sigma
\qquad\text{in }L^1(X,\nu).
\]
Moreover,
\[
\bigl|\eta\,\psi_0(\widetilde{\sigma}^{\,\alpha})\bigr|
\leq
\bigl|\eta\,\widetilde{\sigma}^{\,\alpha}\bigr|+C\bigl|\eta\bigr|.
\]
The right-hand side is a uniformly integrable family. Since
\[
\eta\,\psi_0(\widetilde{\sigma}^{\,\alpha})
\longrightarrow
\eta\,\psi_0(\sigma)
\]
in measure, Vitali's convergence theorem gives
\[
\int_X
\eta\,\psi_0(\widetilde{\sigma}^{\,\alpha})\,d\nu
\longrightarrow
\int_X\eta\,\psi_0(\sigma)\,d\nu.
\]

Similarly,
\[
\left\|
\sigma\bigl(\widetilde{\eta}^{\,\alpha}-\eta\bigr)
\right\|_{L^1}
\leq
\|\sigma\|_{L^2}
\left\|
\widetilde{\eta}^{\,\alpha}-\eta
\right\|_{L^2}
\longrightarrow0.
\]
Hence
\[
\widetilde{\eta}^{\,\alpha}\sigma
\longrightarrow
\eta\sigma
\qquad\text{in }L^1(X,\nu),
\]
and
\[
\bigl|
\widetilde{\eta}^{\,\alpha}\psi_0(\sigma)
\bigr|
\leq
\bigl|\widetilde{\eta}^{\,\alpha}\sigma\bigr|
+
C\bigl|\widetilde{\eta}^{\,\alpha}\bigr|.
\]
Again, the right-hand side is uniformly integrable. Since
\[
\widetilde{\eta}^{\,\alpha}\psi_0(\sigma)
\longrightarrow
\eta\psi_0(\sigma)
\]
in measure, Vitali's theorem yields
\[
\int_X
\widetilde{\eta}^{\,\alpha}\psi_0(\sigma)\,d\nu
\longrightarrow
\int_X\eta\psi_0(\sigma)\,d\nu.
\]

The squeeze theorem now proves the claim.
\end{proof}

\subsection{General Continuous Functions}

\begin{thm}\label{thm:general_szego_lc}
Let $\sigma,\eta\in L^1(X,\nu)\cap L^2(X,\nu)$ be real-valued with $\eta$ non-negative. Assume that the Toeplitz operators $T_\sigma^\alpha$ and $T_\eta^\alpha$ commute for all $\alpha$.

\begin{enumerate}
\item[(i)] If $\sigma$ is bounded and
$\psi$ is continuous on
$[\mathrm{ess\,inf}\,\sigma,\mathrm{ess\,sup}\,\sigma]$,
then
\[
\lim_{\alpha\to\infty}
\frac{1}{c(\alpha)}
\tr\!\big(T_\eta^\alpha \psi(T_\sigma^\alpha)\big)
=
\int_X \eta(x)\psi(\sigma(x))\, d\nu(x).
\]

\item[(ii)] If $\sigma\ge0$ and
$\psi\in C([0,\infty))$ such that the limit
$\lim_{x\to\infty}\psi(x)/x$ exists,
then the same limit holds.
\end{enumerate}
\end{thm}

\begin{proof}
We use Lemma~\ref{lem:diff_convex} together with
Theorem~\ref{thm:convex_szego_lc}.

\medskip
\noindent
\textbf{Case (i).}
Assume $\sigma$ is bounded and $\psi$ is continuous on
$[\mathrm{ess\,inf}\,\sigma,\mathrm{ess\,sup}\,\sigma]$.

By Lemma~\ref{lem:diff_convex}(i), for every $\varepsilon>0$
there exist convex functions $\psi_1,\psi_2$ on this interval such that
\[
\|\psi-(\psi_1-\psi_2)\|_\infty < \varepsilon.
\]
Set $\varphi := \psi-(\psi_1-\psi_2)$.
Then $\|\varphi\|_\infty < \varepsilon$.

By functional calculus,
\[
\psi(T_\sigma^\alpha)
=
(\psi_1-\psi_2)(T_\sigma^\alpha)
+
\varphi(T_\sigma^\alpha).
\]
Hence
\[
\frac{1}{c(\alpha)}
\tr\!\big(T_\eta^\alpha \psi(T_\sigma^\alpha)\big)
=
\frac{1}{c(\alpha)}
\tr\!\big(T_\eta^\alpha (\psi_1-\psi_2)(T_\sigma^\alpha)\big)
+
R_\alpha,
\]
where
\[
R_\alpha
=
\frac{1}{c(\alpha)}
\tr\!\big(T_\eta^\alpha \varphi(T_\sigma^\alpha)\big).
\]

To estimate $R_\alpha$ we use the trace inequality
\[
|\tr(AB)| \le \|B\|\, \tr(|A|),
\]
valid for trace-class $A$ and bounded $B$.
Since $\|\varphi(T_\sigma^\alpha)\|
\le \|\varphi\|_\infty < \varepsilon$,
we obtain
\[
|R_\alpha|
\le
\varepsilon \,
\frac{1}{c(\alpha)} \tr(|T_\eta^\alpha|).
\]

Since $\eta\in L^1(X,\nu)$, write
\[
\eta=\eta_+-\eta_-,
\qquad
\eta_+,\eta_-\geq0.
\]
Then
\[
\begin{aligned}
\|T_\eta^\alpha\|_{\mathfrak S_1}
&\leq
\|T_{\eta_+}^\alpha\|_{\mathfrak S_1}
+
\|T_{\eta_-}^\alpha\|_{\mathfrak S_1}\\
&=
\operatorname{Tr}(T_{\eta_+}^\alpha)
+
\operatorname{Tr}(T_{\eta_-}^\alpha)\\
&=
\int_X|\eta(x)|\,d\nu_\alpha(x)\\
&=
c(\alpha)\int_X|\eta(x)|\,d\nu(x).
\end{aligned}
\]
Consequently,
\[
|R_\alpha|
\leq
\varepsilon\,
\frac{\|T_\eta^\alpha\|_{\mathfrak S_1}}{c(\alpha)}
\leq
\varepsilon\int_X|\eta(x)|\,d\nu(x).
\]

On the other hand, Theorem~\ref{thm:convex_szego_lc} applied to the convex functions
$\psi_1$ and $\psi_2$ gives
\[
\frac{1}{c(\alpha)}
\operatorname{Tr}\!\left(
T_\eta^\alpha
(\psi_1-\psi_2)(T_\sigma^\alpha)
\right)
\longrightarrow
\int_X
\eta(x)(\psi_1-\psi_2)(\sigma(x))
\,d\nu(x).
\]

Since
\[
\left\|
\psi-(\psi_1-\psi_2)
\right\|_\infty<\varepsilon,
\]
we have
\[
\left|
\int_X
\eta(x)
\bigl((\psi_1-\psi_2)(\sigma(x))-\psi(\sigma(x))\bigr)
\,d\nu(x)
\right|
\leq
\varepsilon\int_X|\eta(x)|\,d\nu(x).
\]

Combining this estimate with the bound for $R_\alpha$, we obtain
\[
\begin{aligned}
&\limsup_{\alpha\to\infty}
\left|
\frac{1}{c(\alpha)}
\operatorname{Tr}\!\left(
T_\eta^\alpha\psi(T_\sigma^\alpha)
\right)
-
\int_X\eta(x)\psi(\sigma(x))\,d\nu(x)
\right|
\\
&\qquad\leq
2\varepsilon\int_X|\eta(x)|\,d\nu(x).
\end{aligned}
\]
Since $\varepsilon>0$ is arbitrary, it follows that
\[
\lim_{\alpha\to\infty}
\frac{1}{c(\alpha)}
\operatorname{Tr}\!\left(
T_\eta^\alpha\psi(T_\sigma^\alpha)
\right)
=
\int_X\eta(x)\psi(\sigma(x))\,d\nu(x).
\]


\medskip
\noindent
\textbf{Case (ii).}

Let
\[
\beta:=\lim_{x\to\infty}\frac{\psi(x)}{x},
\qquad
r(x):=\psi(x)-\beta x.
\]
Then
\[
\frac{r(x)}{x}\longrightarrow0
\qquad (x\to\infty).
\]

Fix $\varepsilon>0$. Choose $R>0$ such that
\[
|r(t)|\leq\varepsilon t
\qquad\text{for all }t\geq R.
\]
Let $\chi_R\in C_c([0,\infty))$ satisfy
\[
0\leq\chi_R\leq1,\qquad
\chi_R(t)=1\ \text{for }0\leq t\leq R,\qquad
\chi_R(t)=0\ \text{for }t\geq2R.
\]
Set
\[
r_R:=\chi_R r.
\]
Then $r_R\in C_0([0,\infty))$ and
\[
|r(t)-r_R(t)|\leq\varepsilon t
\qquad\text{for all }t\geq0.
\]

By Lemma~\ref{lem:diff_convex}(ii), for every $\delta>0$ there exist convex functions
$q_1,q_2\in C_0([0,\infty))$ such that, with
$q:=q_1-q_2$,
\[
\|r_R-q\|_\infty<\delta.
\]
Theorem~\ref{thm:convex_szego_lc} applied to $q_1$ and $q_2$ gives
\[
\frac{1}{c(\alpha)}
\operatorname{Tr}\!\left(
T_\eta^\alpha q(T_\sigma^\alpha)
\right)
\longrightarrow
\int_X\eta(x)q(\sigma(x))\,d\nu(x).
\]

Since $\eta\geq0$,
\[
\frac{1}{c(\alpha)}
\left|
\operatorname{Tr}\!\left(
T_\eta^\alpha(r_R-q)(T_\sigma^\alpha)
\right)
\right|
\leq
\delta\,
\frac{\operatorname{Tr}(T_\eta^\alpha)}{c(\alpha)}
=
\delta\int_X\eta\,d\nu,
\]
and
\[
\left|
\int_X\eta(x)(r_R-q)(\sigma(x))\,d\nu(x)
\right|
\leq
\delta\int_X\eta\,d\nu.
\]
Letting $\delta\to0$, we obtain
\[
\lim_{\alpha\to\infty}
\frac{1}{c(\alpha)}
\operatorname{Tr}\!\left(
T_\eta^\alpha r_R(T_\sigma^\alpha)
\right)
=
\int_X\eta(x)r_R(\sigma(x))\,d\nu(x).
\]

Because $T_\eta^\alpha$ and $T_\sigma^\alpha$ commute and
$T_\eta^\alpha\geq0$, they have a common orthonormal eigenbasis.
Using $|r-r_R|\leq\varepsilon\,\mathrm{id}$ on $[0,\infty)$, we obtain
\[
\left|
\operatorname{Tr}\!\left(
T_\eta^\alpha(r-r_R)(T_\sigma^\alpha)
\right)
\right|
\leq
\varepsilon\,
\operatorname{Tr}\!\left(
T_\eta^\alpha T_\sigma^\alpha
\right).
\]
Hence, by the linear case of Theorem~\ref{thm:convex_szego_lc},
\[
\limsup_{\alpha\to\infty}
\frac{1}{c(\alpha)}
\left|
\operatorname{Tr}\!\left(
T_\eta^\alpha(r-r_R)(T_\sigma^\alpha)
\right)
\right|
\leq
\varepsilon\int_X\eta(x)\sigma(x)\,d\nu(x).
\]
Similarly,
\[
\left|
\int_X\eta(x)(r-r_R)(\sigma(x))\,d\nu(x)
\right|
\leq
\varepsilon\int_X\eta(x)\sigma(x)\,d\nu(x).
\]
Letting first $\alpha\to\infty$, then $\delta\to0$, and finally
$\varepsilon\to0$, gives
\[
\lim_{\alpha\to\infty}
\frac{1}{c(\alpha)}
\operatorname{Tr}\!\left(
T_\eta^\alpha r(T_\sigma^\alpha)
\right)
=
\int_X\eta(x)r(\sigma(x))\,d\nu(x).
\]

Finally, applying the linear case to the term $\beta x$ yields
\[
\lim_{\alpha\to\infty}
\frac{1}{c(\alpha)}
\operatorname{Tr}\!\left(
T_\eta^\alpha\psi(T_\sigma^\alpha)
\right)
=
\int_X\eta(x)\psi(\sigma(x))\,d\nu(x).
\]

%%%%%%%%%%%%%%%%%%%%%%%%%%%%%

\end{proof}



%%%%%%%%%%%%%%%%%%%%%%%%%%%%%%%%%%%%%%%%%%%%%%%%

\begin{comment}
\subsection{Eigenvalue Distribution of Toeplitz Operators}
As remarked in the introduction, using the Szeg\H o limits developed above one can derive certain Weyl laws for integrable Toeplitz operators. Below, we will derive Weyl laws for such symbols. The results to follow very much resemble the work of Feichtinger and Nowak \cite{FN01}, and in fact the proofs are essentially identical. We now define the eigen-counting function. 
\begin{definition} 
Let 
    \begin{align*} 
        N_{\alpha}(\sigma,\lambda)=\#\{\lambda_{j}^{\alpha}(\sigma):\lambda_{j}^{\alpha}(\sigma)>\lambda\} 
    \end{align*} 
denote the number of eigenvalues of $T_{\sigma}^{\alpha}$ greater than $\lambda$. 
\end{definition}
We then obtain the following Weyl law.
\begin{thm} 
Let $\sigma\in L^{1}(X,\nu)\cap L^{\infty}(X,\nu)$ be non-negative. Suppose $0<\lambda<\norm{\sigma}_{\infty}$ with  $\nu\left(\{\sigma=\lambda\}\right)=0$. Then, 
    \begin{align*}  
        \lim_{\alpha\to\infty}\frac{1}{c(\alpha)}N_{\alpha}(\sigma,\lambda)=\nu\left(\{\sigma(x)>\lambda\}\right). 
    \end{align*}
\end{thm}
\begin{proof}
Without loss of generality, we assume that $\norm{\sigma}_{\infty}=1$. Let $\chi=\chi_{\lambda}$ be the indicator function on $[\lambda,1)$ and let $\psi(x)=\chi/x$. Let $\varepsilon>0$. There are continuous functions $\psi_{\varepsilon,1},\psi_{\varepsilon,2}$ such that $\psi_{\varepsilon,1}(x)\leq\chi(x)\leq\psi_{\varepsilon,2}(x)\leq \frac{1}{\lambda}$ and $\psi_{\varepsilon,1}$ and $\psi_{\varepsilon,2}$ coincide with $\psi$ outside of the intervals $[\lambda-\varepsilon,\lambda+\varepsilon/2]$. By Corollary 5.2, we have 
    \begin{align*}
        \frac{1}{c(\alpha)}\text{tr}\left(\left(\psi_{\varepsilon,2}-\psi_{\varepsilon,1}\right)\left(T_{\sigma}^{\alpha}\right)\right)\to\int_{X}\sigma(x)\left(\psi_{\varepsilon,2}-\psi_{\varepsilon,1}\right)\left(\sigma(x)\right)d\nu(x), 
    \end{align*} 
        and 
    \begin{align*} 
        \int_{X}\sigma(x)\left(\psi_{\varepsilon,2}-\psi_{\varepsilon,1}\right)\left(\sigma(x)\right)d\nu(x)\leq\nu\left(\{\lambda-\varepsilon<\sigma<\lambda+\varepsilon\}\right)\frac{1}{\lambda}. 
    \end{align*}
        Thus, letting $\varepsilon\to 0$ we see that the limit in the statement of this theorem exists, and must equal the desired quantity. 
\end{proof}

\end{comment}

\begin{comment}

We next extend the locally compact Szeg\H{o} theorem to the case $\eta\equiv 1$ (hence not integrable) by an approximation argument with integrable weights. The extension requires only a mild local linearity at the origin for the test function $\psi$.

\begin{thm}[Szeg\H{o} limit for $\eta\equiv 1$ by $L^1$-approximation]
\label{thm:Szego-eta1}
Assume the standing hypotheses of this section (Assumptions~(I)–(II) and the normalized continuous Parseval frame). Let $\sigma\in L^1(X,\nu)$ be real-valued and bounded, and let $T^\alpha_\sigma$ denote the corresponding Toeplitz operators on $H_\alpha$. Suppose $\psi\colon I\to\mathbb{R}$ is continuous on an interval $I$ containing the essential range of $\sigma$ and satisfies the local linear bound
\begin{equation}\label{eq:psi-linear-local}
|\psi(\lambda)|\le C\,|\lambda|\qquad\text{for }|\lambda|<\delta,
\end{equation}
for some $C,\delta>0$. Then $\psi(T^\alpha_\sigma)$ is trace class for every $\alpha$, and the Szeg\H{o} limit
\[
\lim_{\alpha\to\infty}\frac{1}{c(\alpha)}\,\operatorname{tr}\big(\psi(T^\alpha_\sigma)\big)
\;=\;\int_X \psi(\sigma(x))\,d\nu(x)
\]
holds.
\end{thm}

\begin{proof}
Since $\sigma\in L^1(X,\nu)$ and $d\nu_\alpha=c(\alpha)\,d\nu$, we have $T^\alpha_\sigma\in\mathcal{S}_1$ and $\sum_j|\lambda^{(\alpha)}_j|<\infty$ for the eigenvalues $\{\lambda^{(\alpha)}_j\}$. The local bound \eqref{eq:psi-linear-local} and continuity of $\psi$ on the compact spectral range imply
\[
\sum_j|\psi(\lambda^{(\alpha)}_j)|\;\le\;C\sum_j|\lambda^{(\alpha)}_j|+M\;<\;\infty,
\]
hence $\psi(T^\alpha_\sigma)\in\mathcal{S}_1$.

Choose a monotone increasing family $\{\eta_R\}_R\subset L^1(X,\nu)$ with $0\le \eta_R\le 1$ and $\eta_R\uparrow 1$ pointwise. For each fixed $R$, $T^\alpha_{\eta_R}$ is bounded and, by our Szeg\H{o} theorem above applied to $\eta=\eta_R$ and $\psi$, 
\begin{equation}\label{eq:weighted-limit-R}
\lim_{\alpha\to\infty}\frac{1}{c(\alpha)}\,\operatorname{tr}\!\big(T^\alpha_{\eta_R}\,\psi(T^\alpha_\sigma)\big)
\;=\;\int_X \eta_R(x)\,\psi(\sigma(x))\,d\nu(x).
\end{equation}
(Here we use that $T^\alpha_{\eta_R}$ and $T^\alpha_\sigma$ commute for this application; if such a commuting family $\{\eta_R\}$ is not available, skip \eqref{eq:weighted-limit-R} and argue directly by dominated convergence below.)

Because $0\le T^\alpha_{\eta_R}\le I$ and $T^\alpha_{\eta_R}\to I$ strongly as $R\to\infty$, for $A=\psi(T^\alpha_\sigma)\in\mathcal{S}_1$ we have
\[
\operatorname{tr}(T^\alpha_{\eta_R}A)\;\longrightarrow\;\operatorname{tr}(A)\qquad(R\to\infty),
\]
by the usual trace–dominated convergence (equivalently, by the diagonal integral representation against the frame and monotone convergence). Hence, for each fixed $\alpha$,
\[
\lim_{R\to\infty}\operatorname{tr}\!\big(T^\alpha_{\eta_R}\,\psi(T^\alpha_\sigma)\big)
=\operatorname{tr}\!\big(\psi(T^\alpha_\sigma)\big).
\]
On the right-hand side of \eqref{eq:weighted-limit-R}, monotone convergence yields
\[
\lim_{R\to\infty}\int_X \eta_R\,\psi(\sigma)\,d\nu=\int_X \psi(\sigma)\,d\nu.
\]
Combining these limits with \eqref{eq:weighted-limit-R} gives the claim.
\end{proof}

\begin{remark}
\label{rem:etaR-commuting}
The step \eqref{eq:weighted-limit-R} requires, for each fixed $R$, that $T^\alpha_{\eta_R}$ commute with $T^\alpha_\sigma$. In models where a natural exhaustion $\{\eta_R\}$ has this property (e.g. symbols adapted to the geometry of the frame), the proof above applies verbatim. In full generality one can bypass \eqref{eq:weighted-limit-R} by writing
\[
\frac{1}{c(\alpha)}\,\operatorname{tr}\!\big(\psi(T^\alpha_\sigma)\big)
=\int_X \underbrace{\langle \psi(T^\alpha_\sigma)k_x^\alpha,k_x^\alpha\rangle}_{:=\,\widetilde{\psi(\sigma)}_\alpha(x)}\,d\nu(x),
\]
and using $\big|\widetilde{\psi(\sigma)}_\alpha(x)\big|\le C\,\widetilde{|\sigma|}_\alpha(x)$ (from \eqref{eq:psi-linear-local}) together with $\widetilde{|\sigma|}_\alpha\to|\sigma|$ in $L^1$ by Assumption~(II) to conclude by dominated convergence.
\end{remark}
\end{comment}

%%%%%%%%%%%%%%%%%%%%%%%%%%%%%%%%%%%%%%%%%%%%%%%%%%%%

\begin{prop}\label{diagonal} Assume that $\sigma\in L^1(X,\nu)$ is real-valued and bounded.
Let $\psi$ be continuous on a compact interval containing
$\operatorname{ess\,ran}(\sigma)$ and $0$. Then
\[
\lim_{\alpha\to\infty}
\frac{1}{c(\alpha)}
\operatorname{Tr}
\left(
T_\sigma^\alpha\psi(T_\sigma^\alpha)
\right)
=
\int_X \sigma(x)\psi(\sigma(x))\,d\nu(x).
\]
\end{prop}

\begin{proof}
Set
\[
\varphi(t)=t\psi(t).
\]
Then $\varphi(0)=0$. Fix $\varepsilon>0$. By the Weierstrass
approximation theorem, choose a polynomial $p$ such that
\[
\|\psi-p\|_\infty<\varepsilon
\]
on the chosen compact interval, and set
\[
q(t)=t p(t).
\]
Then
\[
|\varphi(t)-q(t)|\leq \varepsilon |t|.
\]

Since $q(0)=0$, write $q=q_1-q_2$ with $q_1$ and $q_2$
convex and satisfying $q_1(0)=q_2(0)=0$. For example, if
\[
M=\inf q'',
\]
then one may take
\[
q_1=q-\frac{M}{2}t^2,
\qquad
q_2=-\frac{M}{2}t^2
\]
when $M<0$, and $q_1=q$, $q_2=0$ when $M\geq0$.

Applying Proposition~\ref{BL-classical} to $q_1$ and $q_2$, and using
$\widetilde{\sigma}^{\alpha}\to\sigma$ in $L^1(X,\nu)$, gives
\[
\lim_{\alpha\to\infty}
\frac{1}{c(\alpha)}
\operatorname{Tr}q(T_\sigma^\alpha)
=
\int_X q(\sigma(x))\,d\nu(x).
\]

Moreover,
\[
\left|\varphi(t)-q(t)\right|\leq\varepsilon |t|,
\]
so functional calculus gives
\[
\left|
\frac{1}{c(\alpha)}
\operatorname{Tr}
\bigl(\varphi-q\bigr)(T_\sigma^\alpha)
\right|
\leq
\varepsilon\,
\frac{\operatorname{Tr}|T_\sigma^\alpha|}{c(\alpha)}
\leq
\varepsilon\int_X|\sigma|\,d\nu.
\]
The same estimate gives
\[
\left|
\int_X(\varphi-q)(\sigma(x))\,d\nu(x)
\right|
\leq
\varepsilon\int_X|\sigma|\,d\nu.
\]
Letting $\varepsilon\to0$ proves the result.
\end{proof}

\subsection{Eigenvalue distribution: the locally compact case} All the assumptions in this section are still in place. Let $\{\lambda_j^\alpha\}$ denote the eigenvalues of $T_\sigma^\alpha$.

\begin{cor}\label{eigenvalue} Suppose all the assumptions from Theorem~\ref{thm:general_szego_lc} hold. For every $t\in\mathbb R$ such that
\[
\nu(\sigma=t)=0,
\]
and every non-negative $\eta\in L^1(X,\nu)\cap L^2(X,\nu)$,
\[
\lim_{\alpha\to\infty}
\frac{1}{c(\alpha)}
\sum_{\lambda_j^\alpha>t}
\langle T_\eta^\alpha e_j^\alpha,e_j^\alpha\rangle
=
\int_{\{\sigma>t\}} \eta(x)\, d\nu(x),
\]
where $\{e_j^\alpha\}$ is an orthonormal eigenbasis of $T_\sigma^\alpha$.
\end{cor}

\begin{proof} Choose continuous functions $\psi_{1,n},\psi_{2,n}$ such that
\[
0\leq\psi_{1,n}
\leq\mathbf 1_{(t,\infty)}
\leq\psi_{2,n}\leq1,
\]
\[
\psi_{1,n}(x)=\psi_{2,n}(x)
=\mathbf 1_{(t,\infty)}(x)
\qquad\text{for }x\notin
[t-1/n,t+1/n],
\]
and
\[
\psi_{j,n}(x)\longrightarrow
\mathbf 1_{(t,\infty)}(x)
\qquad\text{for every }x\neq t.
\]
It is clear that such functions satisfy the assumptions of Theorem~\ref{thm:general_szego_lc} (ii). Furthermore, note that for each $n$
\begin{align*}
    \text{tr}\left(T_{\eta}^{\alpha}\psi_{1,n}\left(T_{\sigma}^{\alpha}\right)\right)\leq\text{tr}\left(T_{\eta}^{\alpha} 1_{(t,\infty)}\left(T_{\sigma}^{\alpha}\right)\right)\leq\text{tr}\left(T_{\eta}^{\alpha}\psi_{2,n}\left(T_{\sigma}^{\alpha}\right)\right).
\end{align*}
Consider the lower bound. First, note that 
\begin{align*}
    \frac{1}{c(\alpha)}\text{tr}\left(T_{\eta}^{\alpha}\psi_{1,n}\left(T_{\sigma}^{\alpha}\right)\right)-\int_{X}\eta\psi_{1,n}(\sigma)d\nu\leq\frac{1}{c(\alpha)}\text{tr}\left(T_{\eta}^{\alpha} 1_{(t,\infty)}\left(T_{\sigma}^{\alpha}\right)\right)-\int_{X}\eta\psi_{1,n}(\sigma)d\nu.
\end{align*}
Taking the limit as $\alpha\to\infty$ and applying Theorem~\ref{thm:general_szego_lc} (ii), we obtain 
\begin{align*}
    0\leq\liminf_{\alpha\to\infty}\frac{1}{c(\alpha)}\text{tr}\left(T_{\eta}^{\alpha} 1_{(t,\infty)}\left(T_{\sigma}^{\alpha}\right)\right)-\int_{X}\eta\psi_{1,n}(\sigma)d\nu.
\end{align*}
Hence, 
\begin{align*}
    \int_{X}\eta\psi_{1,n}(\sigma)d\nu\leq\liminf_{\alpha\to\infty}\frac{1}{c(\alpha)}\text{tr}\left(T_{\eta}^{\alpha} 1_{(t,\infty)}\left(T_{\sigma}^{\alpha}\right)\right).
\end{align*}
Since $\nu(\sigma=t)=0$,
\[
\psi_{j,n}(\sigma(x))
\longrightarrow
\mathbf 1_{\{\sigma>t\}}(x)
\qquad\text{for $\nu$-almost every }x.
\]
Moreover,
\[
\left|\eta(x)\psi_{j,n}(\sigma(x))\right|
\leq |\eta(x)|.
\]
Therefore, the Dominated Convergence Theorem gives
\[
\int_X\eta(x)\psi_{j,n}(\sigma(x))\,d\nu(x)
\longrightarrow
\int_{\{\sigma>t\}}\eta(x)\,d\nu(x).
\]

For the upper bound, we similarly have
\[
\frac{1}{c(\alpha)}
\operatorname{Tr}\!\left(
T_\eta^\alpha
\mathbf{1}_{(t,\infty)}(T_\sigma^\alpha)
\right)
\le
\frac{1}{c(\alpha)}
\operatorname{Tr}\!\left(
T_\eta^\alpha
\psi_{2,n}(T_\sigma^\alpha)
\right).
\]
Therefore, by Theorem~\ref{thm:general_szego_lc},
\[
\limsup_{\alpha\to\infty}
\frac{1}{c(\alpha)}
\operatorname{Tr}\!\left(
T_\eta^\alpha
\mathbf{1}_{(t,\infty)}(T_\sigma^\alpha)
\right)
\le
\int_X \eta(x)\psi_{2,n}(\sigma(x))\,d\nu(x).
\]
Letting \(n\to\infty\) and applying the dominated convergence theorem gives
\[
\limsup_{\alpha\to\infty}
\frac{1}{c(\alpha)}
\operatorname{Tr}\!\left(
T_\eta^\alpha
\mathbf{1}_{(t,\infty)}(T_\sigma^\alpha)
\right)
\le
\int_{\{\sigma>t\}}\eta(x)\,d\nu(x).
\]
Together with the lower bound, this proves the claimed limit.

\end{proof}
\begin{cor}
\label{cor:eig-count-sigma}
Let $\sigma\in L^1(X,\nu)$ be real-valued. Assume, in addition, $\sigma$ is bounded or non-negative. 
Let $\{\lambda_j^{(\alpha)}\}$ denote the eigenvalues of $T^\alpha_\sigma$ in
decreasing order.  If $t>0$ satisfies $\nu(\{\sigma=t\})=0$, then
\[
\lim_{\alpha\to\infty}
\frac{1}{c(\alpha)}\,\#\{j:\lambda_j^{(\alpha)}>t\}
=
\nu(\sigma>t).
\]
\end{cor}

\begin{proof} 
First suppose that $\sigma\geq0$. We use the following consequence of
Proposition~\ref{BLweighted}: for every $\varphi\in C_0([0,\infty))$,
\begin{equation}\label{prop2.2}
\frac{1}{c(\alpha)}
\operatorname{Tr}\!\left(
T_\sigma^\alpha\varphi(T_\sigma^\alpha)
\right)
\longrightarrow
\int_X\sigma(x)\varphi(\sigma(x))\,d\nu(x).
\end{equation}

Indeed, by Lemma~\ref{lem:diff_convex}(ii), choose convex functions
$q_1,q_2\in C_0([0,\infty))$ such that
\[
\|\varphi-(q_1-q_2)\|_\infty<\varepsilon.
\]
Applying Proposition~\ref{BLweighted} with $\eta=\sigma$ and $\psi=q_j$ gives
\[
\int_X\sigma q_j(\widetilde{\sigma}^{\,\alpha})\,d\nu
\leq
\frac{1}{c(\alpha)}
\operatorname{Tr}\!\left(
T_\sigma^\alpha q_j(T_\sigma^\alpha)
\right)
\leq
\int_X\widetilde{\sigma}^{\,\alpha}q_j(\sigma)\,d\nu.
\]
Since $\widetilde{\sigma}^{\,\alpha}\to\sigma$ in $L^1(X,\nu)$,
the two bounds converge to
\[
\int_X\sigma(x)q_j(\sigma(x))\,d\nu(x).
\]
The uniform approximation then gives \eqref{prop2.2}, because
\[
\left|
\frac{1}{c(\alpha)}
\operatorname{Tr}\!\left(
T_\sigma^\alpha
\bigl(\varphi-(q_1-q_2)\bigr)(T_\sigma^\alpha)
\right)
\right|
\leq
\varepsilon\int_X\sigma\,d\nu.
\]

Choose continuous functions $g_{1,n},g_{2,n}$ such that
\[
0\leq g_{1,n}
\leq \mathbf 1_{(t,\infty)}
\leq g_{2,n}\leq1,
\]
and
\[
g_{j,n}(x)\longrightarrow \mathbf 1_{(t,\infty)}(x)
\qquad\text{for every }x\neq t.
\]
For sufficiently large $n$, define
\[
\varphi_{j,n}(x)
=
\begin{cases}
\dfrac{g_{j,n}(x)}{x},&x>0,\\[2mm]
0,&x=0.
\end{cases}
\]
Since $t>0$, the functions $\varphi_{j,n}$ belong to
$C_0([0,\infty))$. For every eigenvalue $\lambda_j^\alpha\geq0$,
\[
\lambda_j^\alpha\varphi_{1,n}(\lambda_j^\alpha)
\leq
\mathbf 1_{(t,\infty)}(\lambda_j^\alpha)
\leq
\lambda_j^\alpha\varphi_{2,n}(\lambda_j^\alpha).
\]
Therefore,
\[
\frac{1}{c(\alpha)}
\sum_j\lambda_j^\alpha
\varphi_{1,n}(\lambda_j^\alpha)
\leq
\frac{1}{c(\alpha)}
\#\{j:\lambda_j^\alpha>t\}
\leq
\frac{1}{c(\alpha)}
\sum_j\lambda_j^\alpha
\varphi_{2,n}(\lambda_j^\alpha).
\]

For all sufficiently large $n$, we may choose $g_{j,n}$ to vanish on
$[0,t/2]$. Hence, for $j=1,2$,
\[
0\leq \sigma(x)\phi_{j,n}(\sigma(x))
 =g_{j,n}(\sigma(x))
 \leq \mathbf{1}_{\{\sigma>t/2\}}(x)
 \leq \frac{2}{t}\sigma(x).
\]

Using \eqref{prop2.2} and then letting $n\to\infty$, the Dominated Convergence Theorem gives
\[
\lim_{\alpha\to\infty}
\frac{1}{c(\alpha)}
\#\{j:\lambda_j^\alpha>t\}
=
\nu(\{\sigma>t\}),
\]
because $\nu(\{\sigma=t\})=0$.

If \(\sigma\) is bounded but not necessarily non-negative, use
Proposition~\ref{diagonal} instead of the previous corollary. For sufficiently large \(n\),
choose the functions \(g_{j,n}\) so that they vanish in a neighborhood
of \(0\), and define
\[
\phi_{j,n}(x)=
\begin{cases}
g_{j,n}(x)/x, & x>0,\\
0, & x\le 0.
\end{cases}
\]
Then \(\phi_{j,n}\) is continuous on every compact interval containing
the essential range of \(\sigma\) and \(0\). Proposition 5.3 gives
\[
\frac{1}{c(\alpha)}
\sum_j \lambda_j^{(\alpha)}
\phi_{j,n}(\lambda_j^{(\alpha)})
\longrightarrow
\int_X \sigma(x)\phi_{j,n}(\sigma(x))\,d\nu(x).
\]
Since \(t>0\), the same cutoff inequalities and dominated convergence
argument yield
\[
\frac{1}{c(\alpha)}
\#\{j:\lambda_j^{(\alpha)}>t\}
\longrightarrow
\nu(\{\sigma>t\}).
\]
\end{proof}  

% Let $\psi_n$ be continuous functions satisfying
% \[
% 0\le \psi_n\le1,
% \quad
% \psi_n(x)\to \mathbf 1_{(t,\infty)}(x)
% \]
% pointwise.

% Applying the weighted Szeg\H{o} theorem gives
% \[
% \lim_{\alpha\to\infty}
% \frac{1}{c(\alpha)}
% \operatorname{tr}\!\big(
% T_\eta^\alpha \psi_n(T_\sigma^\alpha)
% \big)
% =
% \int_X \eta \psi_n(\sigma)\, d\nu.
% \]

% Letting $n\to\infty$ and using dominated convergence
% together with the assumption $\nu(\sigma=t)=0$
% yields
% \[
% \lim_{\alpha\to\infty}
% \frac{1}{c(\alpha)}
% \operatorname{tr}\!\big(
% T_\eta^\alpha \mathbf 1_{(t,\infty)}(T_\sigma^\alpha)
% \big)
% =
% \int_{\{\sigma>t\}} \eta\, d\nu.
% \]

% Since
% \[
% \operatorname{tr}\!\big(
% T_\eta^\alpha \mathbf 1_{(t,\infty)}(T_\sigma^\alpha)
% \big)
% =
% \sum_{\lambda_j^\alpha>t}
% \langle T_\eta^\alpha e_j^\alpha,e_j^\alpha\rangle,
% \]
% the result follows.
% \end{proof}



%%%%%%%%%%%%%%%%%%%%%%%%%%%%%%%%%%%%%%%%%%%%%%%%%%%%%
\begin{comment}

\subsection{Bergman Space}

Let $\alpha>-1$, and let $\mathcal{H}_{\alpha}=A_{\alpha}^{2}(\mathbb{D})=\{f\in L^{2}(\mathbb{D},(\alpha+1)(1-|z|^{2})^{\alpha}dA(z)):\text{f is holomorphic on }\mathbb{D}\}$. It is known that $A_{\alpha}^{2}(\mathbb{D})$ is a RKHS under the $L^{2}(\mathbb{D},(\alpha+1)(1-|z|^{2})^{\alpha}dA(z))$ inner product, with reproducing kernel
\begin{align*}
    \ip{K_{z}^{\alpha}}{K_{w}^{\alpha}}_{\alpha}=K^{\alpha}(z,w)=\frac{1}{\left(1-z\overline{w}\right)^{\alpha+2}}.
\end{align*}
We have that $\norm{K_{z}^{\alpha}}_{\alpha}^{2}=K^{\alpha}(z,z)=\left(1-|z|^{2}\right)^{-(\alpha+2)}$. Using this, we obtain the normalized reproducing kernel:
\begin{align*}
    |\ip{k_{z}^{\alpha}}{k_{w}^{\alpha}}_{\alpha}|
    =&\frac{1}{\norm{K_{z}^{\alpha}}_{\alpha}\norm{K_{w}^{\alpha}}_{\alpha}}|\ip{K_{z}^{\alpha}}{K_{w}^{\alpha}}_{\alpha}|\\
    =&\left(1-|z|^{2}\right)^{1+\alpha/2}\left(1-|w|^{2}\right)^{1+\alpha/2}\frac{1}{\left|1-z\overline{w}\right|^{\alpha+2}}\\
    =&\left(\frac{\left(1-|z|^{2}\right)\left(1-|w|^{2}\right)}{\left|1-z\overline{w}\right|^2}\right)^{1+\alpha/2}
    =\left(1-\left|\frac{w-z}{1-z\overline{w}}\right|^{2}\right)^{2+\alpha}.
\end{align*}
Thus, 
\begin{align*}
    |\ip{k_{z}^{\alpha}}{k_{w}^{\alpha}}_{\alpha}|^{2}=\left(1-\left|\frac{w-z}{1-z\overline{w}}\right|^{2}\right)^{4+2\alpha}.
\end{align*}
%Now, using a result of Zhu (page 52 of \cite{Z90}), we have that  
%\begin{align*}\left|\ip{k_{z}^{\alpha}}{k_{w}^{\alpha}}_{\alpha}\right|^{2}=\left(\frac{\left(1-|z|^{2}\right)\left(1-|w|^{2}\right)}{\left|1-z\overline{w}\right|^2}\right)^{2+\alpha}=\left(1-\left|\frac{w-z}{1-z\overline{w}}\right|^{2}\right)^{2+\alpha}.\end{align*}
Now, note that the associated measure under this normalized reproducing kernel is 
\begin{align*}
    d\nu_{\alpha}(z)=\norm{K_{z}^{\alpha}}_{\alpha}^{2}(\alpha+1)(1-|z|^{2})^{\alpha}dA(z)=(\alpha+1)\frac{1}{(1-|z|^{2})^{2}}dA(z).
\end{align*}
Hence, $\nu_{\alpha}$ satisfies Assumption (I), by letting $c(\alpha)=\alpha+1$ and $d\nu=(1-|z|^{2})^{-2}dA(z)$. 

We now check Assumption (II). Let $z\in\mathbb{D}$ and $R>0$. Let $\beta(z,w)$ denote the Bergman metric on $\mathbb{D}$, and let $\mathbb{D}(z,R):=\{w\in \mathbb{D}: \beta(w,z)<R\}$ denote the Bergman ball centered at $z$, with radius $R>0$. Using the Mobius transformations and the invariance of the measure and Bergman metric under this group action, it is sufficient to consider the case where $z=0$. Additionally, it is well-known (see e.g. \cite{Z90}) that the Bergman ball $\mathbb{D}(0,R)$ is a Euclidean ball with center at zero and radius $s=\tanh{(R)}$. Thus, we have 
\begin{align*}
    \int_{\mathbb{D}(0,R)^{c}}\left|\ip{k_{0}^{\alpha}}{k_{w}^{\alpha}}_{\alpha}\right|^{2}d\nu_{\alpha}(w)=&\int_{B(0,s)^{c}}\left(1-|w|^{2}\right)^{2\alpha+4}(\alpha+1)\left(1-|w|^{2}\right)^{-2}dA(w)\\
    =&\frac{1}{\pi}(\alpha+1)\int_{0}^{2\pi}\int_{s}^{1}(1-r^{2})^{2\alpha+2}rdrd\theta\\
    =&(\alpha+1)\int_{0}^{1-s^{2}}u^{2\alpha+2}du\\
    =&\frac{(\alpha+1)}{2\alpha+3}(1-s^{2})^{2\alpha+3}\to 0,
\end{align*}
as $\alpha\to\infty$, since $s<1$. Hence, Assumption (II) is satisfied. Therefore, we obtain Theorem~\ref{Sze2} from the introduction as a direct consequence of Corollary~\ref{Sze_Berg}

\end{comment}

%%%%%%%%%%%%%%%%%%%%%%%%%%%%%%%%%%%%%%%%%%%%%%%%%%

\begin{comment}

\subsection{Weighted Bergman Spaces on $\mathbb{D}$}

Let $\alpha>-1$ and consider the weighted Bergman space
\[
\mathcal{H}_{\alpha}
=
A^{2}_{\alpha}(\mathbb{D})
=
\left\{
f\in L^{2}\!\left(\mathbb{D},(\alpha+1)(1-|z|^{2})^{\alpha} dA(z)\right)
:\;
f \text{ holomorphic on } \mathbb{D}
\right\}.
\]
It is well known that $A^{2}_{\alpha}(\mathbb{D})$ is a reproducing kernel Hilbert space with respect to the inner product induced by the measure
\[
(\alpha+1)(1-|z|^{2})^{\alpha} dA(z),
\]
and its reproducing kernel is given by
\[
K^{\alpha}(z,w)
=
\frac{1}{(1-z\overline{w})^{\alpha+2}}.
\]
In particular,
\[
\|K^{\alpha}_{z}\|^{2}_{\alpha}
=
K^{\alpha}(z,z)
=
(1-|z|^{2})^{-(\alpha+2)}.
\]

\medskip

\noindent
The normalized reproducing kernels are therefore
\[
k^{\alpha}_{z}
=
\frac{K^{\alpha}_{z}}{\|K^{\alpha}_{z}\|_{\alpha}},
\]
and a straightforward computation yields
\begin{align*}
\left|
\langle k^{\alpha}_{z},k^{\alpha}_{w}\rangle_{\alpha}
\right|
&=
\frac{|K^{\alpha}(z,w)|}
{\|K^{\alpha}_{z}\|_{\alpha}\|K^{\alpha}_{w}\|_{\alpha}}  \\
&=
\left(
\frac{(1-|z|^{2})(1-|w|^{2})}
{|1-z\overline{w}|^{2}}
\right)^{\frac{\alpha+2}{2}}.
\end{align*}
Using the identity
\[
1-\left|\frac{w-z}{1-z\overline{w}}\right|^{2}
=
\frac{(1-|z|^{2})(1-|w|^{2})}
{|1-z\overline{w}|^{2}},
\]
we obtain
\[
\left|
\langle k^{\alpha}_{z},k^{\alpha}_{w}\rangle_{\alpha}
\right|^{2}
=
\left(
1-
\left|\frac{w-z}{1-z\overline{w}}\right|^{2}
\right)^{\alpha+2}.
\]

\medskip

\noindent
The associated measure appearing in our abstract framework is
\begin{align*}
d\nu_{\alpha}(z)
&=
\|K^{\alpha}_{z}\|^{2}_{\alpha}
(\alpha+1)(1-|z|^{2})^{\alpha} dA(z)  \\
&=
(\alpha+1)\frac{1}{(1-|z|^{2})^{2}}\, dA(z).
\end{align*}
Hence Assumption~(I) is satisfied with
\[
c(\alpha)=\alpha+1,
\qquad
d\nu(z)=\frac{1}{(1-|z|^{2})^{2}}\, dA(z),
\]
which is independent of $\alpha$.

\medskip

\noindent
We now verify Assumption~(II). Let $\beta(z,w)$ denote the Bergman metric on $\mathbb{D}$ and let
\[
\mathbb{D}(z,R)
=
\{w\in\mathbb{D}:\beta(w,z)<R\}
\]
be the Bergman ball of radius $R>0$. By invariance of both the Bergman metric and the measure $d\nu_{\alpha}$ under Möbius automorphisms of $\mathbb{D}$, it suffices to consider the case $z=0$.

It is classical that
\[
\mathbb{D}(0,R)
=
B(0,s),
\qquad
s=\tanh R,
\]
a Euclidean disk centered at the origin. Since
\[
\left|
\langle k^{\alpha}_{0},k^{\alpha}_{w}\rangle_{\alpha}
\right|^{2}
=
(1-|w|^{2})^{\alpha+2},
\]
we compute
\begin{align*}
\int_{\mathbb{D}(0,R)^{c}}
\left|
\langle k^{\alpha}_{0},k^{\alpha}_{w}\rangle_{\alpha}
\right|^{2}
\, d\nu_{\alpha}(w)
&=
(\alpha+1)
\int_{B(0,s)^{c}}
(1-|w|^{2})^{\alpha}
\, dA(w) \\
&=
(\alpha+1)
\int_{s}^{1}
(1-r^{2})^{\alpha} r\, dr \\
&=
\frac{\alpha+1}{2}
\int_{0}^{1-s^{2}}
u^{\alpha}\, du \\
&=
\frac{\alpha+1}{2(\alpha+1)}
(1-s^{2})^{\alpha+1} \\
&=
\frac{1}{2}(1-s^{2})^{\alpha+1}.
\end{align*}
Since $s<1$, we have $(1-s^{2})^{\alpha+1}\to 0$ as $\alpha\to\infty$. Hence Assumption~(II) holds.

\medskip

\noindent
We therefore obtain Theorem~\ref{Sze2} from the introduction as a direct consequence of Theorem~\ref{thm:general_szego_lc}, yielding the Szeg\H{o}-type limit theorem for weighted Bergman spaces on $\mathbb{D}$.

\end{comment}


%%%%%%%%%%%%%%%%%%%%%%%%%%%%%%%%%%%%%%%%%%%%%%%%%%%%%

\subsection{Weighted Bergman Spaces on the Unit Ball $\mathbb{B}^n$}

Let $n\ge 1$ and let $\mathbb{B}^n\subset\mathbb{C}^n$ denote the unit ball.  
For $\alpha>-1$ we consider the weighted Bergman space
\[
A^2_\alpha(\mathbb{B}^n)
=
\left\{
f\in\mathcal{O}(\mathbb{B}^n):
\;
\|f\|^2_{A^2_\alpha}
:=
c_{\alpha,n}
\int_{\mathbb{B}^n} |f(z)|^2 (1-|z|^2)^\alpha\, dV(z)
<\infty
\right\},
\]
where $dV$ is Lebesgue measure on $\mathbb{C}^n$ and
$c_{\alpha,n}$ is the normalizing constant
\[
c_{\alpha,n}
=
\frac{\Gamma(n+\alpha+1)}{\pi^n\,\Gamma(\alpha+1)}
\]
so that $\|1\|_{A^2_\alpha}=1$.
It is well known that $A^2_\alpha(\mathbb{B}^n)$ is a reproducing kernel Hilbert space with reproducing kernel
\begin{equation}
K^\alpha(z,w)
=
\frac{1}{(1-\langle z,w\rangle)^{n+1+\alpha}},
\qquad z,w\in\mathbb{B}^n.
\label{eq:ball-kernel}
\end{equation}
In particular,
\[
\|K^\alpha_z\|^2_{A^2_\alpha}
=
K^\alpha(z,z)
=
(1-|z|^2)^{-(n+1+\alpha)}.
\]

\medskip

\noindent
Thus the normalized reproducing kernels are
\[
k^\alpha_z(w)
=
\frac{K^\alpha(w,z)}{\sqrt{K^\alpha(z,z)}},
\qquad 
z,w\in\mathbb{B}^n.
\]
A direct computation using \eqref{eq:ball-kernel} yields
\[
\left|
\langle k^\alpha_z, k^\alpha_w\rangle_{A^2_\alpha}
\right|
=
\left(
\frac{(1-|z|^2)(1-|w|^2)}
{|1-\langle z,w\rangle|^2}
\right)^{\frac{n+1+\alpha}{2}}.
\]
Recall that the Bergman distance $\beta(z,w)$ satisfies
\[
\cosh^2\!\left(\frac{\beta(z,w)}{2}\right)
=
\frac{|1-\langle z,w\rangle|^2}{(1-|z|^2)(1-|w|^2)},
\]
so that
\[
\frac{(1-|z|^2)(1-|w|^2)}{|1-\langle z,w\rangle|^2}
=
\operatorname{sech}^2\!\left(\frac{\beta(z,w)}{2}\right).
\]
Therefore
\begin{equation}
\left|
\langle k^\alpha_z, k^\alpha_w\rangle_{A^2_\alpha}
\right|^2
=
\operatorname{sech}^{\,2(n+1+\alpha)}\!\left(\frac{\beta(z,w)}{2}\right).
\label{eq:kernel-decay}
\end{equation}

\medskip

\noindent
The associated measures in the abstract framework are
\[
d\nu_\alpha(z)
=
\|K^\alpha_z\|^2_{A^2_\alpha}\,
c_{\alpha,n}(1-|z|^2)^\alpha\, dV(z)
=
c_{\alpha,n}\,(1-|z|^2)^{-n-1}\, dV(z).
\]
Since $c_{\alpha,n}$ depends on $\alpha$ while 
$(1-|z|^2)^{-n-1} dV(z)$ does not,  
Assumption~\textup{(I)} holds with
\[
c(\alpha)=c_{\alpha,n},
\qquad 
d\nu(z)=(1-|z|^2)^{-n-1}dV(z),
\]
independent of $\alpha$.

\medskip

\noindent
We now verify Assumption~\textup{(II)}.
Let $\beta(\cdot,\cdot)$ denote the Bergman distance on $\mathbb{B}^n$, and let
\[
\mathbb{B}^n(z,R)
=
\{w\in\mathbb{B}^n : \beta(w,z)<R\}
\]
be the Bergman ball of radius $R>0$.
Using the invariance of both the Bergman metric and the measure $d\nu_\alpha$
under the automorphism group $\operatorname{Aut}(\mathbb{B}^n)$, it suffices to treat the case $z=0$.
In this case $\mathbb{B}^n(0,R)$ is the Euclidean ball of radius $s=\tanh (R/2)$.

For $z=0$, the inner product formula \eqref{eq:kernel-decay} reduces to
\[
\left|\langle k^\alpha_0, k^\alpha_w\rangle\right|^2
=
(1-|w|^2)^{n+1+\alpha}.
\]
Hence
\begin{align*}
&\int_{\mathbb{B}^n(0,R)^c}
\left|\left\langle k_0^\alpha,k_w^\alpha\right\rangle\right|^2
\,d\nu_\alpha(w) \\
&\qquad=
c_{\alpha,n}\int_{|w|\geq s}(1-|w|^2)^\alpha\,dV(w) \\
&\qquad=
c_{\alpha,n}\sigma_{2n-1}
\int_s^1 (1-r^2)^\alpha r^{2n-1}\,dr \\
&\qquad=
\frac{c_{\alpha,n}\sigma_{2n-1}}{2}
\int_0^{1-s^2}u^\alpha(1-u)^{n-1}\,du \\
&\qquad\leq
\frac{c_{\alpha,n}\sigma_{2n-1}}{2(\alpha+1)}
(1-s^2)^{\alpha+1}.
\end{align*}
Since
\[
c_{\alpha,n}
=
\frac{\Gamma(n+\alpha+1)}{\pi^n\Gamma(\alpha+1)}
=
\frac{1}{\pi^n}\prod_{j=1}^n(\alpha+j)
=
O(\alpha^n),
\]
and $0<1-s^2<1$, the final expression tends to $0$ as
$\alpha\to\infty$. 
Thus, Assumption~\textup{(II)} follows.

\medskip

\noindent
We therefore obtain Theorem~\ref{Sze2} for weighted Bergman spaces on the unit ball $\mathbb{B}^n$
as a direct consequence of the general Szeg\H{o}-type limit theorems, Theorem~\ref{thm:general_szego_lc} and
Proposition~\ref{diagonal}.

%%%%%%%%%%%%%%%%%%%%%%%%%%%%%%%%%%%%%%%%%%%%%%%%%%%%
\begin{comment}

\subsection{Locally Compact Abelian Groups}
The case when $G$ is locally compact and non-compact follows similarly to the compact case. Here, we let $\Gamma_{N}\subset\hat{G}$ be a F\o lner sequence, and define $\mathcal{H}_{N}=\{f\in L^{2}(G):\text{supp}(\hat{f})\subset\Gamma_{N}\}$. It is clear that $\mathcal{H}_{N}$ is a RKHS with reproducing kernel
\begin{align*}
    K_{N}(x,y)=\hat{\chi}_{\Gamma_{N}}(x-y),
\end{align*}
where $\chi_{\Gamma_{N}}$ is the characteristic function on $\Gamma_{N}$, and $\hat{\cdot}$ denotes the Fourier transform. Note that 
\begin{align*}
    \norm{K_{x}^{N}}^{2}=K_{N}(x,x)=\hat{\chi}_{\Gamma_{N}}(0)=\hat{\mu}(\Gamma_{N}).
\end{align*}
We consider the topological space $G$ with Radon measure $\mu_{N}=\hat{\mu}(\Gamma_{N})\mu$. 
Assumption $(I)$ is trivially satisfied. Using the same normalization procedure as in the compact case, we obtain the normalized reproducing kernel as 
\begin{align*}
    \hat{\mu}(\Gamma_{N})\left|\ip{k_{x}^{N}}{k_{y}^{N}}\right|^{2}=\frac{1}{\hat{\mu}(\Gamma_{N})}\left|\hat{\chi}_{\Gamma_{N}}(x-y)\right|^{2}.
\end{align*}
In this case Lemma 3.4 of \cite{coifman1973operators} still applies, giving us Theorem 1.3. \\
An immediate corollary to this is the classical results of Landau and Widom \cite{landau1980eigenvalue}, which is concerned with localization operators. In our setting these operators correspond to Toeplitz operators whose symbol is the indicator of a set, acting on the Paley-Wiener space of band-limited functions. In the view of Theorem 5.5, we obtain the following partial analog to the main result of \cite{landau1980eigenvalue}: 
\begin{cor}
    Let $PW_{\alpha}=\{f\in L^{2}(\R):\text{supp}(\hat{f})\subset[-\alpha,\alpha]\}$, let $\sigma\in L^{1}(\R,dx)\cap L^{\infty}(\R,dx)$ be non-negative, and let $T_{\sigma}^{\alpha}:PW_{\alpha}(\R)\to PW_{\alpha}(\R)$ be the corresponding Toeplitz operator. Then, for $\lambda>0$ such that $|\{\sigma=\lambda\}|=0$
    \begin{align*}
        N_{\alpha}(\lambda)=\frac{\alpha}{\pi}|\{\sigma>\lambda\}|,
    \end{align*}
    as $\alpha\to\infty$. 
\end{cor}
However, it should be noted that the results obtained in \cite{landau1980eigenvalue} are slightly stronger, as they include the second-order term. \\
Finally, using Proposition 3.6 we obtain the following generalization of Wiener's lemma which was recently obtained in \cite{jaming2026wiener}:
\begin{cor}[Jaming, Kellay, and Perez (2026)]
    Let $\mu$ be a complex-valued Borel measure on a locally compact Abelian group $G$ such that $|\mu|$ is a finite measure. Let $\{F_{N}\}_{N=1}^{\infty}$ be a F\o lner sequence in the dual group $\hat{G}$. Then, 
    \begin{align*}
        \mu(\{x\})=\lim_{N\to\infty}\frac{1}{\hat{m}(F_{N})}\int_{F_{N}}\xi(x)\hat{\mu}(\xi)d\hat{m}(\xi),
    \end{align*}
    where $\hat{m}$ is the Haar measure on $\hat{G}$. 
\end{cor}

\end{comment}
%%%%%%%%%%%%%%%%%%%%%%%%%%%%%%%%%%%%%%%%%%%%%%
\subsection{Locally Compact Abelian Groups}

Let $G$ be a locally compact Abelian group with Haar measure $m$ and dual group $\widehat{G}$ equipped with Haar measure $\widehat{m}$.  
Let $\{F_N\}_{N\ge1}\subset \widehat{G}$ be a F\o lner sequence, i.e.
\[
\frac{\widehat{m}((F_N+\gamma)\triangle F_N)}{\widehat{m}(F_N)}
\longrightarrow 0
\qquad \text{for every } \gamma\in\widehat{G}.
\]

Define the Paley–Wiener type spaces
\[
\mathcal{H}_N
=
\{f\in L^2(G): \operatorname{supp}(\widehat{f})\subset F_N\}.
\]
Each $\mathcal{H}_N$ is a reproducing kernel Hilbert space with kernel
\[
K_N(x,y)
=
\widehat{\chi_{F_N}}(x-y),
\]
where $\chi_{F_N}$ denotes the characteristic function of $F_N$ and $\widehat{\cdot}$ denotes the inverse Fourier transform on $\widehat{G}$.

Since
\[
K_N(x,x)
=
\widehat{\chi_{F_N}}(0)
=
\widehat{m}(F_N),
\]
the normalization constant is
\[
c(N)=\widehat{m}(F_N).
\]

Thus Assumption (I),
\[
d\nu_N = c(N)\, dm,
\]
is satisfied.

Moreover, the normalized kernel satisfies
\[
c(N)\, |\langle k_x^N,k_y^N\rangle|^2
=
\frac{1}{\widehat{m}(F_N)}
\big|\widehat{\chi_{F_N}}(x-y)\big|^2,
\]
which forms an approximate identity on $G$ by the F\o lner property.
In particular, the Berezin transforms converge in measure for
$L^1$-symbols.

Therefore, Theorem~\ref{thm:general_szego_lc} and Proposition~\ref{diagonal}
apply in this setting and yields Theorem 1.3.

\medskip


% \medskip

% \subsubsection*{Application to the classical Paley--Wiener space}

% Consider $G=\mathbb{R}$ and
% \[
% PW_\alpha
% =
% \{f\in L^2(\mathbb{R}): \operatorname{supp}(\widehat{f})
% \subset [-\alpha,\alpha]\}.
% \]
% Then $c(\alpha)=2\alpha$ and the Toeplitz operators
% \[
% T_\sigma^\alpha : PW_\alpha \to PW_\alpha
% \]
% correspond to classical time–frequency localization operators.

% As an immediate consequence of
% Theorem~\ref{thm:general_szego_lc}, we obtain:

% \begin{cor}
% Let $\sigma\in L^1(\mathbb{R})\cap L^\infty(\mathbb{R})$
% be non-negative.
% Let $N_\alpha(\lambda)$ denote the number of eigenvalues of
% $T_\sigma^\alpha$ exceeding $\lambda>0$.
% If $|\{\sigma=\lambda\}|=0$, then
% \[
% \lim_{\alpha\to\infty}
% \frac{1}{2\alpha}
% N_\alpha(\lambda)
% =
% |\{x\in\mathbb{R}:\sigma(x)>\lambda\}|.
% \]
% Equivalently,
% \[
% N_\alpha(\lambda)
% \sim
% \frac{\alpha}{\pi}
% \,|\{\sigma>\lambda\}|
% \qquad (\alpha\to\infty).
% \]
% \end{cor}

% This recovers the first-order term in the classical eigenvalue
% asymptotics of
% Landau and Widom~\cite{landau1980eigenvalue}.
% Their result includes a second-order term,
% which lies beyond the scope of the present abstract approach.

\medskip

% \subsubsection*{Wiener-type formula}

% As another application, combining
% Theorem~\ref{thm:general_szego_lc}
% with Proposition~3.6 yields a generalization of
% Wiener’s lemma recently obtained in
% \cite{jaming2026wiener}:

% \begin{cor}[Jaming--Kellay--Perez (2026)]
% Let $\mu$ be a complex Borel measure on $G$
% with $|\mu|(G)<\infty$.
% Let $\{F_N\}$ be a F\o lner sequence in $\widehat{G}$.
% Then for every $x\in G$,
% \[
% \mu(\{x\})
% =
% \lim_{N\to\infty}
% \frac{1}{\widehat{m}(F_N)}
% \int_{F_N}
% \xi(x)\,\widehat{\mu}(\xi)
% \, d\widehat{m}(\xi).
% \]
% \end{cor}

% Thus, in the locally compact Abelian setting,
% the abstract Szeg\H{o} theorem recovers both:

% \begin{itemize}
% \item classical eigenvalue asymptotics for localization operators,
% \item and Fourier-analytic limit formulas of Wiener type.
% \end{itemize}
%%%%%%%%%%%%%%%%%%%%%%%%%%%%%%%%%%%%%%%%%%%%%%%%%
% Below is a polished **LaTeX draft** of a full **“Further Applications”** section, including the **Feichtinger–Nowak application**, written to match the tone and mathematical style of your paper.  
% You can paste it directly into your manuscript (with minor renumbering depending on where you insert it).

% ***

% ```latex
\section{Further Applications}

In this final section we illustrate the scope of our abstract Szeg\H{o} limit theorems by presenting several further applications.  
These include classical results of Landau--Widom for Paley--Wiener spaces, 
a recent Wiener-type lemma of Jaming--Kellay--Pérez on locally compact Abelian groups, 
a Szeg\H{o}-type theorem for Gabor (time--frequency) localization operators originally due to Feichtinger and Nowak, and finally, the necessary density condition for sampling in the Bargmann-Fock space, originally due to Lindholm.  
All of these follow naturally from Theorem~\ref{thm:general_szego_lc} under the appropriate geometric choices of $X$, $\nu$, $c(\alpha)$, and the normalized kernel family $\{k^\alpha_x\}$.


\subsection{Gabor Localization Operators}

Fix any window
\[
g\in L^2(\mathbb R),
\qquad
\|g\|_{L^2(\mathbb R)}=1.
\]
We use the time--frequency shifts
\[
\pi(q,p)f(t)=e^{2\pi i pt}f(t-q).
\]
For $h>0$, let
\[
D_hf(t)=h^{-1/4}f(t/\sqrt h),
\qquad
g_h=D_hg,
\]
and define the semiclassical shifts by
\[
\pi_h(x,\xi)f(t)
=
e^{2\pi i \xi t/h}f(t-x).
\]
Then
\[
\pi_h(x,\xi)
=
D_h\pi(x/\sqrt h,\xi/\sqrt h)D_h^{-1}.
\]

The short-time Fourier transform associated with $g_h$ is
\[
V_{g_h}f(x,\xi)
=
\langle f,\pi_h(x,\xi)g_h\rangle_{L^2(\mathbb R)}.
\]
Let $\mathcal H_h$ denote the range of $V_{g_h}$, equipped with the
inner product for which $V_{g_h}$ is unitary.

The normalized reproducing kernels are
\[
k_z^h(w)
=
\left\langle
\pi_h(z)g_h,\pi_h(w)g_h
\right\rangle_{L^2(\mathbb R)},
\qquad z,w\in\mathbb R^2.
\]
Define
\[
\gamma(\zeta)
=
\left|
\left\langle g,\pi(\zeta)g\right\rangle_{L^2(\mathbb R)}
\right|^2.
\]
Then
\[
\left|
\left\langle k_z^h,k_w^h\right\rangle_{\mathcal H_h}
\right|^2
=
\gamma\left(\frac{w-z}{\sqrt h}\right).
\]
Moreover, the Moyal identity gives
\[
\int_{\mathbb R^2}\gamma(\zeta)\,d\zeta=1.
\]
Consequently, for every $R>0$,
\[
\int_{\{|w-z|>R\}}
\left|
\left\langle k_z^h,k_w^h\right\rangle_{\mathcal H_h}
\right|^2
h^{-1}\,dw
=
\int_{\{|\zeta|>R/\sqrt h\}}\gamma(\zeta)\,d\zeta
\longrightarrow 0
\]
as $h\to0$. Thus the frame measures are
\[
d\nu_h(z)=h^{-1}\,dz,
\qquad
c(h)=h^{-1}.
\]

For a non-negative symbol $\sigma$, define
\[
T_h(\sigma)f
=
\int_{\mathbb R^2}
\sigma(z)
\langle f,k_z^h\rangle_{\mathcal H_h}
k_z^h\,d\nu_h(z).
\]

\begin{thm}\label{FeiNow}
Let
\[
\sigma\in L^1(\mathbb R^2)\cap L^2(\mathbb R^2),
\qquad
\sigma\geq0.
\]
Let $\psi:[0,\infty)\to\mathbb R$ be continuous and suppose that
\[
\lim_{t\to\infty}\frac{\psi(t)}{t}
\]
exists and is finite. Then
\[
\lim_{h\to0}
h\,\operatorname{Tr}
\bigl(T_h(\sigma)\psi(T_h(\sigma))\bigr)
=
\int_{\mathbb R^2}
\sigma(z)\psi(\sigma(z))\,dz.
\]
\end{thm}

This result is related to the theorem of Feichtinger and Nowak
\cite{FN01}. They use the same arbitrary fixed normalized
window $g$, the shifts $\pi(q,p)$, and the dilated symbols
\[
b_R(q,p)=b\left(\frac{q}{R},\frac{p}{R}\right).
\]
Their result is
\[
\lim_{R\to\infty}
\frac{1}{R^2}
\operatorname{Tr}
\bigl(T_{b_R}\psi(T_{b_R})\bigr)
=
\int_{\mathbb R^2}
b(z)\psi(b(z))\,dz,
\]
for $b\in L^1(\mathbb R^2)$ with $0\leq b\leq1$ and
$\psi\in C([0,1])$.

Taking
\[
R=h^{-1/2}
\]
gives $R^{-2}=h$. Under the unitary dilation $D_h$, the operator
$T_h(\sigma)$ is unitarily equivalent to the Feichtinger--Nowak
operator with symbol
\[
\sigma_R(z)=\sigma(z/R).
\]
Hence the preceding semiclassical formula is precisely the
Feichtinger--Nowak limit written in semiclassical variables.

The present theorem extends their result by allowing symbols in
$L^1\cap L^2$ which need not be bounded, and by allowing continuous
functions $\psi$ on $[0,\infty)$ with linear growth at infinity.


\subsection{Paley--Wiener Spaces on $\mathbb{R}^{n}$ and Dilates of Convex Bodies}

We next consider the multidimensional Paley--Wiener setting treated in the classical
work of Landau--Widom~\cite{landau1980eigenvalue}.
Let $\Omega\subset\mathbb{R}^{n}$ be a bounded convex set with nonempty interior 
(and, for convenience, symmetric about the origin).
For $\alpha>0$ define the frequency set
\[
\alpha\Omega = \{ \alpha\xi : \xi\in\Omega\},
\]
and the corresponding Paley--Wiener space
\[
PW_{\alpha\Omega}
=
\left\{
f\in L^{2}(\mathbb{R}^{n}) : \operatorname{supp}(\widehat f)\subset \alpha\Omega
\right\}.
\]
This is a reproducing kernel Hilbert space with reproducing kernel
\[
K_{\alpha}(x,y)
=
(2\pi)^{-n}
\int_{\alpha\Omega} e^{i\xi\cdot(x-y)}\,d\xi,
\qquad x,y\in\mathbb{R}^{n}.
\]
Evaluating at the diagonal gives
\[
K_{\alpha}(x,x)
=
(2\pi)^{-n}\,|\alpha\Omega|
=
(2\pi)^{-n}\,\alpha^{n}|\Omega|.
\]
Consequently, the diagonal measures in our abstract framework take the form
\[
d\nu_{\alpha}(x)
=
\|K_{\alpha,x}\|^{2}\,dx
=
c(\alpha)\,dx,
\qquad
c(\alpha)
=
(2\pi)^{-n}\,\alpha^{n}|\Omega|,
\]
so that Assumption~\textup{(I)} is satisfied.

\medskip

The normalized kernels are
\[
k_{\alpha,x}(y)
=
\frac{K_{\alpha}(y,x)}{\sqrt{K_{\alpha}(x,x)}}
=
\frac{
(2\pi)^{-n}
\int_{\alpha\Omega} e^{i\xi\cdot(y-x)}\,d\xi}
{\sqrt{(2\pi)^{-n}\alpha^{n}|\Omega|}}.
\]

A well-known asymptotic estimate (see, for example, 
\cite[Section~3]{landau1980eigenvalue}) shows that the quantities
$|\langle k_{\alpha,x},k_{\alpha,y}\rangle|^{2}$ form an approximate identity:
for every fixed
$R>0$,
\[
\int_{\|x-y\|>R}
\bigl|\langle k_{\alpha,x},k_{\alpha,y}\rangle\bigr|^2
\,d\nu_\alpha(y)
\longrightarrow0
\qquad(\alpha\to\infty).
\]
Thus, Assumption~\textup{(II)} also holds.

\medskip

Applying Theorem~\ref{thm:general_szego_lc} to $PW_{\alpha\Omega}$ gives the following 
multidimensional generalization of the one--dimensional Paley--Wiener case:

\begin{thm}
Let $\Omega\subset\mathbb{R}^n$ be a bounded convex set with
nonempty interior, and let $PW_{\alpha\Omega}$ be as above. Let
\[
\sigma\in L^1(\mathbb{R}^n)\cap L^2(\mathbb{R}^n)
\]
be non-negative. Let $\psi\in C([0,\infty))$ satisfy $\psi(0)=0$,
and suppose that
\[
\phi(t):=
\begin{cases}
\dfrac{\psi(t)}{t},&t>0,\\[2mm]
\displaystyle\lim_{s\downarrow0}\frac{\psi(s)}{s},&t=0
\end{cases}
\]
is continuous on $[0,\infty)$. Assume also that
\[
\lim_{t\to\infty}\frac{\psi(t)}{t}
\]
exists. Then
\[
\lim_{\alpha\to\infty}
\frac{1}{c(\alpha)}
\operatorname{Tr}\psi\bigl(T_\alpha(\sigma)\bigr)
=
\int_{\mathbb{R}^n}\psi(\sigma(x))\,dx,
\]
where
\[
c(\alpha)=(2\pi)^{-n}\alpha^n|\Omega|.
\]
\end{thm}

\begin{proof} Set $\eta=\sigma$ in Theorem~\ref{thm:general_szego_lc} and apply it to the function
$\phi(x)=\psi(x)/x$. Since $\phi$ is bounded and
continuous, it satisfies the growth condition in Theorem~\ref{thm:general_szego_lc}.
Moreover,
\[
T_\alpha(\sigma)\phi\bigl(T_\alpha(\sigma)\bigr)
=
\psi\bigl(T_\alpha(\sigma)\bigr)
\]
by functional calculus, and
\[
\sigma(x)\phi(\sigma(x))=\psi(\sigma(x)).
\]
The claimed formula follows.
    
\end{proof}
    


This recovers the leading term in the Landau--Widom asymptotic formula 
for time--band limiting operators~\cite{landau1980eigenvalue}.
Their celebrated second-order correction term, 
involving the geometry of the boundary $\partial\Omega$,
lies beyond the reach of our current abstract method.
Nonetheless, the first-order term follows immediately from the
general Szeg\H{o} theorem in a conceptually simple and unified way.

\begin{cor}
\label{cor:slepian}
Let $\Omega=[-1,1]\subset\mathbb{R}$ and let $PW_{\alpha\Omega}$ be the Paley--Wiener
space of $L^{2}$-functions whose Fourier transforms are supported in $[-\alpha,\alpha]$.
Let $\sigma=\chi_{[-1,1]}$ be the indicator function of the time interval $[-1,1]$,
and let $T_{\alpha}(\sigma)$ be the associated time--frequency localization operator
\[
T_{\alpha}(\sigma) f 
=
\int_{-1}^{1} f(x)\, K_{\alpha}(x,\cdot) \, dx,
\]
where $K_{\alpha}$ is the reproducing kernel of $PW_{\alpha\Omega}$.
Let $\lambda^{(\alpha)}_{1}\ge \lambda^{(\alpha)}_{2}\ge \dots$ 
denote the eigenvalues of $T_{\alpha}(\sigma)$.

Then for every $t\in(0,1)$ such that $|\{x:\sigma(x)=t\}|=0$, 
\[
\lim_{\alpha\to\infty}
\frac{\pi}{\alpha}\,
\#\{j:\lambda^{(\alpha)}_{j}>t\}
=
|\{x\in[-1,1]:\sigma(x)>t\}|
=
2.
\]

Equivalently,
\[
\#\{j:\lambda^{(\alpha)}_{j}>t\}
=
\frac{\alpha}{\pi}\left|\{x\in[-1,1]:\sigma(x)>t\}\right|
+o(\alpha)
=
\frac{2\alpha}{\pi}+o(\alpha).
\]
\end{cor}

In particular, the number of eigenvalues of $T_{\alpha}(\sigma)$
that lie near $1$ grows like $\frac{2\alpha}{\pi}$.
These eigenfunctions are the classical prolate spheroidal wave functions (Slepian functions),
and this recovers the sharp Slepian--Pollak--Landau concentration law.


\subsection{A Wiener-Type Limit on LCA Groups}
Let $G$ be a locally compact Abelian group with dual $\widehat{G}$, and let $\{F_N\}$ be a Følner sequence in $\widehat{G}$.  
Section~5.5 shows that each Paley--Wiener type space 
\[
H_N = \{f\in L^{2}(G): \mathrm{supp}(\widehat{f})\subset F_N\}
\]
fits into our abstract setting, with $c(N)=\widehat{m}(F_N)$ and normalized kernels forming an approximate identity. For \(x\neq y\),
\[
\langle k_x^N,k_y^N\rangle
=
\frac{1}{\widehat m(F_N)}
\int_{F_N}\xi(x-y)\,d\widehat m(\xi).
\]
Since \(x-y\neq0\), the function
\(\xi\mapsto\xi(x-y)\) is a nontrivial character of \(\widehat G\).
The Følner property therefore implies
\[
\langle k_x^N,k_y^N\rangle\longrightarrow0.
\]
Thus, Assumption (III) holds. Therefore, Proposition~\ref{lkc} applies and 
we recover the recent Wiener-type lemma of Jaming--Kellay--Pérez \cite{jaming2026wiener}:
\[
\mu(\{x\})
=
\lim_{N\to\infty}
\frac{1}{\widehat{m}(F_N)}
\int_{F_N} \xi(x)\,\widehat{\mu}(\xi)\, d\widehat{m}(\xi),
\]
for any finite complex Borel measure $\mu$.


\subsection{Sampling in the Bargmann--Fock Space}
\label{sec:FockSampling}

In this subsection we show how the classical Lindholm--Landau density condition \cite{lindholm2001sampling}
for sampling in the Bargmann--Fock space follows directly from
our locally compact abstract Szeg\H{o} theorem, through its eigenvalue distribution consequence (Corollary~\ref{cor:eig-count-sigma}).  
\begin{comment}
The key points are:  
(i) Toeplitz operators $T^\alpha_{\chi_{B(a,R)}}$ have spectrum independent of
$a\in\C$;  
(ii) sequences must be appropriately scaled when passing to $\Falpha$;
(iii) all counting takes place in fixed-radius balls $B(a,R)$, with $R$ not scaled.
\end{comment}

Recall that
\[
\Falpha
=\Bigl\{ f\in\mathcal{O}(\C): 
\|f\|_{\Falpha}^{2}
=\frac{\alpha}{\pi}
\int_{\C} |f(z)|^{2}e^{-\alpha|z|^{2}}\,dm(z)<\infty\Bigr\},
\]
The reproducing kernel is
\[
K_\alpha(z,w)=e^{\alpha z\overline w},
\]
and the normalized kernels satisfy
\[
k_z^{(\alpha)}(w)
=e^{\alpha w\overline z-\alpha|z|^2/2},
\qquad
\left|\langle k_z^{(\alpha)},k_w^{(\alpha)}\rangle\right|^2
=e^{-\alpha|z-w|^2}.
\]
Moreover,
\[
d\nu_\alpha(z)=\frac{\alpha}{\pi}\,dm(z),
\qquad
c(\alpha)=\frac{\alpha}{\pi}.
\]
Consequently,
\[
\int_{|w-z|>R}
\left|\langle k_z^{(\alpha)},k_w^{(\alpha)}\rangle\right|^2
\,d\nu_\alpha(w)
=e^{-\alpha R^2}\longrightarrow0.
\]
Thus, the Fock spaces satisfy Assumptions (I) and (II). 

The dilation operator
\[
(D_\alpha f)(z)=f(\sqrt{\alpha}\,z).
\]
is a unitary map $\calF^{2}\to\Falpha$, and a straightforward computation shows the validity of the following lemma.

\begin{lem}
\label{lem:scalingFock}
Let $\Lambda\subset\C$, and define the scaled set
\[
\Lambda_\alpha := \frac{1}{\sqrt{\alpha}}\,\Lambda.
\]
Then $\Lambda$ is sampling (resp.\ interpolating) in $\calF^{2}$ with constants
$A,B$ if and only if $\Lambda_\alpha$ is sampling (resp.\ interpolating) in
$\Falpha$ with the same constants $A,B$.
\end{lem}




For $a\in\C$ and $R>0$ set
\[
T^\alpha_{a,R}:=P_\alpha\,M_{\chi_{B(a,R)}}\,P_\alpha,
\] where $P_\alpha$ is the Bergman projection onto $\Falpha$. 
As in the Fock model, the Weyl translation
\[
(\Weyl_b f)(z)=e^{\alpha(z\overline{b}-\frac12|b|^{2})}f(z-b)
\]
is unitary on $\Falpha$, satisfies $\Weyl_b P_\alpha=P_\alpha\Weyl_b$, and

\[W_a M_{\chi_{B(0,R)}}W_a^*
=M_{\chi_{B(a,R)}},
\qquad
T_{a,R}^\alpha
=W_aT_{0,R}^\alpha W_a^*.\]
Hence

\begin{lem}[Translation covariance]
\label{lem:WeylBall}
For every $a\in\C$, $R>0$, and $\alpha>0$,
\[
T^\alpha_{a,R}
=\Weyl_{a}\,T^\alpha_{0,R}\,\Weyl_{a}^{*}.
\]
Thus, $T^\alpha_{a,R}$ has the same eigenvalues for all $a\in\C$.
\end{lem}

Let $N_\alpha^{(R)}(t)$ denote the number of eigenvalues of $T^\alpha_{a,R}$
exceeding $t\in(0,1)$; by Lemma~\ref{lem:WeylBall} this is independent of $a$.

Applying Corollary~\ref{cor:eig-count-sigma} to the bounded non-negative symbol
\[
\sigma=\chi_{B(a,R)},
\]
we obtain, for every $t\in(0,1)$,
\[
\frac{1}{c(\alpha)}
N_\alpha^{(R)}(t)
\longrightarrow
|\{z\in\mathbb C:\chi_{B(a,R)}(z)>t\}|
=
|B(a,R)|
=
\pi R^2.
\]
Since $c(\alpha)=\alpha/\pi$, this is equivalent to
\begin{equation}\label{eq:FockEigenvalueCounting}
\frac{N_\alpha^{(R)}(t)}{\alpha R^2}
\longrightarrow1.
\end{equation}

%%%%%%%%%%%%%%%%%%%%%%%%%


Assume that $\Lambda$ is a separated sampling set for $F^2$.
Thus, there exists $A>0$ such that
\[
A\|f\|_{F^2}^2
\leq
\sum_{\lambda\in\Lambda}
|\langle f,k_\lambda^{(1)}\rangle_{F^2}|^2,
\qquad
f\in F^2.
\]
By the dilation lemma, $\Lambda_\alpha
=\alpha^{-1/2}\Lambda$ is a sampling set for $F_\alpha^2$ with the
same lower sampling constant:
\[
A\|f\|_{F_\alpha^2}^2
\leq
\sum_{\lambda\in\Lambda_\alpha}
|\langle f,k_\lambda^{(\alpha)}\rangle_{F_\alpha^2}|^2.
\]

Fix $R>0$ and choose $0<\delta<R$. Also choose
\[
0<\rho_0<
\min\left\{\frac12\operatorname{sep}(\Lambda),R\right\},
\qquad
\rho_\alpha=\frac{\rho_0}{\sqrt{\alpha}}.
\]
Since
\[
\operatorname{sep}(\Lambda_\alpha)
=
\frac{\operatorname{sep}(\Lambda)}{\sqrt{\alpha}},
\]
the balls
\[
\{B(\lambda,\rho_\alpha):
\lambda\in\Lambda_\alpha\}
\]
are pairwise disjoint.

A standard Bargmann--Fock mean-value estimate (see \cite{Z12} Lemma 2.32) gives, for
$w\notin B(a,R)$,
\[
|\langle f,k_w^{(\alpha)}\rangle|^2
\leq
\frac{\alpha}{\pi}
\left(1-e^{-\alpha\rho_\alpha^2}\right)^{-1}
\int_{B(w,\rho_\alpha)}
|\langle f,k_u^{(\alpha)}\rangle|^2\,du.
\]
Since $\alpha\rho_\alpha^2=\rho_0^2$, summing over
$\lambda\in\Lambda_\alpha\setminus B(a,R)$ gives
\[
\begin{aligned}
\sum_{\lambda\in\Lambda_\alpha\setminus B(a,R)}
|\langle f,k_\lambda^{(\alpha)}\rangle|^2
&\leq
\left(1-e^{-\rho_0^2}\right)^{-1} \\
&\quad\times
\int_{\mathbb C\setminus B(a,R-\rho_\alpha)}
|\langle f,k_u^{(\alpha)}\rangle|^2
\frac{\alpha}{\pi}\,du.
\end{aligned}
\tag{6.2}
\]

Choose $t\in(0,1)$ sufficiently close to $1$ that
\[
\left(1-e^{-\rho_0^2}\right)^{-1}(1-t)
\leq \frac{A}{2}.
\]
Let $E_\alpha(a,R-\delta;t)$ be the span of the eigenvectors of
$T_{a,R-\delta}^{(\alpha)}$ whose eigenvalues are strictly greater
than $t$. Hence
\[
\dim E_\alpha(a,R-\delta;t)
=
N_\alpha^{(R-\delta)}(t).
\]

For all sufficiently large $\alpha$, we have $\rho_\alpha<\delta$.
If $f\in E_\alpha(a,R-\delta;t)$ and
$\|f\|_{F_\alpha^2}=1$, then
\[
\int_{B(a,R-\delta)}
|\langle f,k_u^{(\alpha)}\rangle|^2
\frac{\alpha}{\pi}\,du
=
\langle T_{a,R-\delta}^{(\alpha)}f,f\rangle
>t.
\]
Using the Parseval frame identity,
\[
\int_{\mathbb C}
|\langle f,k_u^{(\alpha)}\rangle|^2
\frac{\alpha}{\pi}\,du
=
1,
\]
and the inclusion
\[
B(a,R-\delta)\subset B(a,R-\rho_\alpha),
\]
we obtain
\[
\int_{\mathbb C\setminus B(a,R-\rho_\alpha)}
|\langle f,k_u^{(\alpha)}\rangle|^2
\frac{\alpha}{\pi}\,du
\leq 1-t.
\]
Therefore, by \emph{(6.2)},
\[
\begin{aligned}
\sum_{\lambda\in\Lambda_\alpha\cap B(a,R)}
|\langle f,k_\lambda^{(\alpha)}\rangle|^2
&\geq
A-
\left(1-e^{-\rho_0^2}\right)^{-1}(1-t) \\
&\geq \frac{A}{2}.
\end{aligned}
\]

Consequently, the restriction map
\[
\mathcal R_{\alpha,a}:
E_\alpha(a,R-\delta;t)
\longrightarrow
\ell^2\bigl(\Lambda_\alpha\cap B(a,R)\bigr),
\]
defined by
\[
\mathcal R_{\alpha,a}f
=
\bigl(
\langle f,k_\lambda^{(\alpha)}\rangle
\bigr)_{\lambda\in\Lambda_\alpha\cap B(a,R)},
\]
is injective. Hence
\[
N_\alpha^{(R-\delta)}(t)
\leq
\#\bigl(\Lambda_\alpha\cap B(a,R)\bigr),
\qquad
a\in\mathbb C.
\]

By Corollary~\ref{cor:eig-count-sigma}, applied to the symbol
$\chi_{B(a,R-\delta)}$, we have
\[
\frac{N_\alpha^{(R-\delta)}(t)}{c(\alpha)}
\longrightarrow
|B(a,R-\delta)|
=
\pi(R-\delta)^2,
\]
where
\[
c(\alpha)=\frac{\alpha}{\pi}.
\]
Thus
\[
\frac{N_\alpha^{(R-\delta)}(t)}
{\alpha(R-\delta)^2}
\longrightarrow 1.
\]
Taking the infimum over $a$ gives
\[
\liminf_{\alpha\to\infty}
\frac{
\inf_{a\in\mathbb C}
\#\bigl(\Lambda_\alpha\cap B(a,R)\bigr)
}{
\alpha(R-\delta)^2
}
\geq 1.
\]

Since $\Lambda_\alpha=\alpha^{-1/2}\Lambda$,
\[
\#\bigl(\Lambda\cap B(z,R\sqrt{\alpha})\bigr)
=
\#\bigl(\Lambda_\alpha\cap B(z/\sqrt{\alpha},R)\bigr).
\]
Therefore,
\[
\liminf_{\alpha\to\infty}
\frac{
\inf_{z\in\mathbb C}
\#\bigl(\Lambda\cap B(z,R\sqrt{\alpha})\bigr)
}{
\pi(R\sqrt{\alpha})^2
}
\geq
\frac{(R-\delta)^2}{\pi R^2}.
\]
Since $\alpha$ ranges over all positive real numbers, $R\sqrt{\alpha}$
ranges over all sufficiently large radii. Hence
\[
D^-(\Lambda)
\geq
\frac{(R-\delta)^2}{\pi R^2}.
\]
Finally, letting $\delta\downarrow0$ yields
\[
D^-(\Lambda)\geq\frac1\pi.
\]


%%%%%%%%%%%%%%%%%%%%%%%%%%%%%%%%%%%
\begin{comment}

Let $\Lambda$ be separated with $\text{sep}(\Lambda)>0$, and fix $0<r_{0}<\min(\text{sep}(\Lambda), R)$ and let $r_{\alpha}=r_{0}/\sqrt{\alpha}$.
A standard Bargmann-Fock space mean--value argument (see \cite{Z12} Lemma 2.32) gives, for $w\notin B(a,R)$ and all
$f\in\Falpha$,
\begin{align*}
\label{eq:mvFock}
|\langle f,k^{(\alpha)}_{w}\rangle|^{2}
\ \le\ 
\frac{\alpha}{\pi}\left(\left(1-e^{-\pi \alpha r_{\alpha}^{2}}\right)\right)^{-1}\int_{B(w,r_{\alpha})}
|\langle f,k^{(\alpha)}_{u}\rangle|^{2}\,du.\tag{6.2}
\end{align*}
Summing \eqref{eq:mvFock} over
$\lambda\in\Lambda_\alpha\setminus B(a,R)$ (disjointness of the balls
$B(\lambda,r)$ for $\lambda\in\Lambda_\alpha$) yields
\begin{align*}
\label{eq:sumOutside}
\sum_{\lambda\in\Lambda_\alpha\setminus B(a,R)}
|\langle f,k^{(\alpha)}_{\lambda}\rangle|^{2}
\ \le\ 
\left(\left(1-e^{-\pi \alpha r_{\alpha}^{2}}\right)\right)^{-1}
\int_{\C\setminus B(a,R-r_{\alpha})}
|\langle f,k^{(\alpha)}_{u}\rangle|^{2}\frac{\alpha}{\pi}\,du.\tag{6.3}
\end{align*}
%By translating variables, \eqref{eq:sumOutside} holds for balls $B(z,R)$ centered at any $z\in\C$.


Assume $\Lambda$ is sampling in $\calF^{2}$; by
Lemma~\ref{lem:scalingFock}, $\Lambda_\alpha$ is sampling in $\Falpha$ with the
same constants:
\[
A\|f\|_{\Falpha}^{2}
\le \sum_{\lambda\in\Lambda_\alpha}
|\langle f,k^{(\alpha)}_{\lambda}\rangle|^{2}.
\]
Fix $t\in(0,1)$ such that 
\begin{align*}
    (1-t)\left(\left(1-e^{-\pi r_{0}^{2}}\right)\right)^{-1}\leq \frac{A}{2}.
\end{align*}
Let $E_\alpha(a,R-r;t)$ denote the span of eigenvectors of $T^\alpha_{a,R-r}$ with
eigenvalues $\ge t$. Clearly,
\[
\dim E_\alpha(a,R-r;t)=N_\alpha^{(R-r)}(t).
\]

If $f\in E_\alpha(a,R-r;t)$ with $\|f\|=1$, then
\[
\int_{B(a,R-r)}|\langle f,k^{(\alpha)}_u\rangle|^{2}\frac{\alpha}{\pi}\,du\ge t,
\]
and using \eqref{eq:sumOutside} with $r$ small enough gives
\[
\sum_{\lambda\in\Lambda_\alpha\cap B(a,R)}
|\langle f,k^{(\alpha)}_{\lambda}\rangle|^{2}
\ge A-\left(\left(1-e^{-\pi r_{0}^{2}}\right)\right)^{-1}(1-t)\ge \frac{A}{2}.
\]
The restriction map
\[
\mathcal R_{\alpha,a,R}:
E_\alpha(a,R-r;t)
\longrightarrow
\ell^2\bigl(\Lambda_\alpha\cap B(a,R)\bigr),
\]
defined by
\[
\mathcal R_{\alpha,a,R}f
=
\bigl(\langle f,k_\lambda^{(\alpha)}\rangle\bigr)
_{\lambda\in\Lambda_\alpha\cap B(a,R)},
\]
is injective. Consequently,
\[
N_\alpha^{(R-r)}(t)
\leq
\#\bigl(\Lambda_\alpha\cap B(a,R)\bigr).
\]

Taking $\inf_{a}$ and using~\eqref{eq:FockEigenvalueCounting},
\[
\lim_{\alpha\to\infty}\frac{\inf_{a\in\C}\#\bigl(\Lambda_\alpha\cap B(a,R)\bigr)}{\alpha(R-r)^2}
\ \ge\ 1.
\]

Undoing the scaling, using
$\Lambda_\alpha=\alpha^{-1/2}\Lambda$,
\[
\#\bigl(\Lambda\cap B(z,R\sqrt{\alpha})\bigr)
=\#\bigl(\Lambda_\alpha\cap B(z/\sqrt{\alpha},R)\bigr)
,
\]
and hence
\[
\inf_{z\in\C}
\frac{\#(\Lambda\cap B(z,(R\sqrt{\alpha})))}{\pi (R\sqrt{\alpha})^{2}}
=\inf_{z\in\C}\frac{\#\bigl(\Lambda_\alpha\cap B(z/\sqrt{\alpha},R)\bigr)}{ \alpha(R-r)^2}\frac{\alpha(R-r)^2}{\pi(R\sqrt{\alpha})^{2}}.
\]
Letting $\alpha\to\infty$ gives
\[
D^{-}(\Lambda)\ \ge\ \frac{1}{\pi}\frac{(R-r)^2}{R^2}.
\]
Finally, taking $r\to 0$ we get 
\[
D^{-}(\Lambda)\ \ge\ \frac{1}{\pi}.
\]

\end{comment}
% Let $\Lambda$ be separated with $\sep(\Lambda)>0$, and fix $0<r<\min(\sep(\Lambda), R)$.
% A standard Bargmann-Fock space mean--value argument gives, for $w\notin B(a,R)$ and all
% $f\in\Falpha$,
% \begin{equation}
% \label{eq:mvFock}
% |\langle f,k^{(\alpha)}_{w}\rangle|^{2}
% \ \le\ 
% C_{r}\int_{B(w,r)}
% |\langle f,k^{(\alpha)}_{u}\rangle|^{2}\,du.
% \end{equation}
% Summing \eqref{eq:mvFock} over
% $\lambda\in\Lambda_\alpha\setminus B(a,R)$ (disjointness of the balls
% $B(\lambda,r)$ for $\lambda\in\Lambda_\alpha$) yields
% \begin{equation}
% \label{eq:sumOutside}
% \sum_{\lambda\in\Lambda_\alpha\setminus B(a,R)}
% |\langle f,k^{(\alpha)}_{\lambda}\rangle|^{2}
% \ \le\ 
% C_{r}
% \int_{\C\setminus B(a,R-r)}
% |\langle f,k^{(\alpha)}_{u}\rangle|^{2}\,du.
% \end{equation}
% %By translating variables, \eqref{eq:sumOutside} holds for balls $B(z,R)$ centered at any $z\in\C$.


% Assume $\Lambda$ is sampling in $\calF^{2}$; by
% Lemma~\ref{lem:scalingFock}, $\Lambda_\alpha$ is sampling in $\Falpha$ with the
% same constants:
% \[
% A\|f\|_{\Falpha}^{2}
% \le \sum_{\lambda\in\Lambda_\alpha}
% |\langle f,k^{(\alpha)}_{\lambda}\rangle|^{2}.
% \]
% Fix $t\in(0,1)$.  
% Let $E_\alpha(a,R-r;t)$ denote the span of eigenvectors of $T^\alpha_{a,R-r}$ with
% eigenvalues $\ge t$. Clearly,
% \[
% \dim E_\alpha(a,R-r;t)=N_\alpha^{(R-r)}(t).
% \]

% If $f\in E_\alpha(a,R-r;t)$ with $\|f\|=1$, then
% \[
% \int_{B(a,R-r)}|\langle f,k^{(\alpha)}_u\rangle|^{2}\,du\ge t,
% \]
% and using \eqref{eq:sumOutside} with $r$ small enough gives
% \[
% \sum_{\lambda\in\Lambda_\alpha\cap B(a,R)}
% |\langle f,k^{(\alpha)}_{\lambda}\rangle|^{2}
% \ge A-C_r(1-t)\ge \frac{A}{2}.
% \]
% The restriction map $R:E_\alpha(a,R-r;t)\to \ell^{2}(\Lambda_\alpha\cap B(a, R-r))$, given by 
% \[f\mapsto(\langle f,k^{(\alpha)}_\lambda \rangle)_{\lambda\in\Lambda\cap B(a, R-r)}\] is injective, and therefore
% \[
% N_\alpha^{(R-r)}(t)
% \le \#\bigl(\Lambda_\alpha\cap B(a,R)\bigr)
% \qquad\forall a\in\C.
% \]
% Taking $\inf_{a}$ and using \eqref{eq:FockEigenvalueCounting},
% \[
% \lim_{\alpha\to\infty}\frac{\inf_{a\in\C}\#\bigl(\Lambda_\alpha\cap B(a,R)\bigr)}{\alpha(R-r)^2}
% \ \ge\ 1.
% \]

% Undoing the scaling, using
% $\Lambda_\alpha=\alpha^{-1/2}\Lambda$,
% \[
% \#\bigl(\Lambda\cap B(z,R\sqrt{\alpha})\bigr)
% =\#\bigl(\Lambda_\alpha\cap B(z/\sqrt{\alpha},R)\bigr)
% ,
% \]
% % \begin{align*}
% %     \Lambda_{\alpha}=\phi_{a}\left(\frac{1}{\alpha+1}\phi_{a}^{-1}(\Lambda)\right)
% % \end{align*}
% and hence
% \[
% \inf_{z\in\C}
% \frac{\#(\Lambda\cap B(z,(R\sqrt{\alpha})))}{\pi (R\sqrt{\alpha})^{2}}
% =\inf_{z\in\C}\frac{\#\bigl(\Lambda_\alpha\cap B(z/\sqrt{\alpha},R)\bigr)}{ \alpha(R-r)^2}\frac{\alpha(R-r)^2}{\pi(R\sqrt{\alpha})^{2}}.
% \]
% Letting $\alpha\to\infty$ gives
% \[
% D^{-}(\Lambda)\ \ge\ \frac{1}{\pi}\frac{(R-r)^2}{R^2}.
% \]
% Finally, taking $r\to 0$ we get 
% \[
% D^{-}(\Lambda)\ \ge\ \frac{1}{\pi}.
% \]

\begin{comment}
If $\Lambda$ is interpolating for $\calF^{2}$,
then by Lemma~\ref{lem:scalingFock}, $\Lambda_\alpha$ is interpolating in
$\Falpha$.
Applying the duality/surjectivity argument to $T^\alpha_{a,R}$ yields
\[
\#\bigl(\Lambda_\alpha\cap B(a,R)\bigr)
\ \le\ 
\alpha R^{2}+o(\alpha)
\qquad\text{uniformly in }a.
\]
Undoing the scaling as before shows
\[
\#\bigl(\Lambda\cap B(z,R\sqrt{\alpha})\bigr)
\le \alpha R^{2}+o(\alpha).
\]
Dividing by $\pi(R\sqrt{\alpha})^{2}$ and letting $\alpha\to\infty$ yields
\[
D^{+}(\Lambda)\ \le\ \frac{1}{\pi}.
\]
\end{comment}

% \subsection{Necessary Density for Sampling in Bergman Spaces of $\mathbb{B}^{n}$}

% Finally, we discuss an application that is new for all dimensions $n\ge2$.  
% Let $\Lambda\subset\mathbb{B}^n$ be a sampling sequence for the unweighted Bergman space $A^{2}(\mathbb{B}^n)$:
% \[
% A\|f\|^{2}_{A^{2}} 
% \;\le\; 
% \sum_{\lambda\in\Lambda} |\langle f,k_\lambda\rangle|^{2}
% \;\le\; 
% B\|f\|^{2}_{A^{2}},
% \qquad f\in A^{2}(\mathbb{B}^{n}).
% \]
% Define the lower (hyperbolic) Beurling density of $\Lambda$ by
% \[
% D_{-}(\Lambda)
% =
% \liminf_{r\to\infty}
% \inf_{a\in\mathbb{B}^n}
% \frac{\#(\Lambda\cap B(a,r))}{\mathrm{Vol}_{\mathrm{hyp}}(B(a,r))}.
% \]

% Using the Bergman-space version of the Szeg\H{o} limit (Theorem~5.2 with $X=\mathbb{B}^n$), 
% together with a higher-dimensional version of Landau's argument and the exponential localization of Bergman kernels, 
% we obtain the following sharp necessary density condition.

% \begin{thm}[Necessary Density for Bergman Sampling]
% \label{thm:BergmanDensity}
% Let $n\ge2$ and let $\Lambda\subset\mathbb{B}^{n}$ be a sampling sequence for $A^{2}(\mathbb{B}^{n})$.  
% Then
% \[
% D_{-}(\Lambda)\;\ge\; 1,
% \]
% where the hyperbolic volume is normalized so that $\mathrm{Vol}_{\mathrm{hyp}}(B(a,1))=1$.
% \end{thm}

% When $n=1$, this reduces to the classical theorem of Seip.  
% For $n\ge2$, the problem has remained open due to the non-doubling nature of the underlying hyperbolic measure.  
% Our approach circumvents this difficulty by combining the mean-value property for normalized Bergman kernels with the spectral asymptotics of Toeplitz operators via the Szeg\H{o} limit.  
% Full details and several refinements will appear in a companion work.




%%%%%%%%%%%%%%%%%%%%%%%%%%%%%%%%%%%%%%%%%%%%%%%
% \section{Applications}
% In this section, we will discuss several settings which satisfy Assumptions (I) and (II). In particular, it is known that for RKHS's which can be embedded into $L^{2}(X,\nu)$ for some measure space $(X,\nu)$, the associated kernels $\{k_{x}\}$ form a continuous Parseval frame. The first four examples to follow satisfy such an embedding condition. Thus, it is sufficient to show that they satisfy Assumptions (I) and (II). The final example, for which Assumption (II) is not easily verified, will be demonstrated directly by the convergence in measure of the $\alpha$-Berezin transform. 

%This gives the following corollary.
%\begin{cor}Let $\sigma$ be non-negative and $\sigma\in L^{p}(\mathbb{D},(1-|z|^{2})^{-2}dA(z))\cap L^{\infty}(\mathbb{D},(1-|z|^{2})^{-2}dA(z))$, for $p\in [1,\infty)$. Then, for $\psi$ continuous, we have \begin{align*} \frac{1}{\alpha+1}\tr\left(T_{\sigma}\psi\left(T_{\sigma}^{\alpha}\right)\right)\to\int_{\mathbb{D}}\sigma(z)\psi\left(\sigma(z)\right)\frac{1}{(1-|z|^{2})^{2}}dA(z). \end{align*}\end{cor}
           
% \subsection{Segal-Bargmann-Fock Space}

% Let $\alpha>0$, and define $\mathcal{H}_{\alpha}=\mathcal{F}^{2}_{\alpha}(\C):=\{f\in L^{2}(\C,\frac{\alpha}{\pi}e^{-\alpha|z|^{2}}dA(z)):\text{f is entire}\}$, where $dA(z)$ is the area measure on $\C$. It is known that $\mathcal{F}^{2}_{\alpha}(\C)$ is a RKHS \cite{Z12} with the $L^{2}(\C,\frac{\alpha}{\pi}e^{-\alpha|z|^{2}}dA(z))$ inner product and reproducing kernels
% \begin{align*}
%     \ip{K_{z}^{\alpha}}{K_{w}^{\alpha}}_{\alpha}=K^{\alpha}(z,w)=e^{\alpha z\overline{w}}.
% \end{align*}
% We have that $\norm{K_{z}^{\alpha}}_{\alpha}^{2}=K^{\alpha}(z,z)=e^{\alpha|z|^{2}}$. Using this, we obtain the normalized reproducing kernel:
% \begin{align*}
%     \ip{k_{z}^{\alpha}}{k_{w}^{\alpha}}_{\alpha}=\frac{1}{\norm{K_{z}^{\alpha}}_{\alpha}\norm{K_{w}^{\alpha}}_{\alpha}}\ip{K_{z}^{\alpha}}{K_{w}^{\alpha}}_{\alpha}=&e^{\alpha z\overline{w}}e^{-\frac{\alpha}{2}|z|^{2}}e^{-\frac{\alpha}{2}|w|^{2}}\\
%     =&\exp{\left(-\frac{\alpha}{2}\left(|z|^{2}-2z\overline{w}+|w|^{2}\right)\right)}.
% \end{align*}
% Observe that 
% \begin{align*}
%     \left|\ip{k_{z}^{\alpha}}{k_{w}^{\alpha}}_{\alpha}\right|^{2}=&\left|\exp{\left(-\frac{\alpha}{2}\left(|z|^{2}-2z\overline{w}+|w|^{2}\right)\right)}\right|^{2}\\
%     =&\left|\exp{\left(-\frac{\alpha}{2}\left(|z|^{2}-2\text{Re}z\overline{w}-2i\text{Im}z\overline{w}+|w|^{2}\right)\right)}\right|^{2}\\
%     =&\left|\exp{\left(-\alpha i\text{Im}z\overline{w}\right)}\right|^{2}\left|\exp{\left(-\frac{\alpha}{2}\left(|z|^{2}-2\text{Re}z\overline{w}+|w|^{2}\right)\right)}\right|^{2}\\
%     =&\exp{\left(-\alpha |z-w|^{2}\right)}.
% \end{align*}
% Now, note that the associated measure under this normalized reproducing kernel is 
% \begin{align*}
%     d\nu_{\alpha}(z)=\norm{K_{z}^{\alpha}}_{\alpha}^{2}\frac{\alpha}{\pi}e^{-\alpha |z|^{2}}dA(z)=e^{\alpha |z|^{2}}\frac{\alpha}{\pi}e^{-\alpha |z|^{2}}dA(z)=\frac{\alpha}{\pi}dA(z).
% \end{align*}
% Thus, $\nu_{\alpha}$ satisfies Assumption (I), by letting $c(\alpha)=\alpha/\pi$ and $d\nu=dA(z)$. We now check Assumption (II). Let $z\in\C$ and $R>0$. Then we have 
% \begin{align*}
%    0\leq \int_{B(z,R)^{c}}\left|\ip{k_{z}^{\alpha}}{k_{w}^{\alpha}}_{\alpha}\right|^{2}d\nu_{\alpha}(z)=&\int_{B(z,R)^{c}}\frac{\alpha}{\pi}\exp{\left(-\alpha |z-w|^{2}\right)}dA(z)\\
%    =&\int_{B(0,R)^{c}}\frac{\alpha}{\pi}\exp{\left(-\alpha|w|^{2}\right)}dA(z)\\
%    =&2\alpha\int_{R}^{\infty}r\exp{\left(\alpha r^{2}\right)}dr\\
%    =&\exp{\left(-\alpha R^{2}\right)}\to 0,
% \end{align*}
%  as $\alpha\to\infty$. Thus, Assumption (II) is satisfied. Hence, we obtain the following corollary to Corollary~\ref{Sze_Berg}:
% \begin{cor}
% Let $\sigma$ be non-negative and $\sigma\in L^{1}(\C,dA(z))\cap L^{\infty}(\C,dA(z))$. Then, for any $\psi:[0,\norm{\sigma}_{\infty}]\to\mathbb{R}$ continuous, we have 
% \begin{align*}
%     \frac{\pi}{\alpha}\tr\left(T_{\sigma}\psi\left(T_{\sigma}^{\alpha}\right)\right)\to\int_{\C}\sigma(z)\psi\left(\sigma(z)\right)dA(z),
% \end{align*}
% as $\alpha\to\infty$. 
% \end{cor}     
% \subsection{Paley-Wiener Space}

% Let $\alpha>0$, and let $\mathcal{H}_{\alpha}=PW_{\alpha}=\{f\in L^{2}(\R):\text{supp}(\hat{f})\subset[-\alpha,\alpha]\}$, equipped with the Lebesgue measure. It is known that $PW_{\alpha}$ is a RKHS over $\R$, with reproducing kernel
% \begin{align*}
%     \ip{K_{x}^{\alpha}}{K_{y}^{\alpha}}_{\alpha}=K^{\alpha}(x,y)=\begin{cases}
%         \frac{1}{\pi}\frac{\sin{(2\pi\alpha(x-y))}}{x-y} & x\neq y\\
%         2\alpha & x=y.
%     \end{cases}
% \end{align*}   
% Noting that $\norm{K_{x}^{\alpha}}^{2}=2\alpha$, we obtain the normalized kernel
% \begin{align*}
%     \ip{k_{x}^{\alpha}}{k_{y}^{\alpha}}_{\alpha}=\begin{cases}
%         \frac{\sin{(2\pi\alpha(x-y))}}{2\pi\alpha(x-y)} & x\neq y\\
%         1 & x=y,
%     \end{cases}
% \end{align*}
% and the associated measure $d\nu_{\alpha}(x)=2\alpha dx$, where $dx$ is the Lebesgue measure on $\R$. Letting $c(\alpha)=2\alpha$, and $\nu=dx$, we have that Assumption (I) is satisfied. Now, note that 
% \begin{align*}
%     \left|\ip{k_{x}^{\alpha}}{k_{y}^{\alpha}}_{\alpha}\right|^{2}=\begin{cases}
%         \left|\frac{\sin{(2\pi\alpha(x-y))}}{2\pi\alpha(x-y)}\right|^{2} & x\neq y\\
%         1 & x=y,
%     \end{cases}
% \end{align*}
% We now verify Assumption (II). Let $x\in \mathbb{R}$, and let $R>0$. Then, 
% \begin{align*}
%     0\leq \int_{B(x,R)^{c}}2\alpha\left|\ip{k_{x}^{\alpha}}{k_{y}^{\alpha}}_{\alpha}\right|^{2}dy=&\int_{B(x,R)^{c}}2\alpha\left|\frac{\sin{(2\pi\alpha(x-y))}}{2\pi\alpha(x-y)}\right|^{2}dy\\
%     =&\int_{B(0,R)^{c}}2\alpha\left|\frac{\sin{(2\pi\alpha(-y))}}{2\pi\alpha(-y)}\right|^{2}dy\\
%     =&\int_{R}^{\infty}4\alpha\left|\frac{\sin{(2\pi\alpha y)}}{2\pi\alpha y}\right|^{2}dy\\
%     \leq&\int_{R}^{\infty}4\alpha\frac{1}{\left|2\pi\alpha y\right|^{2}}dy\\
%     =&\frac{1}{\pi^{2}\alpha}\int_{R}^{\infty}\frac{1}{y^{2}}dy\to 0,
% \end{align*}
% as $\alpha\to\infty$, and Assumption (II) follows. Applying Corollary~\ref{Sze_Berg} we obtain the following result.
% \begin{cor}
%     Let $\sigma$ be non-negative and $\sigma\in L^{1}(\R,dx)\cap L^{\infty}(\R,dx)$. Then for $\psi:\R\to\R$ continuous, we have 
%     \begin{align*}
%         \frac{1}{2\alpha}\tr\left(T_{\sigma}\psi\left(T_{\sigma}^{\alpha}\right)\right)\to\int_{\R}\sigma(x)\psi\left(\sigma(x)\right)dx.
%     \end{align*}
% \end{cor}
% \begin{subsection}{The Classical Szeg\H{o} Limit Theorem}
%     In this section, we will prove that the classical Szeg\H{o} Limit Theorem is also a consequence of the general theorems we have proved. With this aim, let $\mathcal{K}_{n}=\text{span}\{e^{ik\theta}:k=0,\pm 1,\pm 2,\cdots,\pm n\}$. Since $\{e^{ik\theta}:k\in\mathbb{Z}\}$ form an orthonormal basis for $L^{2}(\mathbb{T},\frac{1}{2\pi}dx)$, $\mathcal{K}_{n}$ is a closed subspace of $L^{2}(\mathbb{T})$. Furthermore, pointwise evaluation is a bounded linear functional. In fact, for $\theta\in\mathbb{T}$ and $p(\theta)\in \mathcal{K}_{n}$,
%     \begin{align*}
%         |p(\theta)|=\left|\sum_{k=-n}^{n}a_{k}e^{ik\theta}\right|\leq\sum_{k=-n}^{n}|a_{k}|\leq\left(\sum_{k=-n}^{n}|a_{k}|^{2}\right)^{1/2}\sqrt{2n+1}=\norm{p}_{2}\sqrt{2n+1}.
%     \end{align*}
%     Hence, the Riesz Representation Theorem applies and $\mathcal{K}_{n}$ is a RKHS with the $L^{2}(\mathbb{T},\frac{1}{2\pi}dx)$ inner product. By standard results of the theory of RKHS \cite{paulsen2016introduction}, the reproducing kernel for this space is 
%     \begin{align*}
%         K_{n}(x,y)=&\sum_{k=-n}^{n}e^{ikx}e^{-iky}=D_{n}(x-y),
%     \end{align*}
%     where $D_{n}$ is the Dirichlet kernel. Now, note that $D_{n}(0)=2n+1$ and hence $K_n(x,x)=2n+1$. Using this, we now check Assumption (II). For this, it is sufficient to check the assumption at $x=0$, since the kernel is symmetric. Let $R>0$. We have 
%     \begin{align*}
%         (2n+1)\left|\ip{k_{0}^{n}}{k_{x}^{n}}_{n}\right|^{2}=&(2n+1)\left|\frac{1}{2n+1}D_{n}(x)\right|^{2}\\
%         =&\frac{1}{2n+1}\left|D_{n}(x)\right|^{2}\\
%         =&\frac{1}{2n+1}\left(\frac{\sin{\left((n+1/2)x\right)}}{\sin{\left(x/2\right)}}\right)^{2},
%     \end{align*}
%     which is the Fej\'er kernel, a well-known approximate identity. Thus, Assumption (II) is satisfied. From this we obtain the following corollary, which is essentially the classical Szeg\H{o} limit theorem. 
%     \begin{cor}
%         Let $\sigma\in L^{\infty}(\mathbb{T})$ be real-valued, and let $\psi:[\inf\sigma,\sup\sigma]\to\mathbb{R}$ be continuous. Then, 
%         \begin{align*}
%             \lim_{n\to\infty}\frac{1}{2n+1}\text{tr}\left(\psi\left(T_{\sigma}^{n}\right)\right)=\int_{0}^{2\pi}\psi\left(\sigma(x)\right)\frac{dx}{2\pi},
%         \end{align*}
%         where $T_{\sigma}^{n}:\mathcal{K}_{n}\to\mathcal{K}_{n}$ is the Toeplitz operator with symbol $\sigma$ acting on $\mathcal{K}_{n}$. 
%     \end{cor}
    %Likewise, we have the following corollary which is similar to \cite{guo2024first}.\begin{cor} Let $\sigma\in L^{1}(\mathbb{T})$ be non-negative. Let $\psi:[0,\infty)\to\mathbb{R}$ be continuous and suppose $\lim_{x\to\infty}\psi(x)/x$ exists. Then, \begin{align*}\lim_{n\to\infty}\frac{1}{2n+1}\text{tr}\left(\psi\left(T_{\sigma}^{n}\right)\right)=\int_{0}^{2\pi}\psi\left(\sigma(x)\right)\frac{dx}{2\pi}.\end{align*}\end{cor}
% \end{subsection}

% \begin{align*}
%     \lim_{n\to\infty}\frac{|K\cdot\Gamma_{N}|}{|\Gamma_{N}|}=1.
% \end{align*}
% We prove the following lemma to obtain the corresponding version of Theorem 1.2 when $G$ is a compact Abelian group. 
% \begin{lem}
%     Let $\{\Gamma_{N}\}\subset G$ be a F\o lner sequence. Then $\{\Gamma_{N}\}$ satisfies the ratio property. 
% \end{lem}
% \begin{proof}
%     Note that 
%     \begin{align*}
%         K\cdot\Gamma_{N}=(K\cdot\Gamma_{N}\cap\Gamma_{N})\cup(K\cdot\Gamma_{N}\cap\Gamma_{N}^{c})=(K\cdot\Gamma_{N}\cap\Gamma_{N})\cup(K\cdot\Gamma_{N}\setminus\Gamma_{N}).
%     \end{align*}
%     Likewise,
%     \begin{align*}
%         \Gamma_{N}=(K\cdot\Gamma_{N}\cap\Gamma_{N})\cup((K\cdot\Gamma_{N})^{c}\cap\Gamma_{N})=(K\cdot\Gamma_{N}\cap\Gamma_{N})\cup(\Gamma_{N}\setminus(K\cdot\Gamma_{N})).
%     \end{align*}
%     Additionally, all of the above unions are disjoint. Hence, 
%     \begin{align*}
%         \left||K\cdot\Gamma_{N}|-|\Gamma_{N}|\right|=\left||K\cdot\Gamma_{N}\setminus\Gamma_{N}|-|\Gamma_{N}\setminus(K\cdot\Gamma_{N})|\right|\leq|(K\cdot\Gamma_{N})\Delta\Gamma_{N}|.
%     \end{align*}
%     Thus, dividing by $|\Gamma_{N}|$ and letting $N\to\infty$ gives the result.
% \end{proof}


% Thus, for $\sigma\in L^{2}(G)$ we see 
% \begin{align*}
%     \Tilde{\sigma}^{N}(x)=\int_{G}\sigma(y)|\ip{k^N_x}{k^N_y}|^2|\Gamma_{N}|d\mu(y)=&\frac{1}{|\Gamma_{N}|}\sum_{\eta\in\Gamma_{N}}\sum_{\xi\in\Gamma_{N}}\hat{\sigma}(\xi\cdot\overline{\eta})\xi\cdot\overline{\eta}(x).
% \end{align*}
% Additionally, using Fourier inversion we see that
% \begin{align*}
%     \sigma(x)=\sum_{\xi\in \hat{G}}\hat{\sigma}(\xi)\xi(x)=\frac{1}{|\Gamma_{N}|}\sum_{\eta\in\Gamma_{N}}\sum_{\xi\in \hat{G}}\hat{\sigma}(\xi\cdot\overline{\eta})\xi\cdot\overline{\eta}(x).
% \end{align*}
% Thus, 
% \begin{align*}
%     |\Tilde{\sigma}^{N}(x)-\sigma(x)|=&\left|\frac{1}{|\Gamma_{N}|}\sum_{\eta\in\Gamma_{N}}\sum_{\gamma\notin\eta\Gamma_{N}}\hat{\sigma}(\gamma)\gamma(x)\right|\\
%     =&\left|\frac{1}{|\Gamma_{N}|}\left(\sum_{\eta\in\Gamma_{N}}\sum_{\gamma\notin\eta\Gamma_{N}}\right)^{'}\ell(\gamma)\hat{\sigma}(\gamma)\gamma(x)\right|,
% \end{align*}
% where $(\sum)^{'}$ denotes summing over unique elements in the sum, and $\ell(\gamma)$ denote the number of repeats of a unique $\gamma$ appearing in the sum. It is clear that $\ell(\gamma)\leq|\Gamma_{N}|$ for each $\gamma$ in the sum. Thus, since the elements of $\hat{G}$ form an orthonormal basis for $L^{2}(G)$ we see that  
% \begin{align*}
%     \norm{\Tilde{\sigma}^{N}-\sigma}_{2}^{2}=&\frac{1}{|\Gamma_{N}|^{2}}\left(\sum_{\eta\in\Gamma_{N}}\sum_{\gamma\notin\eta\Gamma_{N}}\right)^{'}\ell(\gamma)^{2}\left|\hat{\sigma}(\gamma)\right|^{2}\leq\frac{1}{|\Gamma_{N}|}\sum_{\eta\in\Gamma_{N}}\sum_{\gamma\notin\eta\Gamma_{N}}|\hat{\sigma}(\gamma)|^{2}.
% \end{align*}
% Let $\varepsilon>0$. Since $\sigma\in L^{2}(G)$, there is a compact set such that 
% \begin{align*}
%     \sum_{\gamma\in K^{c}}|\hat{\sigma}(\gamma)|^{2}<\varepsilon/2.
% \end{align*}
% Likewise, by the F\o lner condition, there is an $N$ sufficiently large such that 
% \begin{align*}
%     \frac{|K\Gamma_{N}\Delta\Gamma_{N}|}{|\Gamma_{N}|}<\varepsilon/(2\norm{\sigma}_{2}^{2}).
% \end{align*}
% Then, 
% \begin{align*}
%     \frac{1}{|\Gamma_{N}|}\sum_{\eta\in\Gamma_{N}}\sum_{\gamma\notin\eta\Gamma_{N}}|\hat{\sigma}(\gamma)|^{2}=&\frac{1}{|\Gamma_{N}|}\sum_{\eta\in\Gamma_{N}}\sum_{\gamma\in\Gamma_{N}^{c}}|\hat{\sigma}(\gamma\overline{\eta})|^{2}\\
%     =&\frac{1}{|\Gamma_{N}|}\sum_{\eta\in\Gamma_{N}}\sum_{\gamma\in K\Gamma_{N}\setminus\Gamma_{N}}+\sum_{\gamma\in(K\Gamma_{N})^{c}}|\hat{\sigma}(\gamma\overline{\eta})|^{2}\\
%     =&(I)+(II).
% \end{align*}
% For the first sum, we have
% \begin{align*}
%     (I)=\frac{1}{|\Gamma_{N}|}\sum_{\eta\in\Gamma_{N}}\sum_{\gamma\in K\Gamma_{N}\setminus\Gamma_{N}}|\hat{\sigma}(\gamma\overline{\eta})|^{2}\leq\norm{\sigma}_{2}^{2}\frac{|K\Gamma_{N}\setminus\Gamma_{N}|}{|\Gamma_{N}|}\leq\norm{\sigma}_{2}^{2}\frac{|K\Gamma_{N}\Delta\Gamma_{N}|}{|\Gamma_{N}|}<\varepsilon/2.
% \end{align*}
% For the second sum, note that if $\gamma\in(K\Gamma_{N})^{c}$ and $\eta\in\Gamma_{N}$ then $\gamma\overline{\eta}\notin K$. Thus, 
% \begin{align*}
%     (II)=\frac{1}{|\Gamma_{N}|}\sum_{\eta\in\Gamma_{N}}\sum_{\gamma\in(K\Gamma_{N})^{c}}|\hat{\sigma}(\gamma\overline{\eta})|^{2}\leq&\frac{1}{|\Gamma_{N}|}\sum_{\eta\in\Gamma_{N}}\sum_{\gamma\in K^{c}}|\hat{\sigma}(\gamma)|^{2}=\sum_{\gamma\in K^{c}}|\hat{\sigma}(\gamma)|^{2}<\varepsilon/2.
% \end{align*}
% Thus, the convergence is established and we obtain Theorem 1.2.
\bibliographystyle{plain}
\bibliography{refs}
\end{document}



